% Generated mechanically from frozen input p01/ko/sheaves.tex.
% Do not edit this generated file; edit the bound locale source instead.

% OK, start here.
%

\setcounter{chapter}{5}
\chapter{공간 위의 층}



\phantomsection
\label{sheaves-section-phantom}


\section{서론}
\label{sheaves-section-introduction}

\noindent
이 문서에서는 위상공간 위의 층이 지닌 기본 성질을 설명한다.
참고문헌으로 \cite{Godement}가 있다.

\medskip\noindent
뒤에서 다룰 사이트 위의 층에 관한 논의가 이 내용을 대신하게 될 것이다.
그러나 여기서 몇 가지 개념을 간략히 정의해 두는 것도 의미가 있을 듯하다.









\section{기본 개념}
\label{sheaves-section-sheaves-basic}

\noindent
다음은 위상수학의 기본 개념을 모아 놓은 목록이다.

\begin{enumerate}
\item $X$를 위상공간이라 하자. ``$U = \bigcup_{i \in I} U_i$가 열린 덮개라고
하자''라는 표현은 다음을 뜻한다. $I$는 집합이고, 각 $i \in I$에 대하여
열린 부분집합 $U_i \subset X$가 주어지며, $U$는 $U_i$들의 합집합이다.
$I = \emptyset$인 것도 허용한다. 이 경우에는 $U_i$가 없고
$U = \emptyset$이다. 또한 $I \not = \emptyset$인 경우에도 일부 또는 모든
$U_i$가 공집합이어도 된다.
% P01-ERRATA: P01-E0035 -- the frozen reader-facing list contains an editorial placeholder.
\item 기타 등등.
\end{enumerate}









\section{준층}
\label{sheaves-section-presheaves}


\begin{definition}
\label{sheaves-definition-presheaf}
$X$를 위상공간이라 하자.
\begin{enumerate}
\item $X$ 위의 {\it 집합의 준층 $\mathcal{F}$}란, 각 열린집합
$U \subset X$에 집합 $\mathcal{F}(U)$를 대응시키고 각 포함관계
$V \subset U$에 사상
$\rho^U_V : \mathcal{F}(U) \to \mathcal{F}(V)$를 대응시키는 규칙으로서,
$\rho^U_U = \text{id}_{\mathcal{F}(U)}$이고
$W \subset V \subset U$일 때마다
$\rho^U_W = \rho^V_W \circ \rho ^U_V$를 만족하는 것이다.
\item $X$ 위의 집합의 준층 사이의 {\it 사상
$\varphi : \mathcal{F} \to \mathcal{G}$}이란, 각 열린집합
$U \subset X$에 제한 사상과 양립하는 집합의 사상
$\varphi : \mathcal{F}(U) \to \mathcal{G}(U)$를 대응시키는 규칙이다.
즉, $V \subset U \subset X$가 열린집합일 때마다 도식
$$
\xymatrix{
\mathcal{F}(U) \ar[r]^\varphi \ar[d]^{\rho^U_V} &
\mathcal{G}(U) \ar[d]^{\rho^U_V} \\
\mathcal{F}(V) \ar[r]^\varphi & \mathcal{G}(V)
}
$$
은 가환한다.
\item $X$ 위의 집합의 준층들이 이루는 범주를
$\textit{PSh}(X)$로 나타낸다.
\end{enumerate}
\end{definition}

\noindent
집합 $\mathcal{F}(U)$의 원소를 $U$ 위에서의 $\mathcal{F}$의
{\it 단면}이라 한다. 각 $V \subset U$에 대하여 사상
$\rho^U_V : \mathcal{F}(U) \to \mathcal{F}(V)$를
{\it 제한 사상}이라 한다. $s\in \mathcal{F}(U)$이면
$s|_V := \rho^U_V(s)$라는 표기를 사용한다. 이 표기는 위상수학에서 함수의
제한이라는 개념과 양립한다. 실제로 $W \subset V \subset U$이고 $s$가
$U$ 위에서의 $\mathcal{F}$의 단면이면, 위 정의에 나타난 제한 사상의
성질에 따라 $s|_W = (s|_V)|_W$이다.

\medskip\noindent
열린집합 $U$ 위의 단면을 기호 $\Gamma(U, -)$ 또는
$H^0(U, -)$로 나타내는 표기도 흔히 쓰인다. 다시 말해, 다음 등식들은
항상 성립한다.
$$
\Gamma(U, \mathcal{F}) = \mathcal{F}(U) = H^0(U, \mathcal{F}).
$$
이 장에서는 이 표기를 사용하지 않지만, 다른 장에서는 사용할 것이다.

\begin{definition}
\label{sheaves-definition-constant-presheaf}
$X$를 위상공간이라 하고 $A$를 집합이라 하자.
{\it 값이 $A$인 상수 준층}은 각 열린집합 $U \subset X$에
집합 $A$를 대응시키고 모든 제한 사상이 $\text{id}_A$인 준층이다.
\end{definition}







\section{아벨 준층}
\label{sheaves-section-abelian-presheaves}

\noindent
이 절에서는 아벨군의 준층을 정의할 수 있게 해 주는
준층 범주의 몇 가지 성질을 간단히 살펴본다.

\begin{example}
\label{sheaves-example-singleton-presheaf}
$X$를 위상공간이라 하자. $X$의 모든 열린 부분집합에 한원소 집합을
대응시키는 규칙 $\mathcal{F}$를 생각하자. 모든 집합에서 한원소 집합으로
가는 사상은 유일하므로, 유일한 제한 사상 $\rho^U_V$들이 존재한다.
이렇게 얻은 구조는 $X$ 위의 집합의 준층이다. 앞서 언급한 한원소 집합의
성질에 따라 이것은 $X$ 위의 집합의 준층 범주에서 끝 대상이다.
따라서 유일한 동형사상까지 유일하다. 이 준층을 때때로 $*$로 나타낸다.
\end{example}

\begin{lemma}
\label{sheaves-lemma-product-presheaves}
$X$를 위상공간이라 하자. $X$ 위의 집합의 준층 범주에는 곱이 존재한다
(Categories, 정의 \ref{categories-definition-product} 참조).
또한 열린집합 $U$ 위에서 곱 $\mathcal{F} \times \mathcal{G}$의
단면들의 집합은 $U$ 위에서 $\mathcal{F}$와 $\mathcal{G}$의
단면들의 집합의 곱이다.
\end{lemma}

\begin{proof}
구체적으로, $\mathcal{F}$와 $\mathcal{G}$를 위상공간 $X$ 위의
집합의 준층이라 하자. $\mathcal{F} \times \mathcal{G}$로 나타내는 규칙
$U \mapsto \mathcal{F}(U) \times \mathcal{G}(U)$를 생각하자.
$V \subset U \subset X$가 열린집합이면 제한 사상
$$
(\mathcal{F} \times \mathcal{G})(U)
\longrightarrow
(\mathcal{F} \times \mathcal{G})(V)
$$
을 $(s, t) \mapsto (s|_V, t|_V)$로 정의한다. 그러면
$\mathcal{F} \times \mathcal{G}$가 준층임은 즉시 알 수 있다.
또한 사영 사상
$p : \mathcal{F} \times \mathcal{G} \to \mathcal{F}$와
$q : \mathcal{F} \times \mathcal{G} \to \mathcal{G}$가 있다.
임의의 세 번째 준층 $\mathcal{H}$에 대하여
$\Mor(\mathcal{H}, \mathcal{F} \times \mathcal{G})
= \Mor(\mathcal{H}, \mathcal{F}) \times
\Mor(\mathcal{H}, \mathcal{G})$임을 보이는 것은 독자에게 맡긴다.
\end{proof}

\noindent
$(A, + : A \times A \to A, - : A \to A, 0\in A)$가 아벨군이면
영사상과 역원 사상은 덧셈 법칙에 의해 유일하게 결정된다는 것을
상기하자. 다시 말해, ``$(A, +)$를 아벨군이라 하자''라고 말하는 것은
의미가 있다.

\begin{lemma}
\label{sheaves-lemma-abelian-presheaves}
$X$를 위상공간이라 하자.
$\mathcal{F}$를 집합의 준층이라 하자.
$\mathcal{F}$ 위의 다음 구조들을 생각하자.
\begin{enumerate}
\item 각 열린집합 $U$에 대하여 $\mathcal{F}(U)$ 위에 아벨군 구조를
주고, 모든 제한 사상이 아벨군 준동형이 되게 하는 구조.
\item 준층의 사상
$+ : \mathcal{F} \times \mathcal{F} \to \mathcal{F}$,
준층의 사상 $- : \mathcal{F} \to \mathcal{F}$,
그리고 사상 $0 : * \to \mathcal{F}$
(예 \ref{sheaves-example-singleton-presheaf} 참조)로서, 보통의 아벨군에서
$+, -, 0$이 만족하는 모든 공리를 만족하는 것.
\item 준층의 사상
$+ : \mathcal{F} \times \mathcal{F} \to \mathcal{F}$,
준층의 사상 $- : \mathcal{F} \to \mathcal{F}$,
그리고 사상 $0 : * \to \mathcal{F}$로서, 각 열린집합
$U \subset X$에 대하여 사중항
$(\mathcal{F}(U), +, -, 0)$이 아벨군이 되게 하는 것.
\item 준층의 사상 $+ : \mathcal{F} \times \mathcal{F}
\to \mathcal{F}$로서, 모든 열린집합 $U \subset X$에 대하여 사상
$+ : \mathcal{F}(U) \times \mathcal{F}(U) \to \mathcal{F}(U)$가
아벨군 구조를 정의하게 하는 것.
\end{enumerate}
위의 (1)--(4) 유형의 데이터들의 모임 사이에는 자연스러운 전단사가 있다.
\end{lemma}

\begin{proof}
생략한다.
\end{proof}























\noindent
이 보조정리는 준층 범주에서 아벨군 대상 $\mathcal{F}$를 주는 것이,
모든 집합 $\mathcal{F}(U)$에 아벨군 구조가 주어지고 모든 제한 사상이
군 준동형이 되는 집합의 준층 $\mathcal{F}$를 주는 것과 같다는 뜻이다.
대부분의 대수 구조에 대해서는 이러한 대상들의 (준)층을 이와 같은
방식으로 다룰 것이다. 즉, 이러한 대상들의 (준)층을, 모든 단면 집합
$\mathcal{F}(U)$가 제한 사상과 양립하도록 해당 구조를 갖춘
집합의 (준)층 $\mathcal{F}$로 정의할 것이다.

\begin{definition}
\label{sheaves-definition-abelian-presheaves}
$X$를 위상공간이라 하자.
\begin{enumerate}
\item {\it $X$ 위의 아벨군의 준층} 또는
{\it $X$ 위의 아벨 준층}이란 집합의 준층 $\mathcal{F}$로서,
각 열린집합 $U \subset X$에 대하여 집합 $\mathcal{F}(U)$가
아벨군 구조를 갖고 모든 제한 사상 $\rho^U_V$가 아벨군 준동형인 것이다.
위의 보조정리 \ref{sheaves-lemma-abelian-presheaves}를 보라.
\item {\it $X$ 위의 아벨 준층의 사상}
$\varphi : \mathcal{F} \to \mathcal{G}$이란, 모든 열린집합
$U \subset X$에 대하여 아벨군 준동형
$\mathcal{F}(U) \to \mathcal{G}(U)$를 유도하는 집합의 준층의 사상이다.
\item $X$ 위의 아벨군의 준층 범주는 $\textit{PAb}(X)$로 나타낸다.
\end{enumerate}
\end{definition}

\begin{example}
\label{sheaves-example-direct-sum-points}
$X$를 위상공간이라 하자. 각 $x \in X$에 대하여 아벨군 $M_x$가
주어졌다고 하자. 열린집합 $U \subset X$에 대하여
$$
\mathcal{F}(U) = \bigoplus\nolimits_{x \in U} M_x.
$$
로 둔다. 이 아벨군의 전형적인 원소를
$\sum_{i = 1}^n m_{x_i}$로 나타내며, 여기서 $x_i \in U$이고
$m_{x_i} \in M_{x_i}$이다. (물론 $x_1, \ldots, x_n$이 서로
다르도록 표현을 언제나 선택할 수 있다.)
열린집합 $V \subset U \subset X$에 대하여 원소
$s = \sum_{i = 1}^n m_{x_i}$를
$s|_V = \sum_{x_i \in V} m_{x_i}$로 보내는 제한 사상
$\mathcal{F}(U) \to \mathcal{F}(V)$를 정의한다.
이것이 아벨군의 준층임을 확인하는 것은 독자에게 맡긴다.
\end{example}




\section{대수 구조의 준층}
\label{sheaves-section-presheaves-structures}

\noindent
대수 구조의 준층에 관한 정의를 명확히 하자.
$\mathcal{C}$가 범주이고
$F : \mathcal{C} \to \textit{Sets}$가 충실 함자라고 하자.
보통 $F$는 “망각” 함자이다. 대상 $M \in \Ob(\mathcal{C})$에 대하여
$F(M)$을 흔히 대상 $M$의 {\it 바탕 집합}이라 한다.
$M \to M'$이 $\mathcal{C}$의 사상이면
$F(M) \to F(M')$를 {\it 바탕 집합의 사상}이라 한다.
실제로 우리는 흔히 대상과 그 바탕 집합을 구별하지 않으며,
사상에 대해서도 마찬가지이다. 따라서 집합의 사상
$F(M) \to F(M')$가 $\mathcal{C}$의 어떤 사상
$f : M \to M'$에 대한 $F(f)$와 같으면 이를
{\it 대수 구조의 사상}이라 할 것이다.

\medskip\noindent
위의 정의 \ref{sheaves-definition-abelian-presheaves}와 유사하게,
“$\mathcal{C}$의 대상들의 준층”은 다음 데이터로 정의할 수 있다.
\begin{enumerate}
\item 집합의 준층 $\mathcal{F}$, 그리고
\item 모든 열린집합 $U \subset X$에 대하여 택한 대상
$A(U) \in \Ob(\mathcal{C})$.
\end{enumerate}
이 데이터는 다음 조건들을 만족해야 한다(위의 표현을 사용한다).
\begin{enumerate}
\item 모든 열린집합 $U \subset X$에 대하여 집합 $\mathcal{F}(U)$는
$A(U)$의 바탕 집합이고,
\item 모든 열린집합 $V \subset U \subset X$에 대하여 집합의 사상
$\rho_V^U: \mathcal{F}(U) \to \mathcal{F}(V)$는 대수 구조의 사상이다.
\end{enumerate}
% P01-ERRATA: P01-E1233 -- the frozen English sentence has plural/singular disagreement around this one restriction map.
다시 말해, $X$의 모든 열린집합 $V \subset U$에 대하여
제한 사상 $\rho^U_V$는 범주 $\mathcal{C}$의 어떤 유일한 사상
$\alpha^U_V : A(U) \to A(V)$의 상 $F(\alpha^U_V)$이다.
유일성은 $F$가 충실하다는 조건으로부터 강제되며, 또한
$W \subset V \subset U$가 $X$의 열린집합일 때마다
$\alpha^U_W = \alpha^V_W \circ \alpha^U_V$임을 함의한다.
계 $ (A(-), \alpha^U_V)$은 우리가 $X$ 위에서 $\mathcal{C}$에 값을 갖는
준층으로 정의할 대상이다. Sites, 정의
\ref{sites-definition-presheaf}와 비교하라. 규칙
$\mathcal{F}(U) = F(A(U))$와
$\rho_V^U = F(\alpha_V^U)$를 통하여 집합의 준층
$(\mathcal{F}, \rho_V^U)$를 복원한다.

\begin{definition}
\label{sheaves-definition-presheaf-values-in-category}
$X$를 위상공간이라 하자.
$\mathcal{C}$를 범주라 하자.
\begin{enumerate}
% P01-ERRATA: P01-E0036 -- unlike the set-valued definition, the frozen definition omits the identity restriction axiom.
\item {\it $X$ 위에서 $\mathcal{C}$에 값을 갖는 준층 $\mathcal{F}$}란,
모든 열린집합 $U \subset X$에 $\mathcal{C}$의 대상
$\mathcal{F}(U)$를 대응시키고, 각 포함관계 $V \subset U$에
$\mathcal{C}$의 사상
$\rho_V^U : \mathcal{F}(U) \to \mathcal{F}(V)$를 대응시키는 규칙으로서,
$W \subset V \subset U$일 때마다
$\rho_W^U = \rho_W^V \circ \rho_V^U$를 만족하는 것이다.
% P01-ERRATA: P01-E1234 -- the frozen definition leaves the component variable U unquantified.
\item $\mathcal{C}$에 값을 갖는 준층의 {\it 사상
$\varphi : \mathcal{F} \to \mathcal{G}$}이란 제한 사상과 양립하는
$\mathcal{C}$의 사상
$\varphi : \mathcal{F}(U) \to \mathcal{G}(U)$로 주어진다.
\end{enumerate}
\end{definition}

\begin{definition}
\label{sheaves-definition-underlying-presheaf-sets}
$X$를 위상공간이라 하자. $\mathcal{C}$를 범주라 하자.
$F : \mathcal{C} \to \textit{Sets}$를 충실 함자라 하자.
$\mathcal{F}$를 $X$ 위에서 $\mathcal{C}$에 값을 갖는 준층이라 하자.
집합의 준층 $U \mapsto F(\mathcal{F}(U))$를
{\it $\mathcal{F}$의 바탕 집합의 준층}이라 한다.
\end{definition}

\noindent
바탕 집합의 준층을 나타낼 때 같은 문자 $\mathcal{F}$를 사용하는 것이
관례이며, 이는 정의 \ref{sheaves-definition-presheaf-values-in-category} 앞의
논의에 따라 타당하다. 특히 “$s \in \mathcal{F}(U)$라 하자” 또는
“$s$를 $U$ 위에서의 $\mathcal{F}$의 단면이라 하자”라는 말은
$s \in F(\mathcal{F}(U))$를 뜻한다.

\medskip\noindent
이 표기와 정의들은 특히 다음에 적용된다.
{\it (반드시 아벨일 필요는 없는) 군, 환, 고정된 환 위의 가군,
고정된 체 위의 벡터공간 등의 준층}, 그리고
{\it 이들 사이의 사상}.

\section{가군의 준층}
\label{sheaves-section-presheaves-modules}

\noindent
$\mathcal{O}$가 $X$ 위의 환의 준층이라고 하자.
$X$ 위의 $\mathcal{O}$-가군의 준층이라는 개념을 정의하고자 한다.
정의 \ref{sheaves-definition-abelian-presheaves}와 유사하게, 이를 모든 열린집합
$U \subset X$에 대하여 집합 $\mathcal{F}(U)$가
($\mathcal{F}$와 $\mathcal{O}$의) 제한 사상과 양립하는
$\mathcal{O}(U)$-가군 구조를 갖는 집합의 준층 $\mathcal{F}$로
정의하고 싶을 수 있다. 그러나 관례적으로는 다음 정의를 사용하며,
이는 앞의 정의와 동치이다.

\begin{definition}
\label{sheaves-definition-presheaf-modules}
$X$를 위상공간이라 하고, $\mathcal{O}$를 $X$ 위의 환의 준층이라 하자.
\begin{enumerate}
\item {\it $\mathcal{O}$-가군의 준층}은 아벨 준층 $\mathcal{F}$와
집합의 준층의 사상
$$
\mathcal{O} \times \mathcal{F} \longrightarrow \mathcal{F}
$$
으로 주어지며, 모든 열린집합 $U \subset X$에 대하여 사상
$\mathcal{O}(U) \times \mathcal{F}(U) \to \mathcal{F}(U)$가
아벨군 $\mathcal{F}(U)$ 위에 $\mathcal{O}(U)$-가군 구조를 정의해야 한다.
\item $\mathcal{O}$-가군의 준층의 {\it 사상
$\varphi : \mathcal{F} \to \mathcal{G}$}이란 다음 도식이 가환하도록 하는
아벨 준층의 사상 $\varphi : \mathcal{F} \to \mathcal{G}$이다.
$$
\xymatrix{
\mathcal{O} \times \mathcal{F} \ar[r] \ar[d]_{\text{id} \times \varphi} &
\mathcal{F} \ar[d]^{\varphi} \\
\mathcal{O} \times \mathcal{G} \ar[r] &
\mathcal{G}
}
$$
\item 위와 같은 $\mathcal{O}$-가군 사상의 집합은
$\Hom_\mathcal{O}(\mathcal{F}, \mathcal{G})$로 나타낸다.
\item $\mathcal{O}$-가군의 준층 범주는
$\textit{PMod}(\mathcal{O})$로 나타낸다.
\end{enumerate}
\end{definition}

\noindent
$\mathcal{O}_1 \to \mathcal{O}_2$가 $X$ 위의 환의 준층의 사상이라고 하자.
이때 $\mathcal{F}$가 $\mathcal{O}_2$-가군의 준층이면, 합성
$$
\mathcal{O}_1 \times \mathcal{F}
\to
\mathcal{O}_2 \times \mathcal{F}
\to
\mathcal{F}.
$$
을 사용하여 $\mathcal{F}$를 $\mathcal{O}_1$-가군의 준층으로
생각할 수 있다. 환의 제한을 나타내기 위하여 이를 때때로
$\mathcal{F}_{\mathcal{O}_1}$로 쓴다. 이것을 $\mathcal{F}$의
{\it 제한}이라 한다. 제한 함자
$$
\textit{PMod}(\mathcal{O}_2)
\longrightarrow
\textit{PMod}(\mathcal{O}_1)
$$
를 얻는다.

\medskip\noindent
한편 $\mathcal{O}_1$-가군의 준층 $\mathcal{G}$가 주어지면, 규칙
$$
\left(\mathcal{O}_2 \otimes_{p, \mathcal{O}_1} \mathcal{G}\right)(U)
=
\mathcal{O}_2(U) \otimes_{\mathcal{O}_1(U)} \mathcal{G}(U)
$$
을 통하여 $\mathcal{O}_2$-가군의 준층
$\mathcal{O}_2 \otimes_{p, \mathcal{O}_1} \mathcal{G}$를 구성할 수 있다.
첨자 $p$는 “point(점)”가 아니라 “presheaf(준층)”를 뜻한다.
이 준층을 {\it 텐서곱 준층}이라 한다. {\it 환의 변경} 함자
$$
\textit{PMod}(\mathcal{O}_1)
\longrightarrow
\textit{PMod}(\mathcal{O}_2)
$$
를 얻는다.

\begin{lemma}
\label{sheaves-lemma-adjointness-tensor-restrict-presheaves}
위의 $X$, $\mathcal{O}_1$, $\mathcal{O}_2$, $\mathcal{F}$ 및
$\mathcal{G}$에 대하여 표준적인 전단사
$$
\Hom_{\mathcal{O}_1}(\mathcal{G}, \mathcal{F}_{\mathcal{O}_1})
=
\Hom_{\mathcal{O}_2}(
\mathcal{O}_2 \otimes_{p, \mathcal{O}_1} \mathcal{G},
\mathcal{F}
)
$$
가 존재한다. 다시 말해, 제한 함자와 환의 변경 함자는 서로 수반이다.
\end{lemma}

\begin{proof}
이는 환 사상 $A \to B$에 대하여 제한 함자와 환의 변경 함자가
서로 수반이라는 사실로부터 따른다.
\end{proof}





\section{층}
\label{sheaves-section-sheaves}

\noindent
이 절에서는 층 조건을 설명한다.

\begin{definition}
\label{sheaves-definition-sheaf}
$X$를 위상공간이라 하자.
\begin{enumerate}
\item {\it $X$ 위의 집합의 층 $\mathcal{F}$}란 다음의 추가 성질을
만족하는 집합의 준층이다. 임의의 열린 덮개
$U = \bigcup_{i \in I} U_i$와 단면들의 임의의 모임
$s_i \in \mathcal{F}(U_i)$, $i \in I$가 주어지고
$\forall i, j\in I$에 대하여
$$
s_i|_{U_i \cap U_j} = s_j|_{U_i \cap U_j}
$$
이면, 모든 $i \in I$에 대하여 $s_i = s|_{U_i}$를 만족하는
유일한 단면 $s \in \mathcal{F}(U)$가 존재한다.
\item {\it 집합의 층의 사상}은 단순히 집합의 준층의 사상이다.
\item $X$ 위의 집합의 층들의 범주는 $\Sh(X)$로 나타낸다.
\end{enumerate}
\end{definition}

\begin{remark}
\label{sheaves-remark-confusion}
공집합 $\emptyset \subset X$ 위에서 층의 단면들의 집합에 관해
무언가를 말해야 하는지를 두고 언제나 약간의 혼동이 있다.
말해야 하며, 정의를 올바르게 읽었다면 우리는 이미 그렇게 했다.
실제로 공집합은 빈 열린 덮개로 덮인다는 점에 유의하라.
% P01-ERRATA: P01-E1236 -- the frozen English has subject-verb and missing-preposition defects in this empty-product sentence.
따라서 위 정의의 “단면들의 모임 $s_i$”는 실제로 빈 곱의 원소를 이루는데,
그 빈 곱은 층이 값을 갖는 범주의 끝 대상이다. 다시 말해, 정의를
올바르게 읽으면 자동으로
$\mathcal{F}(\emptyset) = \textit{a final object}$를 얻는다.
집합의 층인 경우 이것은 한원소 집합이다. 이 논증이 마음에 들지 않는다면
그저 $\mathcal{F}(\emptyset) = \{*\}$를 요구해도 된다.

\medskip\noindent
특히 이 조건은 열린집합 $U, V \subset X$가 {\it 서로소}이면
$$
\mathcal{F}(U \cup V) = \mathcal{F}(U) \times \mathcal{F}(V).
$$
임을 보장한다. (끝 대상 위의 섬유곱은 곱이기 때문이다.)
\end{remark}

\begin{example}
\label{sheaves-example-basic-continuous-maps}
$X$, $Y$를 위상공간이라 하자.
열린집합 $U \subset X$에 집합
$$
\mathcal{F}(U) = \{ f : U \to Y \mid f \text{ is continuous}\}
$$
을 명백한 제한 사상들과 함께 대응시키는 규칙 $\mathcal{F}$를 생각하자.
$\mathcal{F}$가 층이라고 주장한다. 이를 보이기 위해
$U = \bigcup_{i\in I} U_i$가 열린 덮개이고
$f_i \in \mathcal{F}(U_i)$, $i\in I$이며, 모든 $i, j \in I$에 대하여
$f_i |_{U_i \cap U_j} = f_j|_{U_i \cap U_j}$라고 하자.
이때 $u \in U_i$를 만족하는 임의의 $i \in I$에 대하여
$f(u)$를 $f_i(u)$의 값으로 두어 $f : U \to Y$를 정의한다.
가정에 따라 이는 잘 정의된다. 또한 $f : U \to Y$는 $U_i$로의 제한이
연속 사상 $f_i$와 일치하는 사상이다. 따라서 $f$는 분명히 연속이다!
\end{example}

\noindent
이 예의 결과를 사용하여 상수층을 정의할 수 있다. 실제로 $A$가
집합이라고 하자. $A$에 이산위상을 주자. $U \subset X$를 열린
부분집합이라 하자. 그러면
$$
\{ f : U \to A \mid f\text{ continuous}\}
=
\{ f : U \to A \mid f\text{ locally constant}\}.
$$
% P01-ERRATA: P01-E1237 -- the frozen English omits the head noun after “an open”.
따라서 각 열린집합에 $A$로 가는 모든 국소상수 사상의 집합을
대응시키는 규칙은 층이다.

\begin{definition}
\label{sheaves-definition-constant-sheaf}
$X$를 위상공간이라 하자. $A$를 집합이라 하자.
$\underline{A}$ 또는 $\underline{A}_X$로 나타내는
{\it 값이 $A$인 상수층}은 열린집합 $U \subset X$에
모든 국소상수 사상 $U \to A$의 집합을 대응시키고, 함수의 제한으로
제한 사상을 주는 층이다.
\end{definition}

\begin{example}
\label{sheaves-example-sheaf-product-pointwise}
$X$를 위상공간이라 하자. $(A_x)_{x \in X}$를 점 $x \in X$로
지표화된 집합들 $A_x$의 족이라 하자. 이 데이터로부터 집합의 층
$\Pi$를 구성할 것이다.
% P01-ERRATA: P01-E1238 -- the frozen English phrase “For U subset X open set” omits required predication.
열린집합 $U \subset X$에 대하여
$$
\Pi(U) = \prod\nolimits_{x \in U} A_x.
$$
로 둔다. 열린집합 $V \subset U \subset X$에 대하여 다음 규칙으로
제한 사상을 정의한다. 원소
$s = (a_x)_{x\in U} \in \Pi(U)$의 제한은
$s|_V = (a_x)_{x \in V}$이다. 이것이 집합의 준층을 정의함은 명백하다.
이것이 층이라고 주장한다. 실제로 $U = \bigcup U_i$가 열린 덮개라고 하자.
$s_i \in \Pi(U_i)$이고 $s_i$와 $s_j$가 $U_i \cap U_j$ 위에서
일치한다고 하자. $s_i = (a_{i, x})_{x\in U_i}$로 쓰자. 양립 조건은
$x \in U_i \cap U_j$일 때마다 집합 $A_x$ 안에서
$a_{i, x} = a_{j, x}$임을 뜻한다. 따라서 어떤 $i$에 대하여
$x \in U_i$일 때마다 $a_x = a_{i, x}$를 만족하는 유일한 원소
$s = (a_x)_{x\in U}$가
$\Pi(U) = \prod_{x\in U} A_x$ 안에 존재한다. 물론 이 원소 $s$는
모든 $i$에 대하여 $s|_{U_i} = s_i$를 만족한다.
\end{example}

\begin{example}
\label{sheaves-example-direct-sum-points-not-sheaf}
$X$를 위상공간이라 하자.
각 $x\in X$에 대하여 아벨군 $M_x$가 주어졌다고 하자.
예 \ref{sheaves-example-direct-sum-points}에서 정의한 준층
$\mathcal{F} : U \mapsto \bigoplus_{x \in U} M_x$를 생각하자.
이는 일반적으로 층이 아니다. 예를 들어 $X$가 이산위상을 갖는
무한집합이면, 층 조건은
$\mathcal{F}(X) = \prod_{x\in X} \mathcal{F}(\{x\})$임을 함의하지만,
정의에 따르면
$\mathcal{F}(X)
= \bigoplus_{x \in X} M_x = \bigoplus_{x \in X} \mathcal{F}(\{x\})$이다.
그리고 무한 직합은 일반적으로 무한 직적과 다르다.

\medskip\noindent
그러나 $X$가, $X$의 모든 열린집합이 준콤팩트인 위상공간이면
$\mathcal{F}$는 실제로 층이다. 이는 독자를 위한 연습문제로 남긴다.
\end{example}



\section{아벨층}
\label{sheaves-section-abelian-sheaves}

\begin{definition}
\label{sheaves-definition-abelian-sheaf}
$X$를 위상공간이라 하자.
\begin{enumerate}
\item {\it $X$ 위의 아벨층} 또는
{\it $X$ 위의 아벨군의 층}이란, 그 바탕 집합의 준층이 층인
$X$ 위의 아벨 준층이다.
\item 아벨군의 층들의 범주는 $\textit{Ab}(X)$로 나타낸다.
\end{enumerate}
\end{definition}

\noindent
$X$를 위상공간이라 하자. 아벨 준층 $\mathcal{F}$의 경우, 열린 덮개
$U = \bigcup U_i$에 대한 층 조건은 흔히 아벨군의 복합체
$$
0 \to \mathcal{F}(U)
\to \prod\nolimits_i \mathcal{F}(U_i)
\to \prod\nolimits_{(i_0, i_1)} \mathcal{F}(U_{i_0} \cap U_{i_1})
$$
가 완전하다고 표현한다. 첫 번째 사상은 통상적인 사상이고,
두 번째 사상은 원소 $(s_i)_{i \in I}$를 원소
$$
(
s_{i_0}|_{U_{i_0} \cap U_{i_1}} -
s_{i_1}|_{U_{i_0} \cap U_{i_1}}
)_{(i_0, i_1)}
\in \prod\nolimits_{(i_0, i_1)} \mathcal{F}(U_{i_0} \cap U_{i_1})
$$
로 보낸다.

\section{대수 구조의 층}
\label{sheaves-section-sheaves-structures}

\noindent
몇 가지 종류의 구조를 갖는 층의 정의를 명확히 하자.
먼저 층 조건을 다시 표현하자. $\mathcal{F}$가 위상공간 $X$ 위의
집합의 준층이라고 하자. 층 조건은 다음과 같이 다시 쓸 수 있다.
$U = \bigcup_{i\in I} U_i$를 열린 덮개라 하고, 도식
$$
\xymatrix{
\mathcal{F}(U) \ar[r]
&
\prod\nolimits_{i\in I}
\mathcal{F}(U_i)
\ar@<1ex>[r] \ar@<-1ex>[r]
&
\prod\nolimits_{(i_0, i_1) \in I \times I}
\mathcal{F}(U_{i_0} \cap U_{i_1})
}
$$
을 생각하자. 여기서 왼쪽 사상은 규칙
% P01-ERRATA: P01-E1239 -- the frozen source uses product operators where tuple/family notation is required.
$s \mapsto \prod_{i \in I} s|_{U_i}$로 정의된다. 오른쪽의 두 사상은
$$
\prod\nolimits_i s_i
\mapsto
\prod\nolimits_{(i_0, i_1)} s_{i_0}|_{U_{i_0} \cap U_{i_1}}
\text{ resp. }
\prod\nolimits_i s_i
\mapsto
\prod\nolimits_{(i_0, i_1)} s_{i_1}|_{U_{i_0} \cap U_{i_1}}.
$$
이다. 층 조건은 정확히 왼쪽 화살표가 오른쪽 두 화살표의 등화자임을
말한다. 이 설명은 그 범주에 곱이 존재하기만 하면 범주에 값을 갖는
준층의 경우로 곧바로 일반화된다.

\begin{definition}
\label{sheaves-definition-sheaf-values-in-category}
$X$를 위상공간이라 하자. $\mathcal{C}$를 곱을 갖는 범주라 하자.
$X$ 위에서 $\mathcal{C}$에 값을 갖는 준층 $\mathcal{F}$가
{\it 층}이라는 것은, 모든 열린 덮개에 대하여 도식
$$
\xymatrix{
\mathcal{F}(U) \ar[r]
&
\prod\nolimits_{i\in I}
\mathcal{F}(U_i)
\ar@<1ex>[r] \ar@<-1ex>[r]
&
\prod\nolimits_{(i_0, i_1) \in I \times I}
\mathcal{F}(U_{i_0} \cap U_{i_1})
}
$$
이 범주 $\mathcal{C}$에서 등화자 도식이라는 뜻이다.
\end{definition}

\noindent
$\mathcal{C}$가 범주이고
$F : \mathcal{C} \to \textit{Sets}$가 충실 함자라고 하자.
염두에 둘 좋은 예는 $\mathcal{C}$가 아벨군의 범주이고
$F$가 망각 함자인 경우이다.
$\mathcal{F}$를 $X$ 위에서 $\mathcal{C}$에 값을 갖는 준층이라 하자.
위 조건을 바탕 집합의 준층
(정의 \ref{sheaves-definition-underlying-presheaf-sets})을 사용하여
다시 표현하고자 한다. 바탕 집합의 준층이 집합의 층일 필요충분조건은
집합의 모든 도식
$$
\xymatrix{
F(\mathcal{F}(U)) \ar[r]
&
\prod\nolimits_{i\in I}
F(\mathcal{F}(U_i))
\ar@<1ex>[r] \ar@<-1ex>[r]
&
\prod\nolimits_{(i_0, i_1) \in I \times I}
F(\mathcal{F}(U_{i_0} \cap U_{i_1}))
}
$$
이 망각 함자 $F$를 적용한 뒤에도 등화자 도식인 것이다.
따라서 $\mathcal{C}$에는 곱과 등화자가 존재하고 $F$가 이들과
가환하기를 원한다. 이는 $\mathcal{C}$가 극한을 갖고 $F$가 극한과
가환한다는 조건과 동치이다. Categories, 보조정리
\ref{categories-lemma-limits-products-equalizers}를 보라.
그러나 이것만으로는 아직 충분하지 않다
(예 \ref{sheaves-example-sheaves-topological-spaces} 참조).
또한 $F$가 {\it 동형사상을 반영}해야 한다. 이 성질은
$\mathcal{C}$의 사상 $f : A \to A'$가 주어졌을 때,
$F(f)$가 전단사이면 (그리고 그때에만) $f$가 동형사상이라는 뜻이다.

\begin{lemma}
\label{sheaves-lemma-sheaves-structure}
범주 $\mathcal{C}$와 함자
$F : \mathcal{C} \to \textit{Sets}$가 다음 성질들을 갖는다고 하자.
\begin{enumerate}
\item $F$는 충실하고,
\item $\mathcal{C}$는 극한을 가지며 $F$는 극한과 가환하고,
\item 함자 $F$는 동형사상을 반영한다.
\end{enumerate}
$X$를 위상공간이라 하자. $\mathcal{F}$를 $\mathcal{C}$에 값을 갖는
준층이라 하자. 그러면 $\mathcal{F}$가 층일 필요충분조건은 그 바탕
집합의 준층이 층인 것이다.
\end{lemma}

\begin{proof}
$\mathcal{F}$가 층이라고 가정하자. 그러면 $\mathcal{F}(U)$는
위 도식의 등화자이고, 가정에 따라 $F(\mathcal{F}(U))$는 대응하는
집합의 도식의 등화자이다. 따라서 $F(\mathcal{F})$는 집합의 층이다.

\medskip\noindent
$F(\mathcal{F})$가 층이라고 가정하자.
$E \in \Ob(\mathcal{C})$를 정의
\ref{sheaves-definition-sheaf-values-in-category}의 평행한 두 화살표의
등화자라 하자. $\mathcal{F}$가 준층이라는 사실만으로도 표준적인 사상
$\mathcal{F}(U) \to E$를 얻는다. 가정에 따라 유도된 사상
$F(\mathcal{F}(U)) \to F(E)$는 동형사상이다. 실제로 $F(E)$는
대응하는 집합의 도식의 등화자이다. 그러므로 보조정리의 조건 (3)에 의해
$\mathcal{F}(U) \to E$가 동형사상임을 알 수 있다.
\end{proof}

\noindent
이 보조정리는 특히
{\it 군, 환, 고정된 환 위의 대수, 고정된 환 위의 가군,
고정된 체 위의 벡터공간 등의 층}에 적용된다.
다시 말해, 이들은 그 바탕 집합의 준층이 층인, 군, 환, 고정된 환 위의
가군, 고정된 체 위의 벡터공간 등의 준층이다.

\begin{example}
\label{sheaves-example-C0-sheaf-rings}
$X$를 위상공간이라 하자. 각 열린집합 $U \subset X$에 대하여
$\mathbf{R}$-대수
$\mathcal{C}^{0}(U) = \{ f : U \to \mathbf{R} \mid f\text{ is continuous}\}$
를 생각하자. 명백한 제한 사상들을 주면 이는 $X$ 위의
$\mathbf{R}$-대수의 준층이 된다. 예
\ref{sheaves-example-basic-continuous-maps}에 따라 이는 집합의 층이다.
따라서 보조정리 \ref{sheaves-lemma-sheaves-structure}에 따라 이는
$X$ 위의 $\mathbf{R}$-대수의 층이다.
\end{example}

\begin{example}
\label{sheaves-example-sheaves-topological-spaces}
위상공간의 범주 $\textit{Top}$을 생각하자.
곱과 등화자와 가환하는 자연스러운 충실 함자
$\textit{Top} \to \textit{Sets}$가 존재한다. 그러나 이 함자는
동형사상을 반영하지 않는다. 실제로 보조정리
\ref{sheaves-lemma-sheaves-structure}의 유사 명제는 성립하지 않는다.
$X = \mathbf{N}$에 이산위상을 주었다고 하자.
% P01-ERRATA: P01-E1240 -- the frozen English omits a comma around its interrupting index phrase.
각 $i \in \mathbf{N}$에 대하여 $A_i$를 이산위상공간이라 하자.
임의의 부분집합 $U \subset \mathbf{N}$에 대하여
$\mathcal{F}(U) = \prod_{i\in U} A_i$에 이산위상을 주어 정의하자.
그러면 이것은 바탕 집합의 준층이 층인 위상공간의 준층이다.
예 \ref{sheaves-example-sheaf-product-pointwise}를 보라. 그러나 각 $A_i$가
적어도 두 원소를 가지면 정의
\ref{sheaves-definition-sheaf-values-in-category}에 따라 이것은
위상공간의 층이 아니다. 각
$\mathcal{F}(U) = \prod_{i\in U} A_i$에 {\it 곱위상}을 주면 실제로
$X$ 위의 위상공간의 층이 된다는 것은 독자가 확인할 수 있다.
\end{example}


\section{가군의 층}
\label{sheaves-section-sheaves-modules}

\begin{definition}
\label{sheaves-definition-sheaf-modules}
$X$를 위상공간이라 하자.
$\mathcal{O}$를 $X$ 위의 환의 층이라 하자.
\begin{enumerate}
\item {\it $\mathcal{O}$-가군의 층}이란 $\mathcal{O}$-가군의
준층 $\mathcal{F}$로서(정의 \ref{sheaves-definition-presheaf-modules} 참조),
그 바탕 아벨군의 준층 $\mathcal{F}$가 층인 것이다.
\item {\it $\mathcal{O}$-가군의 층의 사상}은
$\mathcal{O}$-가군의 준층의 사상이다.
% P01-ERRATA: P01-E1316 -- the frozen English says “set of morphism” instead of “set of morphisms”.
\item $\mathcal{O}$-가군의 층 $\mathcal{F}$와 $\mathcal{G}$가 주어졌을 때,
$\Hom_\mathcal{O}(\mathcal{F}, \mathcal{G})$로
$\mathcal{O}$-가군의 층의 사상들의 집합을 나타낸다.
\item $\mathcal{O}$-가군의 층들의 범주는
$\textit{Mod}(\mathcal{O})$로 나타낸다.
\end{enumerate}
\end{definition}

\noindent
이 정의는 $\mathcal{O}$가 단지 환의 준층인 경우에도 어느 정도 의미가 있다.
하지만 이것이 유용한 예는 알지 못하며, $\mathcal{O}$가 환의 층이 아닌
경우에는 “$\mathcal{O}$-가군의 층”이라는 용어를 사용하지 않을 것이다.







\section{줄기}
\label{sheaves-section-stalks}

\noindent
$X$를 위상공간이라 하자. $x \in X$를 점이라 하자.
$\mathcal{F}$를 $X$ 위의 집합의 준층이라 하자.
{\it $x$에서 $\mathcal{F}$의 줄기}란 집합
$$
\mathcal{F}_x
=
\colim_{x\in U} \mathcal{F}(U)
$$
이다. 여기서 여극한은 $X$에서 $x$의 열린 근방 $U$들의 집합에
대하여 취한다. 열린 근방들의 집합에는 (역)포함관계로 부분순서를 준다.
$U \geq U' \Leftrightarrow U \subset U'$라고 한다.
이 계의 전이 사상들은 $\mathcal{F}$의 제한 사상들로 주어진다.
계에 대한 (여)극한의 표기와 용어는 Categories, 절
\ref{categories-section-posets-limits}를 보라. 이 여극한은 유향
여극한임에 유의하라. 따라서 $\mathcal{F}_x$를 쉽게 기술할 수 있다.
구체적으로,
$$
\mathcal{F}_x
=
\{
(U, s)
\mid
x\in U, s\in \mathcal{F}(U)
\}/\sim
$$
이고, 동치관계는 다음과 같이 주어진다.
$(U, s) \sim (U', s')$일 필요충분조건은
$x \in U''$인 열린집합 $U'' \subset U \cap U'$가 존재하여
$s|_{U''} = s'|_{U''}$인 것이다.
% P01-ERRATA: P01-E1241 -- the frozen English omits “by” in its notation sentence.
쌍 $(U, s)$가 주어지면, $(U, s)$의 동치류에 대응하는
$\mathcal{F}_x$의 원소를 때때로 $s_x$로 나타낸다.
때로는 $s_x$를 가리켜 “$\mathcal{F}_x$에서 $s$의 상”이라는 표현을
사용한다. 예를 들어 두 쌍 $(U, s)$와 $(U', s')$가 주어졌을 때,
$s_x = s'_x$임을 나타내기 위해 “$\mathcal{F}_x$에서 $s$와 $s'$가
같다”라고 말하기도 한다. 다른 저자들은 “$x$에서 $s$의 {\it 싹}”이라는
용어를 사용한다.

\medskip\noindent
이 정의의 명백한 귀결로서, 임의의 열린집합 $U \subset X$에 대하여
표준적인 사상
$$
\mathcal{F}(U)
\longrightarrow
\prod\nolimits_{x \in U} \mathcal{F}_x
$$
이 존재한다. 이는
% P01-ERRATA: P01-E1242 -- the frozen source uses a product operator instead of tuple/family notation.
$s \mapsto \prod_{x \in U} (U, s)$로 정의된다. 생각해 보라!

\begin{lemma}
\label{sheaves-lemma-sheaf-subset-stalks}
$\mathcal{F}$를 위상공간 $X$ 위의 집합의 층이라 하자.
모든 열린집합 $U \subset X$에 대하여 사상
$$
\mathcal{F}(U)
\longrightarrow
\prod\nolimits_{x \in U} \mathcal{F}_x
$$
은 단사이다.
\end{lemma}

\begin{proof}
$s, s' \in \mathcal{F}(U)$가 모든 $x \in U$에 대하여 각 줄기
$\mathcal{F}_x$에서 같은 원소로 간다고 하자. 이는 모든 $x \in U$에
대하여 $x \in V^x$이고 $s|_{V^x} = s'|_{V^x}$인 열린집합
$V^x \subset U$가 존재한다는 뜻이다. 그런데
$U = \bigcup_{x \in U} V^x$는 열린 덮개이다. 따라서 층 조건의
유일성에 의해 $s = s'$임을 알 수 있다.
\end{proof}

\begin{definition}
\label{sheaves-definition-separated}
$X$를 위상공간이라 하자. $X$ 위의 집합의 준층 $\mathcal{F}$가
{\it 분리되어 있다}는 것은, 모든 열린집합 $U \subset X$에 대하여 사상
$\mathcal{F}(U) \to \prod_{x \in U} \mathcal{F}_x$가 단사라는 뜻이다.
\end{definition}

\noindent
또 하나의 관찰은 줄기 $\mathcal{F}_x$의 구성이 준층 $\mathcal{F}$에
대하여 함자적이라는 것이다. 다시 말해, 이는 함자
$$
\textit{PSh}(X) \longrightarrow \textit{Sets},
\ \mathcal{F} \longmapsto \mathcal{F}_x.
$$
를 준다. 이 함자를 {\it 줄기 함자}라 한다. 즉,
$\varphi : \mathcal{F} \to \mathcal{G}$가 준층의 사상이면
$\varphi_x : \mathcal{F}_x \to \mathcal{G}_x$를 규칙
$(U, s) \mapsto (U, \varphi(s))$로 정의한다.
이것이 잘 작동함을 보이려면, $\mathcal{F}_x$에서
$(U, s) = (U', s')$이면 $\mathcal{G}_x$에서도
$(U, \varphi(s)) = (U', \varphi(s'))$임을 확인해야 한다.
$\varphi$가 제한 사상들과 양립하므로 이는 명백하다.

\begin{example}
\label{sheaves-example-stalk-constant-presheaf}
$X$를 위상공간이라 하자. $A$를 집합이라 하자.
값이 $A$인 상수 준층을 잠시 $A_p$로 나타내자
($p$는 point가 아니라 presheaf를 뜻한다).
값이 $A$인 상수층으로 가는 준층의 표준적인 사상
$A_p \to \underline{A}$가 있다. 모든 점에 대하여 표준적인 전단사
$A = (A_p)_x = \underline{A}_x$가 있으며, 여기서 두 번째 사상은
$A_p \to \underline{A}$로부터 함자성에 의해 유도된다.
\end{example}



\begin{example}
\label{sheaves-example-germs-functions}
$X = \mathbf{R}^n$에 유클리드 위상을 주었다고 하자.
$X$ 위의 $\mathcal{C}^\infty$ 함수들의 준층을 생각하고
$\mathcal{C}^\infty_{\mathbf{R}^n}$으로 나타내자. 다시 말해,
$\mathcal{C}^\infty_{\mathbf{R}^n}(U)$는
$\mathcal{C}^\infty$ 함수 $f : U \to \mathbf{R}$들의 집합이다.
예 \ref{sheaves-example-basic-continuous-maps}에서와 같이 이것이 층임을
쉽게 보일 수 있다. 실제로 이것은 $\mathbf{R}$-벡터공간의 층이다.

\medskip\noindent
이제 $x \in X = \mathbf{R}^n$를 점이라 하자. 줄기
$\mathcal{C}^\infty_{\mathbf{R}^n, x}$의 원소를 어떻게 생각해야 하는가?
그러한 원소는 정의역이 $x$를 포함하는 $\mathcal{C}^\infty$ 함수 $f$로
주어진다.
% P01-ERRATA: P01-E1243 -- the frozen English has singular “pair” with plural “determine”.
그러한 함수들의 한 쌍 $f$, $g$는 $x$의 한 근방에서 일치할 때
줄기의 같은 원소를 정한다. 다시 말해,
$\mathcal{C}^\infty_{\mathbf{R}^n, x}$의 원소는 때때로
{\it $x$에서 $\mathcal{C}^\infty$ 함수의 싹}이라고 부르는 것과 같다.
\end{example}

\begin{example}
\label{sheaves-example-sheaf-product-pointwise-stalk}
$X$를 위상공간이라 하자. 각 $x \in X$에 대하여 $A_x$를 집합이라 하자.
예 \ref{sheaves-example-sheaf-product-pointwise}의 층
$\mathcal{F} : U \mapsto \prod_{x\in U} A_x$를 생각하자.
여기서는 $x$에서 $\mathcal{F}$의 줄기 $\mathcal{F}_x$가 일반적으로
집합 $A_x$와 같지 {\it 않다}는 점을 지적하고자 한다. 물론 사상
$\mathcal{F}_x \to A_x$가 있지만 일반적으로 이것이 말할 수 있는 전부이다.
예를 들어 모든 $n \not = m$에 대하여 $x_n \not = x_m$이고
$x = \lim x_n$이며, 모든 $y \in X$에 대하여
$A_y = \{0, 1\}$이라고 하자. 그러면 $\mathcal{F}_x$는 $0$과 $1$로
이루어진 수열들의 꼬리들의 (무한) 집합 위로 전사한다.
실제로 $x$의 모든 열린 근방은 거의 모든 $x_n$을 포함한다.
한편 $x$의 모든 근방이 $A_y = \emptyset$인 점 $y$를 포함하면
$\mathcal{F}_x = \emptyset$이다.
\end{example}



\section{아벨 준층의 줄기}
\label{sheaves-section-stalks-abelian-presheaves}

\noindent
먼저 일반적인 경우의 모형으로서 아벨군의 경우를 다룬다.

\begin{lemma}
\label{sheaves-lemma-stalk-abelian-presheaf}
$X$를 위상공간이라 하자. $\mathcal{F}$를 $X$ 위의 아벨군의
준층이라 하자. 모든 열린집합 $U \subset X$, $x\in U$에 대하여 사상
$\mathcal{F}(U) \to \mathcal{F}_x$가 군 준동형이 되도록 하는
$\mathcal{F}_x$ 위의 아벨군 구조가 유일하게 존재한다. 또한
$$
\mathcal{F}_x
=
\colim_{x\in U} \mathcal{F}(U)
$$
가 아벨군의 범주에서 성립한다.
\end{lemma}

\begin{proof}
원소들의 쌍 $(U, s)$와 $(V, t)$의 합을
쌍 $(U \cap V, s|_{U\cap V} +
t|_{U \cap V})$로 정의한다. 나머지는 쉽게 확인된다.
\end{proof}

\noindent
위 증명에서 핵심은 열린 근방들의 부분순서집합이 유향 집합이라는
사실이다(Categories, 정의
\ref{categories-definition-directed-set} 참조). 실제로 두 아벨군
$A, B$의 쌍대곱은 직합 $A \oplus B$이지만, 집합의 범주에서
쌍대곱은 서로소 합집합 $A \amalg B$이다. 이는 일반적으로
아벨군의 범주의 여극한이 집합의 범주의 여극한과 일치하지 않음을 보여 준다.


\section{대수 구조의 준층의 줄기}
\label{sheaves-section-stalks-presheaves-structures}

\noindent
보조정리 \ref{sheaves-lemma-stalk-abelian-presheaf}의 증명은 유향 여극한이
망각 함자와 가환하는 모든 종류의 대수 구조에 적용된다.

\begin{lemma}
\label{sheaves-lemma-stalk-presheaf-values-in-category}
$\mathcal{C}$를 범주라 하자.
$F : \mathcal{C} \to \textit{Sets}$를 함자라 하자.
다음을 가정한다.
\begin{enumerate}
\item $F$는 충실하고,
\item $\mathcal{C}$에는 유향 여극한이 존재하며 $F$는 그것들과 가환한다.
\end{enumerate}
$X$를 위상공간이라 하자. $x \in X$라 하자.
$\mathcal{F}$를 $\mathcal{C}$에 값을 갖는 준층이라 하자. 그러면
$$
\mathcal{F}_x = \colim_{x\in U} \mathcal{F}(U)
$$
가 $\mathcal{C}$에서 존재한다. 그 바탕 집합은 $\mathcal{F}$의
바탕 집합의 준층의 줄기와 같다. 또한 구성
$\mathcal{F} \mapsto \mathcal{F}_x$는 $\mathcal{C}$에 값을 갖는
준층의 범주에서 $\mathcal{C}$로 가는 함자이다.
\end{lemma}

\begin{proof}
생략한다.
\end{proof}

\noindent
정의 자체에 의해 모든 사상
$\mathcal{F}(U) \to \mathcal{F}_x$는 범주 $\mathcal{C}$의 사상이며,
망각 함자 $F$를 적용하면 바탕 집합의 층에 대한 대응하는 사상이 된다.
늘 그렇듯이 $\mathcal{C}$의 사상과 그 바탕 집합의 사상을 구별하지
않을 것이다. $F$가 충실하므로 이렇게 해도 된다.

\medskip\noindent
이 보조정리는 특히 다음에 적용된다.
{\it (반드시 아벨일 필요는 없는) 군, 환, 고정된 환 위의 가군,
고정된 체 위의 벡터공간의 준층}.

\section{가군 준층의 줄기}
\label{sheaves-section-stalk-presheaves-modules}

\begin{lemma}
\label{sheaves-lemma-stalk-module}
$X$를 위상공간이라 하자.
$\mathcal{O}$를 $X$ 위의 환의 준층이라 하자.
$\mathcal{F}$를 $\mathcal{O}$-가군의 준층이라 하자.
$x \in X$라 하자.
% P01-ERRATA: P01-E1244 -- the frozen source calls the module action an internal multiplication map.
곱셈 사상 $\mathcal{O} \times \mathcal{F} \to \mathcal{F}$로부터 오는
표준적인 사상 $\mathcal{O}_x \times \mathcal{F}_x
\to \mathcal{F}_x$는 아벨군 $\mathcal{F}_x$ 위에
$\mathcal{O}_x$-가군 구조를 정의한다.
\end{lemma}

\begin{proof}
생략한다.
\end{proof}

\begin{lemma}
\label{sheaves-lemma-stalk-tensor-presheaf-modules}
$X$를 위상공간이라 하자.
$\mathcal{O} \to \mathcal{O}'$를 $X$ 위의 환의 준층의 사상이라 하자.
$\mathcal{F}$를 $\mathcal{O}$-가군의 준층이라 하자.
$x \in X$라 하자. 다음 등식이 성립한다.
$$
\mathcal{F}_x \otimes_{\mathcal{O}_x} \mathcal{O}'_x
=
(\mathcal{F} \otimes_{p, \mathcal{O}} \mathcal{O}')_x
$$
이는 $\mathcal{O}'_x$-가군으로서의 등식이다.
\end{lemma}

\begin{proof}
생략한다.
\end{proof}


\section{대수 구조}
\label{sheaves-section-algebraic-structures}

\noindent
이 절에서는 앞의 절들에서 만난 개념들을 가볍게 형식화한다.

\begin{definition}
\label{sheaves-definition-algebraic-structure}
{\it 대수 구조의 유형}은 범주 $\mathcal{C}$와 다음 성질들을 갖는
함자 $F : \mathcal{C} \to \textit{Sets}$로 주어진다.
\begin{enumerate}
\item $F$는 충실하고,
\item $\mathcal{C}$는 극한을 가지며 $F$는 극한과 가환하고,
\item $\mathcal{C}$는 여과 여극한을 가지며 $F$는 그것들과 가환하고,
\item $F$는 동형사상을 반영한다.
\end{enumerate}
\end{definition}

\noindent
우리는 아래의 여러 논증에서 사용할 성질들을 드러내기 위하여 이 정의를
도입한다. 다만 Categories, 주석
\ref{categories-remark-big-categories}에 열거된 것들을 제외하면,
관례상 “큰” 범주를 연구할 수 없으므로 이 개념 자체를 아주 자세히
연구하지는 않을 것이다. 열거된 범주 가운데 다음 범주들은 필요한
성질을 갖는다.

\begin{lemma}
\label{sheaves-lemma-list-algebraic-structures}
다음 범주들에 명백한 망각 함자를 주면 대수 구조의 유형들이 정의된다.
\begin{enumerate}
\item 기점이 지정된 집합들의 범주.
\item 아벨군들의 범주.
\item 군들의 범주.
\item 모노이드들의 범주.
\item 환들의 범주.
\item 고정된 환 $R$에 대한 $R$-가군들의 범주.
\item 고정된 체 위의 리 대수들의 범주.
\end{enumerate}
\end{lemma}

\begin{proof}
생략한다.
\end{proof}

\noindent
이제부터 대수 구조의 (준)층과 그 줄기는 바탕 집합의 (준)층을 통하여
생각할 것이다. 보조정리 \ref{sheaves-lemma-sheaves-structure}와
\ref{sheaves-lemma-stalk-presheaf-values-in-category}에 따라 이렇게 해도 된다.

\medskip\noindent
이 절의 나머지 부분에서는 앞으로 유용하게 쓸 대수 구조에 관한 몇 가지
결과를 지적한다.

\begin{lemma}
\label{sheaves-lemma-properties-algebraic-structures}
$(\mathcal{C}, F)$를 대수 구조의 유형이라 하자.
\begin{enumerate}
\item $\mathcal{C}$는 끝 대상 $0$을 가지며 $F(0) = \{ * \}$이다.
\item $\mathcal{C}$는 곱을 가지며
$F(\prod A_i) = \prod F(A_i)$이다.
\item $\mathcal{C}$는 섬유곱을 가지며
$F(A \times_B C) = F(A)\times_{F(B)}F(C)$이다.
\item $\mathcal{C}$는 등화자를 가진다. $E \to A$가
$a, b : A \to B$의 등화자이면 $F(E) \to F(A)$는
$F(a), F(b) : F(A) \to F(B)$의 등화자이다.
\item $A \to B$가 단사사상일 필요충분조건은
$F(A) \to F(B)$가 단사인 것이다.
\item $F(a) : F(A) \to F(B)$가 전사이면 $a$는 전사사상이다.
\item $A_1 \to A_2 \to A_3 \to \ldots$가 주어지면
$\colim A_i$가 존재하고 $F(\colim A_i) = \colim F(A_i)$이며,
더 일반적으로 모든 여과 여극한에 대해서도 마찬가지이다.
\end{enumerate}
\end{lemma}

\begin{proof}
생략한다. 흥미로운 명제는 (5)뿐이다. 이는 $A \to B$가
단사사상일 필요충분조건이 $A \to A \times_B A$가 동형사상이라는
사실과, 이어서 $F$가 동형사상을 반영한다는 사실을 적용하면 따른다.
\end{proof}

\begin{lemma}
\label{sheaves-lemma-image-contained-in}
$(\mathcal{C}, F)$를 대수 구조의 유형이라 하자.
$A, B, C \in \Ob(\mathcal{C})$라고 하자.
$f : A \to B$와 $g : C \to B$를 $\mathcal{C}$의 사상이라 하자.
$F(g)$가 단사이고 $\Im(F(f)) \subset \Im(F(g))$이면,
어떤 사상 $t : A \to C$에 대하여 $f$는
$f = g \circ t$로 분해된다.
\end{lemma}

\begin{proof}
$A \times_B C$를 생각하자. 가정은
$F(A \times_B C) = F(A) \times_{F(B)} F(C) = F(A)$임을 함의한다.
% P01-ERRATA: P01-E1245 -- reflection proves an isomorphism of the canonical projection, not literal object equality.
따라서 $F$가 동형사상을 반영하므로 $A = A \times_B C$이다.
결론이 따른다.
\end{proof}

\begin{example}
\label{sheaves-example-application-lemma-image-contained-in}
이 보조정리는 다음 상황에 자주 적용될 것이다. $\mathcal{C}$의 도식
$$
\xymatrix{
A \ar[r] & B \ar[d] \\
C \ar[r] & D
}
$$
이 주어졌다고 하자. $C \to D$가 바탕 집합에서 단사이고,
합성 $A \to B \to D$의 바탕 집합에서의 상이
$C \to D$의 상 안에 있다고 하자. 그러면 $\mathcal{C}$에서
가환 도식
$$
\xymatrix{
A \ar[r] \ar[d] & B \ar[d] \\
C \ar[r] & D
}
$$
을 얻는다.
\end{example}

\begin{example}
\label{sheaves-example-sheaf-product-pointwise-algebraic-structure}
% P01-ERRATA: P01-E0039 -- the frozen sentence calls F, rather than the pair (C,F), a type of algebraic structure.
$F : \mathcal{C} \to \textit{Sets}$를 대수 구조의 유형이라 하자.
$X$를 위상공간이라 하자. 모든 $x \in X$에 대하여 대상
$A_x \in \Ob(\mathcal{C})$가 주어졌다고 하자.
$X$ 위에서 $\mathcal{C}$에 값을 갖고 규칙
$\Pi(U) = \prod_{x \in U} A_x$로 정의되는 준층 $\Pi$를 생각하자
(제한 사상은 명백하다). 연관된 집합의 준층
$U \mapsto F(\Pi(U)) = \prod_{x \in U} F(A_x)$는 예
\ref{sheaves-example-sheaf-product-pointwise}에 따라 층임에 유의하라.
따라서 $\Pi$는 유형 $(\mathcal{C} , F)$의 대수 구조의 층이다.
이로부터 아벨군, 군, 환 등의 층의 많은 예를 얻는다.
\end{example}













\section{완전성과 점}
\label{sheaves-section-exactness-points}

\noindent
임의의 범주에는 전사사상, 단사사상, 동형사상 등의 개념이 있다.

\begin{lemma}
\label{sheaves-lemma-points-exactness}
$X$를 위상공간이라 하자.
$\varphi : \mathcal{F} \to \mathcal{G}$를 $X$ 위의 집합의 층의
사상이라 하자.
\begin{enumerate}
\item 사상 $\varphi$가 층의 범주에서 단사사상일 필요충분조건은
모든 $x \in X$에 대하여 사상
$\varphi_x : \mathcal{F}_x \to \mathcal{G}_x$가 단사인 것이다.
\item 사상 $\varphi$가 층의 범주에서 전사사상일 필요충분조건은
모든 $x \in X$에 대하여 사상
$\varphi_x : \mathcal{F}_x \to \mathcal{G}_x$가 전사인 것이다.
\item 사상 $\varphi$가 층의 범주에서 동형사상일 필요충분조건은
모든 $x \in X$에 대하여 사상
$\varphi_x : \mathcal{F}_x \to \mathcal{G}_x$가 전단사인 것이다.
\end{enumerate}
\end{lemma}

\begin{proof}
생략한다.
\end{proof}

\noindent
따라서 집합의 층의 범주에서 전사사상과 단사사상의 개념은 다음과 같이
기술할 수 있다.

\begin{definition}
\label{sheaves-definition-injective-surjective}
$X$를 위상공간이라 하자.
\begin{enumerate}
\item 준층 $\mathcal{F}$가 준층 $\mathcal{G}$의 {\it 부분준층}이라는
것은 모든 열린집합 $U \subset X$에 대하여
$\mathcal{F}(U) \subset \mathcal{G}(U)$이고 $\mathcal{G}$의
제한 사상들이 $\mathcal{F}$의 제한 사상들을 유도한다는 뜻이다.
$\mathcal{F}$와 $\mathcal{G}$가 층이면 $\mathcal{F}$를
$\mathcal{G}$의 {\it 부분층}이라 한다. 이를 때때로
$\mathcal{F} \subset \mathcal{G}$로 나타낸다.
\item $X$ 위의 집합의 준층의 사상
$\varphi : \mathcal{F} \to \mathcal{G}$가 {\it 단사}라는 것은
모든 $X$의 열린집합 $U$에 대하여
$\mathcal{F}(U) \to \mathcal{G}(U)$가 단사라는 것과 동치이다.
\item $X$ 위의 집합의 준층의 사상
$\varphi : \mathcal{F} \to \mathcal{G}$가 {\it 전사}라는 것은
모든 $X$의 열린집합 $U$에 대하여
$\mathcal{F}(U) \to \mathcal{G}(U)$가 전사라는 것과 동치이다.
\item $X$ 위의 집합의 층의 사상
$\varphi : \mathcal{F} \to \mathcal{G}$가 {\it 단사}라는 것은
모든 $X$의 열린집합 $U$에 대하여
$\mathcal{F}(U) \to \mathcal{G}(U)$가 단사라는 것과 동치이다.
\item $X$ 위의 집합의 층의 사상
$\varphi : \mathcal{F} \to \mathcal{G}$가 {\it 전사}라는 것은
$X$의 모든 열린집합 $U$와 $\mathcal{G}(U)$의 모든 단면 $s$에 대하여,
$s|_{U_i}$가 모든 $i$에 대해
$\mathcal{F}(U_i) \to \mathcal{G}(U_i)$의 상에 속하도록 하는
열린 덮개 $U = \bigcup U_i$가 존재한다는 것과 동치이다.
\end{enumerate}
\end{definition}


\begin{lemma}
\label{sheaves-lemma-characterize-epi-mono}
$X$를 위상공간이라 하자.
\begin{enumerate}
\item 준층의 범주에서 전사사상(각각 단사사상)은 정확히 준층의
전사 사상(각각 단사 사상)이다.
\item 층의 범주에서 전사사상(각각 단사사상)은 정확히 층의
전사 사상(각각 단사 사상)이며, 정확히 모든 줄기에서
전사(각각 단사)인 사상이다.
\end{enumerate}
\end{lemma}

\begin{proof}
생략한다.
\end{proof}


\begin{lemma}
\label{sheaves-lemma-check-homomorphism-stalks}
$X$를 위상공간이라 하자.
$(\mathcal{C}, F)$를 대수 구조의 유형이라 하자.
$\mathcal{F}$, $\mathcal{G}$가 $X$ 위에서 $\mathcal{C}$에 값을 갖는
층이라고 하자.
$\varphi : \mathcal{F} \to \mathcal{G}$를 그 바탕 집합의 층들의
사상이라 하자. 모든 점 $x \in X$에 대하여 사상
$\mathcal{F}_x \to \mathcal{G}_x$가 대수 구조의 사상이면,
$\varphi$는 대수 구조의 층들의 사상이다.
\end{lemma}

\begin{proof}
$U$를 $X$의 열린 부분집합이라 하자. (바탕) 집합의 도식
$$
\xymatrix{
\mathcal{F}(U) \ar[r] \ar[d] &
\prod_{x \in U} \mathcal{F}_x \ar[d] \\
\mathcal{G}(U) \ar[r] &
\prod_{x \in U} \mathcal{G}_x
}
$$
을 생각하자. 가정과 앞의 결과들에 따라 왼쪽 수직 화살표를 제외한
모든 화살표는 대수 구조의 사상이다. 또한 아래쪽 수평 화살표는
단사이다. 보조정리 \ref{sheaves-lemma-sheaf-subset-stalks}를 보라.
따라서 보조정리 \ref{sheaves-lemma-image-contained-in}로부터 결론을 얻는다.
예 \ref{sheaves-example-application-lemma-image-contained-in}도 보라.
% P01-ERRATA: P01-E1246 -- the frozen proof lacks final punctuation after the last reference.
\end{proof}

\noindent
아벨층의 짧은 완전열 등은 가군의 층을 다루는 장에서 논의할 것이다.
Modules, 절 \ref{modules-section-kernels}를 보라.


\section{층화}
\label{sheaves-section-sheafification}

\noindent
이 절에서는 위상공간 위의 준층을 층화하는 방법을 설명한다.
이 경우에는 줄기를 이용하여 층화를 기술할 것이다. 이는 Sites, 절
\ref{sites-section-sheafification}에서 설명한 일반적인 절차와는 다르며,
아마도 어느 정도 더 이해하기 쉽다.

\medskip\noindent
기본 구성은 다음과 같다. $\mathcal{F}$를 위상공간 $X$ 위의 집합의
준층이라 하자. 각 열린집합 $U \subset X$에 대하여
$$
\mathcal{F}^{\#}(U)
=
\{
(s_u) \in \prod\nolimits_{u \in U} \mathcal{F}_u
\text{ such that }(*)
\}
$$
로 정의한다. 여기서 $(*)$는 다음 성질이다.
\begin{itemize}
\item[(*)] 모든 $u \in U$에 대하여 열린 근방 $u \in V \subset U$와
단면 $\sigma \in \mathcal{F}(V)$가 존재하여, 모든 $v \in V$에 대해
$\mathcal{F}_v$에서 $s_v = (V, \sigma)$가 성립한다.
\end{itemize}
$(*)$는 각 $u \in U$에 대한 조건이며, $u \in U$가 주어졌을 때
이 조건의 참 여부는 $u$의 임의의 열린 근방에 놓인 $v$들에 대한 값
$s_v$에만 의존한다. 따라서 $V \subset U \subset X$가 열린집합이면
사영 사상들
$$
\prod\nolimits_{u \in U} \mathcal{F}_u
\longrightarrow
\prod\nolimits_{v \in V} \mathcal{F}_v
$$
은 $\mathcal{F}^{\#}(U)$의 원소를 $\mathcal{F}^{\#}(V)$의 원소로
보낸다는 것이 분명하다.
% P01-ERRATA: P01-E1247 -- the frozen source has plural subject "projection maps" with singular verb "maps".
이 사상들을 제한 사상으로 사용하면 $\mathcal{F}^\#$는 $X$ 위의
집합의 준층이 된다.

\medskip\noindent
더 나아가, 절 \ref{sheaves-section-stalks}에서 설명한 사상
$\mathcal{F}(U) \to \prod_{u \in U} \mathcal{F}_u$의 상은 분명히
$\mathcal{F}^{\#}(U)$에 포함된다. 또한 $V \subset U \subset X$가
열린집합이면 다음 가환도식이 있다.
$$
\xymatrix{
\mathcal{F}(U) \ar[r] \ar[d] &
\mathcal{F}^{\#}(U) \ar[r] \ar[d] &
\prod_{u\in U} \mathcal{F}_u \ar[d] \\
\mathcal{F}(V) \ar[r] &
\mathcal{F}^{\#}(V) \ar[r] &
\prod_{v\in V} \mathcal{F}_v
}
$$
여기서 수직 사상들은 제한 사상들로부터 유도된다. 따라서 준층의
표준적인 사상 $\mathcal{F} \to \mathcal{F}^{\#}$가 있음을 알 수 있다.

\medskip\noindent
예 \ref{sheaves-example-sheaf-product-pointwise}에서 자명한 제한 사상들을 갖는
규칙 $\Pi(\mathcal{F}) : U \mapsto \prod_{u\in U} \mathcal{F}_u$가
층임을 보았다. 구성에 의해 $\mathcal{F}^{\#}$는 이 층의 부분준층이다.
달리 말하면 준층의 사상들
$$
\mathcal{F} \to \mathcal{F}^\# \to \Pi(\mathcal{F}).
$$
이 있다. 또한 $\mathcal{F}$에 위의 열을 대응시키는 규칙은 준층
$\mathcal{F}$에 대해 분명히 함자적이다. 아래 보조정리들의 증명에서
이 표기를 사용할 것이다.

\begin{lemma}
\label{sheaves-lemma-sheafification-sheaf}
준층 $\mathcal{F}^{\#}$는 층이다.
\end{lemma}

\begin{proof}
독자에게는 여기의 증명을 읽는 것보다 스스로 설명을 찾아내는 편이
더 나을지도 모른다. 실제로 이 보조정리가 참인 이유는 연속함수의
준층이 층인 이유와 같다. 예 \ref{sheaves-example-basic-continuous-maps}를 보라
(그리고 이 비유는 ``espace \'etal\'e''를 사용하여 엄밀하게 만들 수 있다).
% P01-ERRATA: P01-E1248 -- the frozen source says "for the same reason as why".
% P01-ERRATA: P01-E1249 -- the frozen source says "presheaf of continuous function" in the singular.

\medskip\noindent
어쨌든 $U = \bigcup U_i$를 열린 덮개라 하자.
$s_i = (s_{i, u})_{u \in U_i} \in \mathcal{F}^{\#}(U_i)$이고,
$s_i$와 $s_j$가 $U_i \cap U_j$ 위에서 일치한다고 하자.
$\Pi(\mathcal{F})$는 층이므로 $\prod_{u\in U} \mathcal{F}_u$ 안에서,
$U_i$에 대한 제한이 $s_i$인 원소 $s = (s_u)_{u\in U}$를 얻는다.
성질 $(*)$를 확인해야 한다. $u \in U$를 택하자. 그러면 어떤 $i$에
대해 $u \in U_i$이다. 따라서 $s_i$에 대한 $(*)$에 의해 열린집합
$V$와 $u \in V \subset U_i$, 그리고 $\sigma \in \mathcal{F}(V)$가
존재하여, 모든 $v \in V$에 대해 $\mathcal{F}_v$에서
$s_{i, v} = (V, \sigma)$가 성립한다. $s_{i, v} = s_v$이므로 $s$에
대한 $(*)$를 얻는다.
\end{proof}

\begin{lemma}
\label{sheaves-lemma-stalk-sheafification}
$X$를 위상공간이라 하자. $\mathcal{F}$를 $X$ 위의 집합의 준층이라
하고 $x \in X$라 하자. 그러면 $\mathcal{F}_x = \mathcal{F}^\#_x$이다.
\end{lemma}

\begin{proof}
사상 $\mathcal{F}_x \to \mathcal{F}^\#_x$는 단사이다. 실제로 이미
사상 $\mathcal{F}_x \to \Pi(\mathcal{F})_x$가 단사이기 때문이다.
즉, 사상 $\mathcal{F}_x \to \Pi(\mathcal{F})_x$의 왼쪽 역인 표준적인
사상 $\Pi(\mathcal{F})_x \to \mathcal{F}_x$가 있다. 예
\ref{sheaves-example-sheaf-product-pointwise-stalk}를 보라.
전사임을 보이기 위해 $\overline{s} \in \mathcal{F}^\#_x$라 하자.
$\overline{s}$가 $s \in \mathcal{F}^\#(U)$인 $(U, s)$의 동치류가
되도록 하는 $x$의 열린 근방 $U$를 찾을 수 있다. 정의에 의해 이는
$x$의 열린 근방 $V \subset U$와 단면 $\sigma \in \mathcal{F}(V)$가
존재하여, $s|_V$가 $\Pi(\mathcal{F})(V)$에서 $\sigma$의 상이 됨을
뜻한다. 분명히 $(V, \sigma)$의 류는 $\overline{s}$로 가는
$\mathcal{F}_x$의 원소를 정의한다.
\end{proof}

\begin{lemma}
\label{sheaves-lemma-sheafify-universal}
$\mathcal{F}$를 $X$ 위의 집합의 준층이라 하자. 집합의 층으로 가는
임의의 사상 $\mathcal{F} \to \mathcal{G}$는 유일하게
$\mathcal{F} \to \mathcal{F}^\# \to \mathcal{G}$로 분해된다.
\end{lemma}

\begin{proof}
분명히 가환도식
$$
\xymatrix{
\mathcal{F} \ar[r] \ar[d] &
\mathcal{F}^\# \ar[r] \ar[d] &
\Pi(\mathcal{F}) \ar[d] \\
\mathcal{G} \ar[r] &
\mathcal{G}^\# \ar[r] &
\Pi(\mathcal{G}) \\
}
$$
이 있다. 따라서 $\mathcal{G} = \mathcal{G}^\#$임을 보이면 충분하다.
이를 위해서는 보조정리 \ref{sheaves-lemma-points-exactness}에 의해 모든 점
$x \in X$에 대하여 사상 $\mathcal{G}_x \to \mathcal{G}^\#_x$가
전단사임을 보이면 충분하다. 이는 위의 보조정리
\ref{sheaves-lemma-stalk-sheafification}이다.
\end{proof}

\noindent
이 보조정리가 실제로 말하는 바는 다음 함자들의 수반쌍이 있다는 것이다.
$i : \Sh(X) \to \textit{PSh}(X)$ (포함)와
$\# : \textit{PSh}(X) \to \Sh(X)$ (층화)이다. 공식은
$$
\Mor_{\textit{PSh}(X)}(\mathcal{F}, i(\mathcal{G}))
=
\Mor_{\Sh(X)}(\mathcal{F}^\#, \mathcal{G})
$$
이며, 이는 층화가 포함 함자의 왼쪽 수반임을 말한다.
% P01-ERRATA: P01-E1250 -- the frozen source says "a left adjoint of" rather than "left adjoint to".
Categories, 절 \ref{categories-section-adjoint}를 보라.

\begin{example}
\label{sheaves-example-sheafify-constant}
표기에 대해서는 예 \ref{sheaves-example-stalk-constant-presheaf}를 보라.
사상 $A_p \to \underline{A}$는 사상 $A_p^\# \to \underline{A}$를
유도한다. 이것이 동형사상임은 쉽게 알 수 있다. 말로 하면, 값이 $A$인
상수 준층의 층화는 값이 $A$인 상수층이다.
\end{example}

\begin{lemma}
\label{sheaves-lemma-separated-presheaf-into-sheaf}
$X$를 위상공간이라 하자. 준층 $\mathcal{F}$가 분리되어 있다는 것
(정의 \ref{sheaves-definition-separated}을 보라)은 표준적인 사상
$\mathcal{F} \to \mathcal{F}^\#$가 단사인 것과 동치이다.
\end{lemma}

\begin{proof}
이는 이 절에서 한 $\mathcal{F}^\#$의 구성으로부터 분명하다.
\end{proof}

\begin{lemma}
\label{sheaves-lemma-sheafify-epi-mono}
$X$를 위상공간이라 하자. 집합의 준층들의 전사 사상(각각 단사 사상)을
층화하면 전사 사상(각각 단사 사상)이 된다.
\end{lemma}

\begin{proof}
생략한다.
\end{proof}


\section{아벨 준층의 층화}
\label{sheaves-section-sheafify-abelian-presheaves}

\noindent
다음의 다소 이상해 보이는 보조정리는 아마 필요하지 않을 것이지만,
대수 구조의 준층을 층화할 때 매우 편리하다.

\begin{lemma}
\label{sheaves-lemma-diagram-fibre-product}
$X$를 위상공간이라 하고, $\mathcal{F}$를 $X$ 위의 집합의 준층이라 하자.
$U \subset X$가 열린집합이라 하자. 표준적인 섬유곱 도식
$$
\xymatrix{
\mathcal{F}^\#(U) \ar[d] \ar[r] &
\Pi(\mathcal{F})(U) \ar[d] \\
\prod_{x \in U} \mathcal{F}_x
\ar[r] &
\prod_{x \in U} \Pi(\mathcal{F})_x
}
$$
이 있다. 여기서 사상들은 다음과 같다.
\begin{enumerate}
\item 왼쪽 수직 사상의 성분들은
$\mathcal{F}^\#(U) \to \mathcal{F}^\#_x = \mathcal{F}_x$이다.
여기서 등식은 보조정리 \ref{sheaves-lemma-stalk-sheafification}이다.
\item 위쪽 수평 사상은 절 \ref{sheaves-section-sheafification}에서 설명한
준층의 사상 $\mathcal{F} \to \Pi(\mathcal{F})$로부터 온다.
\item 오른쪽 수직 사상은 자명한 성분 사상들
$\Pi(\mathcal{F})(U) \to \Pi(\mathcal{F})_x$를 갖는다.
\item 아래쪽 수평 사상은 성분들
$\mathcal{F}_x \to \Pi(\mathcal{F})_x$를 갖는데, 이들은 절
\ref{sheaves-section-sheafification}에서 설명한 준층의 사상
$\mathcal{F} \to \Pi(\mathcal{F})$로부터 온다.
\end{enumerate}
\end{lemma}

\begin{proof}
도식이 가환함은 분명하다. 이것이 섬유곱 도식임을 보여야 한다.
모든 사상 $\mathcal{F}_x \to \Pi(\mathcal{F})_x$가 단사이므로
(보조정리 \ref{sheaves-lemma-stalk-sheafification}의 증명 첫 부분을 보라),
아래쪽 수평 화살표는 단사이다. 단면
$s \in \Pi(\mathcal{F})(U)$가 $\mathcal{F}^\#$에 속하는 것은
$(*)$가 성립하는 것과 동치이다. 그런데 $(*)$는 모든 점의 근방에서
단면 $s$가 $\mathcal{F}$의 단면으로부터 온다는 뜻이다. 줄기 함자들의
정의에 따라, 이는 모든 줄기 $\Pi(\mathcal{F})_x$에서 $s$의 값이 줄기
$\mathcal{F}_x$의 원소로부터 온다는 것과 동치이다. 이로써 증명된다.
\end{proof}


\begin{lemma}
\label{sheaves-lemma-sheafify-abelian-presheaf}
$X$를 위상공간이라 하고, $\mathcal{F}$를 $X$ 위의 아벨 준층이라 하자.
$\mathcal{F} \to \mathcal{F}^\#$가 아벨 준층의 사상이 되도록 하는
$\mathcal{F}^\#$ 위의 아벨층 구조가 유일하게 존재한다. 더 나아가
다음 수반 성질이 성립한다.
$$
\Mor_{\textit{PAb}(X)}(\mathcal{F}, i(\mathcal{G}))
=
\Mor_{\textit{Ab}(X)}(\mathcal{F}^\#, \mathcal{G}).
$$
\end{lemma}

\begin{proof}
절 \ref{sheaves-section-sheafification}에서 정의한 집합의 층
$\Pi(\mathcal{F})$를 상기하자. 모든 줄기 $\mathcal{F}_x$는 아벨군이다.
보조정리 \ref{sheaves-lemma-stalk-abelian-presheaf}를 보라. 따라서 예
\ref{sheaves-example-sheaf-product-pointwise-algebraic-structure}에 의해
$\Pi(\mathcal{F})$는 아벨군의 층이다. 또한 사상
$\mathcal{F} \to \Pi(\mathcal{F})$가 아벨 준층의 사상임은 분명하다.
모든 열린 부분집합 $U \subset X$에 대하여 절
\ref{sheaves-section-sheafification}의 조건 $(*)$가
$\Pi(\mathcal{F})(U)$의 부분군을 정의함을 보이면, $\mathcal{F}^\#$는
표준적으로 아벨층의 구조를 물려받는다. 이를 직접 보이는 것은 매우
쉬우며, 간단하고 좋은 논증을 찾는 일은 독자에게 맡긴다. 여기서 사용할
논증은 대수 구조의 준층으로 일반화되며 다음과 같다. 보조정리
\ref{sheaves-lemma-diagram-fibre-product}는 $\mathcal{F}^\#(U)$가 아벨군들의
도식의 섬유곱임을 보여 준다. 따라서 원하는 대로 $\mathcal{F}^\#$는
아벨 부분군이다.
% P01-ERRATA: P01-E0041 -- the frozen source has subject-verb disagreement: "Lemma ... show".

\medskip\noindent
이 시점에서 보조정리 \ref{sheaves-lemma-stalk-abelian-presheaf}에 의해
$\mathcal{F}^\#_x$는 아벨군이고, $\mathcal{F}_x \to \mathcal{F}^\#_x$는
전단사(보조정리 \ref{sheaves-lemma-stalk-sheafification})이자 아벨군의
준동형임을 주목하자. 따라서 $\mathcal{F}_x \to \mathcal{F}^\#_x$는
아벨군의 동형사상이다. 아래에서는 이를 더 언급하지 않고 사용할 것이다.

\medskip\noindent
수반 성질을 증명하기 위해 집합의 준층을 층화하는 데 대한 수반 성질을
사용한다. 예를 들어 $\psi : \mathcal{F} \to i(\mathcal{G})$가 준층의
사상이면 층의 사상 $\psi' : \mathcal{F}^\# \to \mathcal{G}$를 얻는다.
% P01-ERRATA: P01-E1251 -- the frozen source omits the article before "morphism".
이것이 아벨층의 사상임을 확인해야 한다. 예컨대 보조정리
\ref{sheaves-lemma-stalk-sheafification}에 의해 줄기에서 이것이 참임을 주목한 뒤,
위의 보조정리 \ref{sheaves-lemma-check-homomorphism-stalks}를 사용하면 된다.
\end{proof}



\section{대수 구조의 준층의 층화}
\label{sheaves-section-sheafification-presheaves-structures}


\begin{lemma}
\label{sheaves-lemma-sheafify-presheaf-structures}
$X$를 위상공간이라 하자. $(\mathcal{C}, F)$를 대수 구조의 유형이라
하고, $\mathcal{F}$를 $X$ 위에서 $\mathcal{C}$에 값을 갖는 준층이라
하자. 다음 성질들을 갖는 $\mathcal{C}$에 값을 갖는 층
$\mathcal{F}^\#$와 $\mathcal{C}$에 값을 갖는 준층의 사상
$\mathcal{F} \to \mathcal{F}^\#$가 존재한다.
\begin{enumerate}
\item 사상 $\mathcal{F} \to \mathcal{F}^\#$는
$\mathcal{F}^\#$의 바탕 집합의 층을 $\mathcal{F}$의 바탕 집합의
준층을 층화한 것과 동일시한다.
\item $\mathcal{G}$가 $\mathcal{C}$에 값을 갖는 층일 때, 임의의 사상
$\mathcal{F} \to \mathcal{G}$에 대하여 유일한 분해
$\mathcal{F} \to \mathcal{F}^\# \to \mathcal{G}$가 존재한다.
\end{enumerate}
\end{lemma}

\begin{proof}
증명은 보조정리 \ref{sheaves-lemma-sheafify-abelian-presheaf}의 증명과 같으며,
보조정리 \ref{sheaves-lemma-image-contained-in}를 반복해서 적용한다
(예 \ref{sheaves-example-application-lemma-image-contained-in}도 보라).
그러나 주된 발상은 $\mathcal{F}^\#(U)$를 $\mathcal{C}$에서의 다음
도식의 섬유곱으로 정의하는 것이다.
% P01-ERRATA: P01-E1252 -- the frozen source lacks the comma before parenthetical "however".
$$
\xymatrix{
 &
\Pi(\mathcal{F})(U) \ar[d] \\
\prod_{x \in U} \mathcal{F}_x
\ar[r] &
\prod_{x \in U} \Pi(\mathcal{F})_x
}
$$
보조정리 \ref{sheaves-lemma-diagram-fibre-product}와 비교하라.
\end{proof}

\section{가군의 준층의 층화}
\label{sheaves-section-sheafification-presheaves-modules}

\begin{lemma}
\label{sheaves-lemma-sheafification-presheaf-modules}
$X$를 위상공간이라 하자. $\mathcal{O}$를 $X$ 위의 환의 준층이라 하고,
$\mathcal{F}$를 $\mathcal{O}$-가군의 준층이라 하자.
$\mathcal{O}^\#$를 $\mathcal{O}$의 층화라 하고, $\mathcal{F}^\#$를
아벨군의 준층으로서 $\mathcal{F}$를 층화한 것이라 하자. 도식
$$
\mathcal{O}^\# \times \mathcal{F}^\#
\longrightarrow
\mathcal{F}^\#
$$
을 가환으로 만들고 $\mathcal{F}^\#$를 $\mathcal{O}^\#$-가군의 층으로
만드는 집합의 층의 사상이 존재한다.
$$
\xymatrix{
\mathcal{O} \times \mathcal{F} \ar[r] \ar[d] &
\mathcal{F} \ar[d] \\
\mathcal{O}^\# \times \mathcal{F}^\# \ar[r] &
\mathcal{F}^\#
}
$$
또한 $\mathcal{G}$가 $\mathcal{O}^\#$-가군의 층이면,
$\mathcal{O}$-가군의 준층의 임의의 사상
$\mathcal{F} \to \mathcal{G}$는 ($\mathcal{G}$를 $\mathcal{O}$-가군으로
제한한 것에 이르는 사상으로 보아) 유일하게
$\mathcal{F} \to \mathcal{F}^\# \to \mathcal{G}$로 분해된다.
% P01-ERRATA: P01-E1253 -- the frozen source's parenthetical description of restriction of scalars is malformed.
여기서 $\mathcal{F}^\# \to \mathcal{G}$는
$\mathcal{O}^\#$-가군의 사상이다.
\end{lemma}

\begin{proof}
생략한다.
\end{proof}

\noindent
이는 실제로 함자
$i : \textit{Mod}(\mathcal{O}^\#) \to \textit{PMod}(\mathcal{O})$
(스칼라 제한과 층을 준층에 포함하는 것을 합친 것)와 보조정리의 층화 함자
${}^\# : \textit{PMod}(\mathcal{O}) \to \textit{Mod}(\mathcal{O}^\#)$가
수반임을 뜻한다. 공식으로는
$$
\Mor_{\textit{PMod}(\mathcal{O})}(\mathcal{F}, i\mathcal{G})
=
\Mor_{\textit{Mod}(\mathcal{O}^\#)}(\mathcal{F}^\#, \mathcal{G})
$$
이다.

\medskip\noindent
$X$를 위상공간이라 하자. $\mathcal{O}_1 \to \mathcal{O}_2$를 $X$ 위의
환의 층들의 사상이라 하자. 절 \ref{sheaves-section-presheaves-modules}에서
이 상황에 딸린 가군의 준층들에 대한 제한 함자와 환 변경 함자를 정의했다.

\medskip\noindent
$\mathcal{F}$가 $\mathcal{O}_2$-가군의 층이면 $\mathcal{F}$의 제한
$\mathcal{F}_{\mathcal{O}_1}$은 분명히 $\mathcal{O}_1$-가군의 층이다.
따라서 제한 함자
$$
\textit{Mod}(\mathcal{O}_2)
\longrightarrow
\textit{Mod}(\mathcal{O}_1)
$$
를 얻는다.

\medskip\noindent
한편 $\mathcal{O}_1$-가군의 층 $\mathcal{G}$가 주어졌을 때,
$\mathcal{O}_2$-가군의 준층
$\mathcal{O}_2 \otimes_{p, \mathcal{O}_1} \mathcal{G}$는 일반적으로
층이 아니다. 따라서 준층에 대한 우리의 구성을 층화하여
{\it 텐서곱 층} $\mathcal{O}_2 \otimes_{\mathcal{O}_1} \mathcal{G}$를
공식
$$
\mathcal{O}_2 \otimes_{\mathcal{O}_1} \mathcal{G}
=
(\mathcal{O}_2 \otimes_{p, \mathcal{O}_1} \mathcal{G})^\#
$$
으로 정의한다. 환 변경 함자
$$
\textit{Mod}(\mathcal{O}_1)
\longrightarrow
\textit{Mod}(\mathcal{O}_2)
$$
를 얻는다.

\begin{lemma}
\label{sheaves-lemma-adjointness-tensor-restrict}
위와 같은 $X$, $\mathcal{O}_1$, $\mathcal{O}_2$, $\mathcal{F}$,
$\mathcal{G}$에 대하여 표준적인 전단사
$$
\Hom_{\mathcal{O}_1}(\mathcal{G}, \mathcal{F}_{\mathcal{O}_1})
=
\Hom_{\mathcal{O}_2}(
\mathcal{O}_2 \otimes_{\mathcal{O}_1} \mathcal{G},
\mathcal{F}
)
$$
가 존재한다. 달리 말하면 제한 함자와 환 변경 함자는 서로 수반이다.
\end{lemma}

\begin{proof}
이는 보조정리 \ref{sheaves-lemma-adjointness-tensor-restrict-presheaves}와,
$\mathcal{F}$가 층이므로 성립하는 등식
$\Hom_{\mathcal{O}_2}(
\mathcal{O}_2 \otimes_{\mathcal{O}_1} \mathcal{G},
\mathcal{F}
)
=
\Hom_{\mathcal{O}_2}(
\mathcal{O}_2 \otimes_{p, \mathcal{O}_1} \mathcal{G},
\mathcal{F}
)$
로부터 따른다.
\end{proof}

\begin{lemma}
\label{sheaves-lemma-stalk-tensor-sheaf-modules}
$X$를 위상공간이라 하자. $\mathcal{O} \to \mathcal{O}'$를 $X$ 위의
환의 층들의 사상이라 하고, $\mathcal{F}$를 $\mathcal{O}$-가군의 층이라
하자. $x \in X$라 하자. $\mathcal{O}'_x$-가군으로서
$$
\mathcal{F}_x \otimes_{\mathcal{O}_x} \mathcal{O}'_x
=
(\mathcal{F} \otimes_\mathcal{O} \mathcal{O}')_x
$$
이다.
% P01-ERRATA: P01-E0042 -- the frozen source omits "of" in "a sheaf O-modules".
\end{lemma}

\begin{proof}
이는 보조정리 \ref{sheaves-lemma-stalk-tensor-presheaf-modules}와 줄기를 취하는 것이
층화와 가환한다는 사실로부터 바로 따른다.
\end{proof}


\section{연속사상과 층}
\label{sheaves-section-presheaves-functorial}

\noindent
$f : X \to Y$를 위상공간들의 연속사상이라 하자. 준층과 층에 대한
직접상 함자와 역상 함자를 정의할 것이다.

\medskip\noindent
$\mathcal{F}$를 $X$ 위의 집합의 준층이라 하자. 임의의 열린집합
$V \subset Y$에 대하여
$$
f_*\mathcal{F}(V) = \mathcal{F}(f^{-1}(V))
$$
로 놓아 $\mathcal{F}$의 {\it 직접상}을 정의한다.
열린집합 $V_1 \subset V_2 \subset Y$가 주어졌을 때 제한 사상은
도식
$$
\xymatrix{
f_*\mathcal{F}(V_2) \ar[d] \ar@{=}[r] &
\mathcal{F}(f^{-1}(V_2)) \ar[d]^{\text{restriction for }\mathcal{F}} \\
f_*\mathcal{F}(V_1) \ar@{=}[r] &
\mathcal{F}(f^{-1}(V_1))
}
$$
의 가환성으로 주어진다. 이것이 집합의 준층을 정의함은 분명하다.
이 구성은 준층 $\mathcal{F}$에 대해 분명히 함자적이므로 함자
$$
f_* : \textit{PSh}(X) \longrightarrow \textit{PSh}(Y).
$$
를 얻는다.

\begin{lemma}
\label{sheaves-lemma-pushforward-sheaf}
$f : X \to Y$를 연속사상이라 하고, $\mathcal{F}$를 $X$ 위의 집합의
층이라 하자. 그러면 $f_*\mathcal{F}$는 $Y$ 위의 층이다.
\end{lemma}

\begin{proof}
$V = \bigcup V_j$가 $Y$의 열린 덮개이면
$f^{-1}(V) = \bigcup f^{-1}(V_j)$가 $X$의 열린 덮개라는 사실로부터
즉시 따른다.
\end{proof}

\noindent
그 결과 함자
$$
f_* : \Sh(X) \longrightarrow \Sh(Y).
$$
를 얻는다. 이는 다음의 강한 의미에서 합성과 양립한다.

\begin{lemma}
\label{sheaves-lemma-pushforward-composition}
$f : X \to Y$와 $g : Y \to Z$를 위상공간들의 연속사상이라 하자.
함자 $(g \circ f)_*$와 $g_* \circ f_*$는 같다(집합의 준층에 대해서도
층에 대해서도).
\end{lemma}

\begin{proof}
이는 $(g \circ f)_*\mathcal{F}(W) =
\mathcal{F}((g \circ f)^{-1}W)$이고,
$(g_* \circ f_*)\mathcal{F}(W) = \mathcal{F}(f^{-1} g^{-1} W)$이며,
$(g \circ f)^{-1}W = f^{-1} g^{-1} W$이기 때문이다.
\end{proof}

\noindent
$\mathcal{G}$를 $Y$ 위의 집합의 준층이라 하자. 주어진 준층
$\mathcal{G}$의 {\it 역상 준층} $f_p\mathcal{G}$는 준층에 대한
직접상 $f_*$의 왼쪽 수반으로 정의한다. 달리 말하면 이는 $X$ 위의 준층
$f_p \mathcal{G}$로서
$$
\Mor_{\textit{PSh}(X)}(f_p\mathcal{G}, \mathcal{F})
=
\Mor_{\textit{PSh}(Y)}(\mathcal{G}, f_*\mathcal{F}).
$$
을 만족해야 한다. Yoneda 보조정리에 의해 이는 역상을 유일하게 결정한다.
실제로 이러한 역상이 존재한다.

\begin{lemma}
\label{sheaves-lemma-pullback-presheaves}
$f : X \to Y$를 연속사상이라 하자. $f_*$의 왼쪽 수반인 함자
$f_p : \textit{PSh}(Y) \to \textit{PSh}(X)$가 존재한다. 준층
$\mathcal{G}$에 대하여 이 함자는 규칙
$$
f_p\mathcal{G}(U) = \colim_{f(U) \subset V} \mathcal{G}(V)
$$
으로 결정된다. 여기서 여극한은 $Y$ 안에서 $f(U)$의 열린근방들
$V$의 모임에 대해 취한다. 이 여극한들은 유향 부분순서집합에 대해
취해진다. ($f_p\mathcal{G}$의 제한 사상은 증명에서 설명한다.)
\end{lemma}

\begin{proof}
여극한은 $f(U)$를 포함하는 열린 부분집합 $V \subset Y$들로 이루어진
부분순서집합에 대해 취하며, 순서는 포함관계의 역순이다. $V, V'$가
그 안에 있으면 $V \cap V'$도 그 안에 있으므로 이는 유향
부분순서집합이다. 더 나아가 $U_1 \subset U_2$이면 $f(U_2)$의 모든
열린근방은 $f(U_1)$의 열린근방이다. 따라서 $f_p\mathcal{G}(U_2)$를
정의하는 계는 $f_p\mathcal{G}(U_1)$을 정의하는 계의 부분계이며,
제한 사상을 얻는다(예컨대 Categories, 보조정리
\ref{categories-lemma-functorial-colimit}의 일반론을 적용한다).

\medskip\noindent
여극한의 구성은 $\mathcal{G}$에 대해 분명히 함자적이고 제한 사상들도
마찬가지임을 주목하자. 따라서 $f_p$를 함자로 정의했다.

\medskip\noindent
작지만 유용한 관찰로, 표준적인 사상
$\mathcal{G}(U) \to f_p\mathcal{G}(f^{-1}(U))$가 존재한다.
이는 $f(f^{-1}(U))$의 열린근방들의 계가 원소 $U$를 포함하기 때문이다.
이 사상은 제한 사상들과 양립한다. 달리 말하면 표준적인 사상
$i_\mathcal{G} : \mathcal{G} \to f_* f_p \mathcal{G}$가 있다.

\medskip\noindent
$\mathcal{F}$를 $X$ 위의 집합의 준층이라 하자. 집합의 준층의 사상
$\psi : f_p\mathcal{G} \to \mathcal{F}$가 주어졌다고 하자.
이에 대응하는 사상 $\mathcal{G} \to f_*\mathcal{F}$는
$f_*\psi \circ i_\mathcal{G} :
\mathcal{G} \to f_* f_p \mathcal{G} \to f_* \mathcal{F}$이다.

\medskip\noindent
또 하나의 작지만 유용한 관찰로, 표준적인 사상
$c_\mathcal{F} : f_p f_* \mathcal{F} \to \mathcal{F}$가 존재한다.
실제로 $U \subset X$가 열려 있다고 하자. $Y$에서 그 상을 포함하는
모든 열린근방 $V \supset f(U)$에 대해 사상
$f_*\mathcal{F}(V) = \mathcal{F}(f^{-1}(V))\to \mathcal{F}(U)$가
존재하는데, 바로 $\mathcal{F}$의 제한 사상이다. 이는 $f(U)$를 포함하는
여러 열린집합의 $f^{-1}$ 위에서의 $\mathcal{F}$의 값들 사이의 제한
사상들과 양립한다. 따라서 표준적인 사상
$f_p f_* \mathcal{F}(U) \to \mathcal{F}(U)$를 얻는다. 또 하나의
자명한 확인으로 이 사상들이 제한 사상들과 양립하여 집합의 준층의 사상
$c_\mathcal{F}$를 정의함을 알 수 있다.

\medskip\noindent
$\varphi : \mathcal{G} \to f_*\mathcal{F}$가 집합의 준층의 사상이라고
하자. $f_p\varphi : f_p \mathcal{G} \to f_p f_* \mathcal{F}$를
생각하자. 뒤에 $c_\mathcal{F}$를 합성하면 원하는 사상
$c_\mathcal{F} \circ f_p\varphi : f_p\mathcal{G} \to \mathcal{F}$를
얻는다. 이 구성이 위에서 준 반대 방향의 구성의 역이라는 확인은 생략한다.
\end{proof}

\begin{lemma}
\label{sheaves-lemma-stalk-pullback-presheaf}
$f : X \to Y$를 연속사상이라 하자. $x \in X$라 하고,
$\mathcal{G}$를 $Y$ 위의 집합의 준층이라 하자. 줄기들의 표준적인 전단사
$(f_p\mathcal{G})_x = \mathcal{G}_{f(x)}$가 있다.
\end{lemma}

\begin{proof}
이는 다음과 같이 알 수 있다.
\begin{eqnarray*}
(f_p\mathcal{G})_x
& = &
\colim_{x \in U} f_p\mathcal{G}(U) \\
& = &
\colim_{x \in U} \colim_{f(U) \subset V} \mathcal{G}(V) \\
& = &
\colim_{f(x) \in V} \mathcal{G}(V) \\
& = &
\mathcal{G}_{f(x)}
\end{eqnarray*}
여기서 Categories, 보조정리 \ref{categories-lemma-colimits-commute}와,
$f(x)$를 포함하는 $Y$의 임의의 열린집합 $V$가 위의 세 번째 표현에
나타난다는 사실을 사용했다. 자세한 내용은 생략한다.
\end{proof}

\noindent
$\mathcal{G}$를 $Y$ 위의 집합의 층이라 하자. {\it 역상층}
$f^{-1}\mathcal{G}$를 공식
$$
f^{-1}\mathcal{G} = (f_p\mathcal{G})^\# .
$$
으로 정의한다. 역상 $f^{-1}$은 층에 대한 직접상의 왼쪽 수반이다.
% P01-ERRATA: P01-E1254 -- the frozen source says "a left adjoint of" rather than "left adjoint to".
달리 말하면
$$
\Mor_{\Sh(X)}(f^{-1}\mathcal{G}, \mathcal{F})
=
\Mor_{\Sh(Y)}(\mathcal{G}, f_*\mathcal{F}).
$$
이다. 실제로
\begin{eqnarray*}
\Mor_{\Sh(X)}(f^{-1}\mathcal{G}, \mathcal{F})
& = &
\Mor_{\textit{PSh}(X)}(f_p\mathcal{G}, \mathcal{F}) \\
& = &
\Mor_{\textit{PSh}(Y)}(\mathcal{G}, f_*\mathcal{F}) \\
& = &
\Mor_{\Sh(Y)}(\mathcal{G}, f_*\mathcal{F})
\end{eqnarray*}
이다. 첫 번째 등식에는 층화가 층을 준층에 포함하는 함자의 왼쪽
수반이라는 사실을 사용한다. 두 번째 등식에는 $f_p$가 준층에서
$f_*$의 왼쪽 수반이라는 사실을 사용한다. 보조정리
\ref{sheaves-lemma-f-map}의 증명에서 이 명제로 돌아올 것이다.

\begin{lemma}
\label{sheaves-lemma-stalk-pullback}
$x \in X$라 하고 $\mathcal{G}$를 $Y$ 위의 집합의 층이라 하자.
줄기들의 표준적인 전단사
$(f^{-1}\mathcal{G})_x = \mathcal{G}_{f(x)}$가 있다.
\end{lemma}

\begin{proof}
이는 보조정리 \ref{sheaves-lemma-stalk-sheafification}와
\ref{sheaves-lemma-stalk-pullback-presheaf}를 결합한 것이다.
\end{proof}

\begin{lemma}
\label{sheaves-lemma-pullback-composition}
$f : X \to Y$와 $g : Y \to Z$를 위상공간들의 연속사상이라 하자.
함자 $(g \circ f)^{-1}$과 $f^{-1} \circ g^{-1}$은 표준적으로
동형이다. 마찬가지로 준층에서
$(g \circ f)_p \cong f_p \circ g_p$이다.
\end{lemma}

\begin{proof}
수반함자가 유일한 동형사상을 제외하고 유일하다는 사실과 보조정리
\ref{sheaves-lemma-pushforward-composition}를 사용하면 알 수 있다.
\end{proof}

\begin{definition}
\label{sheaves-definition-f-map}
$f : X \to Y$를 연속사상이라 하자. $\mathcal{F}$를 $X$ 위의 집합의
층이라 하고 $\mathcal{G}$를 $Y$ 위의 집합의 층이라 하자.
{\it $f$-사상 $\xi : \mathcal{G} \to \mathcal{F}$}이란 열린 부분집합
$V \subset Y$로 첨자화된 사상들의 모임
$\xi_V : \mathcal{G}(V) \to \mathcal{F}(f^{-1}(V))$로서, 모든 열린집합
$V' \subset V \subset Y$에 대하여 도식
$$
\xymatrix{
\mathcal{G}(V) \ar[r]_{\xi_V} \ar[d]_{\text{restriction of }\mathcal{G}} &
\mathcal{F}(f^{-1}V) \ar[d]^{\text{restriction of }\mathcal{F}} \\
\mathcal{G}(V') \ar[r]^{\xi_{V'}} &
\mathcal{F}(f^{-1}V')
}
$$
이 가환하도록 하는 것이다.
\end{definition}

\noindent
문헌에서는 때때로 이를 보조정리 \ref{sheaves-lemma-f-map}의 (4)처럼 다르게
정의하지만, 그 보조정리가 보여 주듯 실제 차이는 없다.

\begin{lemma}
\label{sheaves-lemma-f-map}
$f : X \to Y$를 연속사상이라 하자. 다음 네 집합 사이에 전단사들이 있다.
\begin{enumerate}
\item 사상 $\mathcal{G} \to f_*\mathcal{F}$들의 집합,
\item 사상 $f^{-1}\mathcal{G} \to \mathcal{F}$들의 집합,
\item $f$-사상 $\xi : \mathcal{G} \to \mathcal{F}$들의 집합, 그리고
\item 모든 열린집합 $U \subset X$와 $V \subset Y$로서
$f(U) \subset V$인 것들에 대한 사상
$\xi_{U, V} : \mathcal{G}(V) \to \mathcal{F}(U)$들의 모임 중에서
모든 제한 사상과 양립하는 것들의 집합.
\end{enumerate}
이 전단사들은 $\mathcal{F} \in \Sh(X)$와
$\mathcal{G} \in \Sh(Y)$에 대해 함자적이다.
\end{lemma}

\begin{proof}
층의 사상 $a : \mathcal{G} \to f_*\mathcal{F}$는 정의상 $Y$의 각
열린집합 $V$에 사상 $a_V : \mathcal{G}(V) \to f_*\mathcal{F}(V)$를
대응시키는 규칙이며,
$f_*\mathcal{F}(V) = \mathcal{F}(f^{-1}(V))$이다. 따라서 적어도
그 “데이터”는 $\mathcal{G}$에서 $\mathcal{F}$로 가는 $f$-사상
$\xi$에 필요한 것과 정확히 일치한다. 집합 (1)과 (3)이 전단사임을
보이기 위해, $a$가 층의 사상인 것은 이에 대응하는 사상족 $a_V$가
정의 \ref{sheaves-definition-f-map}의 조건을 만족하는 것과 동치임을 관찰한다.
% P01-ERRATA: P01-E1255 -- the frozen source omits "the" before "corresponding family".

\medskip\noindent
$f^{-1}\mathcal{G}$가 $f_p\mathcal{G}$의 층화임을 상기하자.
층화의 보편 성질에 의해 층의 사상
$b : f^{-1}\mathcal{G} \to \mathcal{F}$를 주는 것은 준층의 사상
$b_p : f_p\mathcal{G} \to \mathcal{F}$를 주는 것과 같다. 여기서
$f_p$는 이 절 앞부분에서 정의한 함자이다. 이러한 사상 $b_p$를 주려면
$X$의 각 열린집합 $U$에 대하여 제한 사상들과 양립하는 사상
$$
b_{p, U} :
\colim_{f(U) \subset V} \mathcal{G}(V)
\longrightarrow
\mathcal{F}(U)
$$
을 지정해야 한다. 우리는 $b_{p, U}$를 $Y$에서 $f(U) \subset V$인
모든 열린집합 $V$에 대한 사상
$b_{p, U, V} : \mathcal{G}(V) \to \mathcal{F}(U)$들의 모임으로
보아도 되고 그렇게 본다. 이 사상들은 생각할 수 있는 모든 제한 사상과
양립해야 한다. 달리 말하면 $b_p$는 (4)와 같은 사상들의 모임에
대응함을 알 수 있다. 물론 역으로 그러한 모임은 사상 $b_p$를 정의하고,
따라서 사상 $b : f^{-1}\mathcal{G} \to \mathcal{F}$를 정의한다.

\medskip\noindent
보조정리의 증명을 끝내려면, (4)와 같은 모임 $\xi_{U, V}$에
$\xi_V = \xi_{f^{-1}(V), V}$인 $f$-사상 $\xi$를 대응시키는
“구조를 잊는” 규칙이 전단사임을 보여야 한다. 이를 위해 $\xi$가 보통의
$f$-사상이면 $\tilde \xi_{U, V}$를
$\xi_V : \mathcal{G}(V) \to \mathcal{F}(f^{-1}(V))$에 제한 사상
$\mathcal{F}(f^{-1}(V)) \to \mathcal{F}(U)$를 합성한 것으로 정의한다.
% P01-ERRATA: P01-E0043 -- the frozen source misspells "restriction" as "restruction".
이는 정확히 $f(U) \subset V$, 즉 $U \subset f^{-1}(V)$이므로 말이 된다.
이로써 증명이 끝난다.
\end{proof}

\noindent
$X$ 위 또는 $Y$ 위의 층들 사이의 사상 대신 $f$-사상을 생각하는 것이
때때로 편리하다. $f$-사상의 합성을 다음과 같이 정의한다.

\begin{definition}
\label{sheaves-definition-composition-f-maps}
$f : X \to Y$와 $g : Y \to Z$를 위상공간들의 연속사상이라 하자.
$\mathcal{F}$를 $X$ 위의 층, $\mathcal{G}$를 $Y$ 위의 층,
$\mathcal{H}$를 $Z$ 위의 층이라 하자. $\varphi : \mathcal{G} \to
\mathcal{F}$를 $f$-사상이라 하고, $\psi : \mathcal{H} \to
\mathcal{G}$를 $g$-사상이라 하자.
% P01-ERRATA: P01-E0044 -- the frozen source uses "an g-map" here and again in lemma-compose-f-maps-stalks.
{\it $\varphi$와 $\psi$의 합성}은 다음 도식들의 가환성으로 정의되는
$(g \circ f)$-사상 $\varphi \circ \psi$이다.
$$
\xymatrix{
\mathcal{H}(W) \ar[rr]_{(\varphi \circ \psi)_W}
\ar[rd]_{\psi_W} & &
\mathcal{F}(f^{-1}g^{-1}W) \\
&
\mathcal{G}(g^{-1}W)
\ar[ru]_{\varphi_{g^{-1}W}}
}
$$
\end{definition}

\noindent
이 정의가 잘 작동한다는 확인은 독자에게 맡긴다. 달리 생각하면
$\varphi \circ \psi$를 합성
$$
\mathcal{H}
\xrightarrow{\psi}
g_*\mathcal{G}
\xrightarrow{g_*\varphi}
g_* f_* \mathcal{F} = (g \circ f)_* \mathcal{F}
$$
으로 볼 수 있다. 이제 $f$-사상을 생각하는 편이 왠지 더 쉬워 보이지
않는가?

\medskip\noindent
마지막으로 연속사상 $f : X \to Y$와 $f$-사상
$\varphi : \mathcal{G} \to \mathcal{F}$가 주어졌을 때, 모든
$x \in X$에 대하여 줄기 위의 자연스러운 사상
$$
\varphi_x : \mathcal{G}_{f(x)} \longrightarrow \mathcal{F}_x
$$
이 있다. $\mathcal{G}_{f(x)}$의 원소를 나타내는 대표 $(V, s)$의 상은
대표 $(f^{-1}V, \varphi_V(s))$를 갖는 $\mathcal{F}_x$의 원소이다.
이것이 잘 정의됨을 확인하는 일은 독자에게 맡긴다. 달리 말하면, 이는
모든 열린집합 $f(x) \in V \subset Y$에 대해 도식
$$
\xymatrix{
\mathcal{F}(f^{-1}V) \ar[r] &
\mathcal{F}_x \\
\mathcal{G}(V) \ar[r] \ar[u]^{\varphi_V} &
\mathcal{G}_{f(x)} \ar[u]^{\varphi_x}
}
$$
이 가환하도록 하는 유일한 사상이다.

\begin{lemma}
\label{sheaves-lemma-compose-f-maps-stalks}
$f : X \to Y$와 $g : Y \to Z$를 위상공간들의 연속사상이라 하자.
$\mathcal{F}$를 $X$ 위의 층, $\mathcal{G}$를 $Y$ 위의 층,
$\mathcal{H}$를 $Z$ 위의 층이라 하자. $\varphi : \mathcal{G} \to
\mathcal{F}$를 $f$-사상이라 하고, $\psi : \mathcal{H} \to
\mathcal{G}$를 $g$-사상이라 하자. $x \in X$를 점이라 하자.
줄기 위의 사상
$(\varphi \circ \psi)_x : \mathcal{H}_{g(f(x))}
\to \mathcal{F}_x$는 합성
$$
\mathcal{H}_{g(f(x))}
\xrightarrow{\psi_{f(x)}}
\mathcal{G}_{f(x)}
\xrightarrow{\varphi_x}
\mathcal{F}_x
$$
이다.
\end{lemma}

\begin{proof}
정의 \ref{sheaves-definition-composition-f-maps}와 위에서 정의한 줄기 위의
사상으로부터 즉시 따른다.
\end{proof}


\section{연속사상과 아벨층}
\label{sheaves-section-abelian-presheaves-functorial}

\noindent
$f : X \to Y$를 연속사상이라 하자. 절
\ref{sheaves-section-presheaves-functorial}의 대응물과 비슷한 성질을 갖는 함자들
\begin{eqnarray*}
f_* : \textit{PAb}(X) & \longrightarrow & \textit{PAb}(Y) \\
f_* : \textit{Ab}(X) & \longrightarrow & \textit{Ab}(Y) \\
f_p : \textit{PAb}(Y) & \longrightarrow & \textit{PAb}(X) \\
f^{-1} : \textit{Ab}(Y) & \longrightarrow & \textit{Ab}(X)
\end{eqnarray*}
이 있다고 주장한다. 이를 보기 위해 다음과 같이 논증한다.

\medskip\noindent
각 함자는 절 \ref{sheaves-section-presheaves-functorial}의 대응하는 함자와 같은
방식으로 구성할 것이다. 그 절의 모든 여극한이 유향 여극한이므로 이
방법이 작동한다(하지만 아래에서 그 과정을 살펴볼 것이다).

\medskip\noindent
먼저 $\mathcal{F}$를 $X$ 위의 아벨 준층, $\mathcal{G}$를 $Y$ 위의
아벨 준층이라 하고, 아벨군으로서
\begin{eqnarray*}
f_*\mathcal{F}(V) & = & \mathcal{F}(f^{-1}(V)) \\
f_p\mathcal{G}(U) & = & \colim_{f(U) \subset V} \mathcal{G}(V)
\end{eqnarray*}
로 정의한다. 제한 사상들은 집합의 준층에 대한 제한 사상들과 같고,
모두 아벨군 준동형이다.

\medskip\noindent
대응 $\mathcal{F} \mapsto f_*\mathcal{F}$와
$\mathcal{G} \to f_p\mathcal{G}$는 아벨군의 준층 범주 위의 함자이다.
% P01-ERRATA: P01-E0046 -- the frozen source uses \to rather than \mapsto for its second object assignment; P01-E1256 independently records the same locus.
이는 분명하다. 예를 들어 아벨 준층의 사상
$\mathcal{G}_1 \to \mathcal{G}_2$는 모든 사상이 준동형인 유향계의 사상
$\{\mathcal{G}_1(V)\}_{f(U) \subset V} \to
\{\mathcal{G}_2(V)\}_{f(U) \subset V}$를 낳고, 따라서 아벨군 준동형
$f_p\mathcal{G}_1(U) \to f_p\mathcal{G}_2(U)$를 낳는다.

\medskip\noindent
함자 $f_*$와 $f_p$는 아벨군의 준층 범주에서 수반이다. 즉,
$$
\Mor_{\textit{PAb}(X)}(f_p\mathcal{G}, \mathcal{F})
=
\Mor_{\textit{PAb}(Y)}(\mathcal{G}, f_*\mathcal{F}).
$$
이다. 이를 증명하기 위해, 보조정리 \ref{sheaves-lemma-pullback-presheaves}의
증명에 나오는 사상
$i_\mathcal{G} : \mathcal{G} \to f_* f_p\mathcal{G}$가 아벨 준층의
사상임을 주목하자. 따라서 $\psi : f_p\mathcal{G} \to \mathcal{F}$가
아벨 준층의 사상이면, 이에 대응하는 사상
$\mathcal{G} \to f_*\mathcal{F}$는 사상
$f_*\psi \circ i_\mathcal{G} :
\mathcal{G} \to f_* f_p \mathcal{G} \to f_* \mathcal{F}$이고,
이 역시 아벨 준층의 사상이다.
% P01-ERRATA: P01-E1257 -- the frozen English sentence has two competing finite predicates.
반대 방향에 대해서는 보조정리 \ref{sheaves-lemma-pullback-presheaves}의 증명에
나오는 사상 $c_\mathcal{F} : f_p f_* \mathcal{F} \to \mathcal{F}$도
아벨 준층의 사상임을 지적한다(이는 모두 준동형인 $\mathcal{F}$의 제한
사상들로 만들어지기 때문이다). 따라서 아벨 준층의 사상
$\varphi : \mathcal{G} \to f_*\mathcal{F}$가 주어지면 사상
$c_\mathcal{F} \circ f_p\varphi : f_p\mathcal{G} \to \mathcal{F}$도
아벨 준층의 사상이다. 구성들 $\psi \mapsto f_*\psi$와
$\varphi \mapsto c_\mathcal{F} \circ f_p\varphi$는 집합의 준층의
사상에 대한 구성으로서 서로 역이므로, 아벨 준층의 사상에 대해서도
서로 역임을 알 수 있다.
% P01-ERRATA: P01-E1258 -- the frozen source's first stated construction omits precomposition with the unit.

\medskip\noindent
$\mathcal{F}$가 $Y$ 위의 아벨층이면 $f_*\mathcal{F}$는 $X$ 위의
아벨층이다.
% P01-ERRATA: P01-E0045 -- the frozen source reverses X and Y here despite f:X->Y; the canonical Korean text retains that source statement.
이는 아벨층의 정의와 집합의 층에 대해 이것이 참이라는 사실로부터
따른다. 보조정리 \ref{sheaves-lemma-pushforward-sheaf}를 보라. 이로써 아벨층의
범주 위의 함자 $f_*$를 정의한다.

\medskip\noindent
앞에서와 같이 $f^{-1}\mathcal{G} = (f_p\mathcal{G})^\#$로 정의한다.
$f_*$와 $f^{-1}$의 수반성은 집합의 준층의 경우처럼 형식적으로 따른다.
그 논증은 다음과 같다.
\begin{eqnarray*}
\Mor_{\textit{Ab}(X)}(f^{-1}\mathcal{G}, \mathcal{F})
& = &
\Mor_{\textit{PAb}(X)}(f_p\mathcal{G}, \mathcal{F}) \\
& = &
\Mor_{\textit{PAb}(Y)}(\mathcal{G}, f_*\mathcal{F}) \\
& = &
\Mor_{\textit{Ab}(Y)}(\mathcal{G}, f_*\mathcal{F})
\end{eqnarray*}

\begin{lemma}
\label{sheaves-lemma-pullback-abelian-stalk}
$f : X \to Y$를 연속사상이라 하자.
\begin{enumerate}
\item $\mathcal{G}$를 $Y$ 위의 아벨 준층이라 하고 $x \in X$라 하자.
보조정리 \ref{sheaves-lemma-stalk-pullback-presheaf}의 전단사
$\mathcal{G}_{f(x)} \to (f_p\mathcal{G})_x$는 아벨군의 동형사상이다.
\item $\mathcal{G}$를 $Y$ 위의 아벨층이라 하고 $x \in X$라 하자.
보조정리 \ref{sheaves-lemma-stalk-pullback}의 전단사
$\mathcal{G}_{f(x)} \to (f^{-1}\mathcal{G})_x$는 아벨군의 동형사상이다.
\end{enumerate}
\end{lemma}

\begin{proof}
생략한다.
\end{proof}

\noindent
연속사상 $f : X \to Y$와 $X$ 위의 아벨군의 층 $\mathcal{F}$,
$Y$ 위의 아벨군의 층 $\mathcal{G}$가 주어지면, 아벨군의 층의
{\it $f$-사상 $\mathcal{G} \to \mathcal{F}$}이라는 개념이 의미를
갖는다. 정의 \ref{sheaves-definition-f-map}에서 집합의 사상을 아벨군의
준동형으로 바꾸어 정확히 같은 방식으로 정의할 수도 있고, 단순히 아벨층의
사상 $\mathcal{G} \to f_*\mathcal{F}$와 같은 것이라고 말할 수도 있다.
아래에서는 이 개념을 자유롭게 사용할 것이다. $\mathcal{G}$와
$\mathcal{F}$ 사이의 $f$-사상들의 군은 군들
$\Mor_{\textit{Ab}(X)}(f^{-1}\mathcal{G}, \mathcal{F})$ 및
$\Mor_{\textit{Ab}(Y)}(\mathcal{G}, f_*\mathcal{F})$와 표준적으로
전단사이다.

\medskip\noindent
$f$-사상의 합성은 집합의 층의 $f$-사상의 경우와 정확히 같은 방식으로
정의한다. 더 나아가 위와 같은 $f$-사상
$\mathcal{G} \to \mathcal{F}$가 주어졌을 때 유도되는 줄기 위의 사상들
$$
\varphi_x : \mathcal{G}_{f(x)} \longrightarrow \mathcal{F}_x
$$
은 아벨군 준동형이다.


\section{연속사상과 대수 구조의 층}
\label{sheaves-section-presheaves-structures-functorial}

\noindent
$(\mathcal{C}, F)$를 대수 구조의 유형이라 하자. 위상공간 $X$에 대하여
다음 표기를 도입한다.
\begin{enumerate}
\item $\textit{PSh}(X, \mathcal{C})$는 $\mathcal{C}$에 값을 갖는
준층들의 범주이다.
\item $\Sh(X, \mathcal{C})$는 $\mathcal{C}$에 값을 갖는 층들의
범주이다.
\end{enumerate}
$f : X \to Y$를 위상공간들의 연속사상이라 하자. 앞 절과 같은 논증으로,
아벨 (준)층에 대해 구성한 함자들과 같은 방식으로 구성되고 같은 성질을
갖는 함자들
\begin{eqnarray*}
f_* : \textit{PSh}(X, \mathcal{C}) &
\longrightarrow &
\textit{PSh}(Y, \mathcal{C}) \\
f_* : \Sh(X, \mathcal{C}) &
\longrightarrow &
\Sh(Y, \mathcal{C}) \\
f_p : \textit{PSh}(Y, \mathcal{C}) &
\longrightarrow &
\textit{PSh}(X, \mathcal{C}) \\
f^{-1} : \Sh(Y, \mathcal{C}) &
\longrightarrow &
\Sh(X, \mathcal{C})
\end{eqnarray*}
이 있음을 알 수 있다. 특히 가환도식
$$
\xymatrix{
\textit{PSh}(X, \mathcal{C}) \ar[r]^{f_*} \ar[d]^F &
\textit{PSh}(Y, \mathcal{C}) \ar[d]^F &
\Sh(X, \mathcal{C}) \ar[r]^{f_*} \ar[d]^F &
\Sh(Y, \mathcal{C}) \ar[d]^F
\\
\textit{PSh}(X) \ar[r]^{f_*} &
\textit{PSh}(Y) &
\Sh(X) \ar[r]^{f_*} &
\Sh(Y)
\\
\textit{PSh}(Y, \mathcal{C}) \ar[r]^{f_p} \ar[d]^F &
\textit{PSh}(X, \mathcal{C}) \ar[d]^F &
\Sh(Y, \mathcal{C}) \ar[r]^{f^{-1}} \ar[d]^F &
\Sh(X, \mathcal{C}) \ar[d]^F
\\
\textit{PSh}(Y) \ar[r]^{f_p} &
\textit{PSh}(X) &
\Sh(Y) \ar[r]^{f^{-1}} &
\Sh(X)
}
$$
들이 있다.

\medskip\noindent
기억해야 할 주된 공식들은 다음과 같다.
\begin{eqnarray*}
f_*\mathcal{F}(V) & = & \mathcal{F}(f^{-1}(V)) \\
f_p\mathcal{G}(U) & = & \colim_{f(U) \subset V} \mathcal{G}(V) \\
f^{-1}\mathcal{G} & = & (f_p\mathcal{G})^\# \\
(f_p\mathcal{G})_x & = & \mathcal{G}_{f(x)} \\
(f^{-1}\mathcal{G})_x & = & \mathcal{G}_{f(x)}
\end{eqnarray*}
각 공식은 범주 $\mathcal{C}$에서 성립하고, 바탕 집합을 취하면 집합의
준층에 대한 대응하는 공식을 얻는다는 성질을 갖는다. 또한 수반 성질
\begin{eqnarray*}
\Mor_{\textit{PSh}(X, \mathcal{C})}(f_p\mathcal{G}, \mathcal{F})
& = &
\Mor_{\textit{PSh}(Y, \mathcal{C})}(\mathcal{G}, f_*\mathcal{F}) \\
\Mor_{\Sh(X, \mathcal{C})}(f^{-1}\mathcal{G}, \mathcal{F})
& = &
\Mor_{\Sh(Y, \mathcal{C})}(\mathcal{G}, f_*\mathcal{F}).
\end{eqnarray*}
들이 있다. 이를 증명하는 주된 단계는 보조정리
\ref{sheaves-lemma-pullback-presheaves}의 증명에 나오는 사상들
$$
i_\mathcal{G} : \mathcal{G} \longrightarrow f_*f_p\mathcal{G}
$$
과
$$
c_\mathcal{F} : f_p f_* \mathcal{F} \longrightarrow \mathcal{F}
$$
을 $\mathcal{C}$에 값을 갖는 준층들의 사상으로 구성하는 것이다.
그 구성은 집합의 준층의 경우와 정확히 같으므로 독자에게 맡겨도 안전하다.

\medskip\noindent
연속사상 $f : X \to Y$와 $X$ 위의 대수 구조의 층 $\mathcal{F}$,
$Y$ 위의 대수 구조의 층 $\mathcal{G}$가 주어지면 대수 구조의 층의
{\it $f$-사상 $\mathcal{G} \to \mathcal{F}$}이라는 개념이 의미를
갖는다. 정의 \ref{sheaves-definition-f-map}에서 집합의 사상을
$\mathcal{C}$의 사상으로 바꾸어 정확히 같은 방식으로 정의할 수도 있고,
단순히 대수 구조의 층의 사상
$\mathcal{G} \to f_*\mathcal{F}$와 같은 것이라고 말할 수도 있다.
아래에서는 이 개념을 자유롭게 사용할 것이다. $\mathcal{G}$와
$\mathcal{F}$ 사이의 $f$-사상들의 집합은 집합들
$\Mor_{\Sh(X, \mathcal{C})}(f^{-1}\mathcal{G}, \mathcal{F})$ 및
$\Mor_{\Sh(Y, \mathcal{C})}(\mathcal{G}, f_*\mathcal{F})$와
표준적으로 전단사이다.

\medskip\noindent
$f$-사상의 합성은 집합의 층의 $f$-사상의 경우와 정확히 같은 방식으로
정의한다. 더 나아가 위와 같은 $f$-사상
$\mathcal{G} \to \mathcal{F}$가 주어졌을 때 유도되는 줄기 위의 사상들
$$
\varphi_x : \mathcal{G}_{f(x)} \longrightarrow \mathcal{F}_x
$$
은 대수 구조의 준동형이다.


\begin{lemma}
\label{sheaves-lemma-f-map-sets-algebraic-structures}
$f : X \to Y$를 위상공간들의 연속사상이라 하자. $X$ 위의 대수 구조의
층 $\mathcal{F}$와 $Y$ 위의 대수 구조의 층 $\mathcal{G}$가 주어졌다고
하자. $\varphi : \mathcal{G} \to \mathcal{F}$를 바탕 집합의 층들의
$f$-사상이라 하자. 모든 열린집합 $V \subset Y$에 대하여 집합의 사상
$\varphi_V : \mathcal{G}(V) \to \mathcal{F}(f^{-1}V)$가
$\mathcal{C}$의 어떤 사상이 바탕 집합에 미치는 작용이면, $\varphi$는
대수 구조의 층들 사이의 유일한 $f$-사상으로부터 온다.
\end{lemma}

\begin{proof}
생략한다.
\end{proof}


\section{연속사상과 가군의 층}
\label{sheaves-section-presheaves-modules-functorial}

\noindent
가군의 층의 경우는 더 복잡하다. 그 이유는 역상 함자와 직접상 함자를
정의하는 자연스러운 무대가 아래에서 정의할 환 달린 공간이기 때문이다.
먼저 몇 가지 자명한 보조정리를 서술한다.

\begin{lemma}
\label{sheaves-lemma-pushforward-presheaf-module}
$f : X \to Y$를 위상공간들의 연속사상이라 하자. $\mathcal{O}$를 $X$
위의 환의 준층이라 하고, $\mathcal{F}$를 $\mathcal{O}$-가군의 준층이라
하자. 바탕 집합의 준층들의 자연스러운 사상
$$
f_*\mathcal{O} \times f_*\mathcal{F}
\longrightarrow
f_*\mathcal{F}
$$
이 있으며, 이는 $f_*\mathcal{F}$를 $f_*\mathcal{O}$-가군의 준층으로
만든다. 이 구성은 $\mathcal{F}$에 대해 함자적이다.
\end{lemma}

\begin{proof}
$V \subset Y$가 열려 있다고 하자. 보조정리의 사상을 사상
% P01-ERRATA: P01-E0047 -- the frozen source has "Let V ... is open" here and at line 2742; P01-E1259 independently records both loci.
$$
f_*\mathcal{O}(V) \times f_*\mathcal{F}(V)
=
\mathcal{O}(f^{-1}V) \times \mathcal{F}(f^{-1}V)
\to
\mathcal{F}(f^{-1}V)
=
f_*\mathcal{F}(V).
$$
으로 정의한다. 여기서 가운데 화살표는 $X$ 위의 곱셈 사상이다.
% P01-ERRATA: P01-E1260 -- this and the parallel map on Y are scalar action maps, though the frozen source calls them multiplication maps.
이것이 제한 사상들과 양립하고 $f_*\mathcal{F}$ 위에
$f_*\mathcal{O}$-가군 구조를 정의함을 확인하는 것은 독자에게 맡긴다.
\end{proof}

\begin{lemma}
\label{sheaves-lemma-pullback-presheaf-module}
$f : X \to Y$를 위상공간들의 연속사상이라 하자. $\mathcal{O}$를 $Y$
위의 환의 준층이라 하고, $\mathcal{G}$를 $\mathcal{O}$-가군의 준층이라
하자. 바탕 집합의 준층들의 자연스러운 사상
$$
f_p\mathcal{O} \times f_p\mathcal{G}
\longrightarrow
f_p\mathcal{G}
$$
이 있으며, 이는 $f_p\mathcal{G}$를 $f_p\mathcal{O}$-가군의 준층으로
만든다. 이 구성은 $\mathcal{G}$에 대해 함자적이다.
\end{lemma}

\begin{proof}
$U \subset X$가 열려 있다고 하자. 보조정리의 사상을 사상
% P01-ERRATA: P01-E0047 -- second frozen "Let U ... is open" locus.
\begin{eqnarray*}
f_p\mathcal{O}(U) \times f_p\mathcal{G}(U)
& = &
\colim_{f(U) \subset V} \mathcal{O}(V)
\times
\colim_{f(U) \subset V} \mathcal{G}(V) \\
& = &
\colim_{f(U) \subset V} (\mathcal{O}(V)\times \mathcal{G}(V)) \\
& \to &
\colim_{f(U) \subset V} \mathcal{G}(V) \\
& = &
f_p\mathcal{G}(U).
\end{eqnarray*}
으로 정의한다. 여기서 가운데 화살표는 $Y$ 위의 곱셈 사상이다.
두 번째 등식은 유향 여극한이 유한 극한과 가환하기 때문에 성립한다.
Categories, 보조정리 \ref{categories-lemma-directed-commutes}를 보라.
이것이 제한 사상들과 양립하고 $f_p\mathcal{G}$ 위에
$f_p\mathcal{O}$-가군 구조를 정의함을 확인하는 것은 독자에게 맡긴다.
\end{proof}

\noindent
$f : X \to Y$를 연속사상이라 하자. $\mathcal{O}_X$를 $X$ 위의 환의
준층이라 하고 $\mathcal{O}_Y$를 $Y$ 위의 환의 준층이라 하자.
현재 함자들
\begin{eqnarray*}
f_* : \textit{PMod}(\mathcal{O}_X) &
\longrightarrow &
\textit{PMod}(f_*\mathcal{O}_X) \\
f_p : \textit{PMod}(\mathcal{O}_Y) &
\longrightarrow &
\textit{PMod}(f_p\mathcal{O}_Y)
\end{eqnarray*}
을 정의했다. 이들은 다음과 같은 양립성을 만족한다.

\begin{lemma}
\label{sheaves-lemma-adjoint-push-pull-presheaves-modules}
$f : X \to Y$를 위상공간들의 연속사상이라 하자. $\mathcal{O}$를 $Y$
위의 환의 준층, $\mathcal{G}$를 $\mathcal{O}$-가군의 준층,
$\mathcal{F}$를 $f_p\mathcal{O}$-가군의 준층이라 하자. 그러면
$$
\Mor_{\textit{PMod}(f_p\mathcal{O})}(f_p\mathcal{G}, \mathcal{F})
=
\Mor_{\textit{PMod}(\mathcal{O})}(\mathcal{G}, f_*\mathcal{F}).
$$
이다. 여기서 보조정리 \ref{sheaves-lemma-pullback-presheaf-module}와
\ref{sheaves-lemma-pushforward-presheaf-module}를 사용하며, $f_*\mathcal{F}$를
사상 $i_\mathcal{O} : \mathcal{O} \to f_*f_p\mathcal{O}$를 통한
$\mathcal{O}$-가군으로 본다(이 사상은 보조정리
\ref{sheaves-lemma-pullback-presheaves}의 증명에서 처음 정의했다).
\end{lemma}

\begin{proof}
절 \ref{sheaves-section-abelian-presheaves-functorial}에 따라
$$
\Mor_{\textit{PAb}(X)}(f_p\mathcal{G}, \mathcal{F})
=
\Mor_{\textit{PAb}(Y)}(\mathcal{G}, f_*\mathcal{F}).
$$
이다. 따라서 이 대응 아래에서 가군 사상들의 부분집합이 대응함을
증명해야 한다. 또한 이 대응은 규칙
$$
(\psi : f_p\mathcal{G} \to \mathcal{F})
\longmapsto
(f_*\psi \circ i_\mathcal{G} :
\mathcal{G} \to f_* \mathcal{F})
$$
과 반대 방향의 규칙
$$
(\varphi : \mathcal{G} \to f_* \mathcal{F})
\longmapsto
(c_\mathcal{F} \circ f_p\varphi : f_p\mathcal{G} \to \mathcal{F})
$$
으로 결정된다. 여기서 $i_\mathcal{G}$와 $c_\mathcal{F}$는 절
\ref{sheaves-section-abelian-presheaves-functorial}에서와 같다. 따라서 $f_*$와
$f_p$의 함자성을 사용하면, 사상들
$i_\mathcal{G} : \mathcal{G} \to f_* f_p \mathcal{G}$와
$c_\mathcal{F} : f_p f_* \mathcal{F} \to \mathcal{F}$가 가군 구조와
양립함을 확인하는 것으로 충분하다. 이 확인은 독자에게 맡긴다.
\end{proof}

\begin{lemma}
\label{sheaves-lemma-adjoint-pull-push-presheaves-modules}
$f : X \to Y$를 위상공간들의 연속사상이라 하자. $\mathcal{O}$를 $X$
위의 환의 준층, $\mathcal{F}$를 $\mathcal{O}$-가군의 준층,
$\mathcal{G}$를 $f_*\mathcal{O}$-가군의 준층이라 하자. 그러면
$$
\Mor_{\textit{PMod}(\mathcal{O})}(
\mathcal{O} \otimes_{p, f_pf_*\mathcal{O}} f_p\mathcal{G}, \mathcal{F})
=
\Mor_{\textit{PMod}(f_*\mathcal{O})}(\mathcal{G}, f_*\mathcal{F}).
$$
이다. 여기서 보조정리 \ref{sheaves-lemma-pullback-presheaf-module}와
\ref{sheaves-lemma-pushforward-presheaf-module}를 사용하고, 텐서곱의 정의에서
사상 $c_\mathcal{O} : f_pf_*\mathcal{O} \to \mathcal{O}$를 사용한다.
\end{lemma}

\begin{proof}
이는 등식들
\begin{eqnarray*}
\Mor_{\textit{PMod}(\mathcal{O})}(
\mathcal{O} \otimes_{p, f_pf_*\mathcal{O}} f_p\mathcal{G}, \mathcal{F})
& = &
\Mor_{\textit{PMod}(f_pf_*\mathcal{O})}(
f_p\mathcal{G}, \mathcal{F}_{f_pf_*\mathcal{O}}) \\
& = &
\Mor_{\textit{PMod}(f_*\mathcal{O})}(\mathcal{G},
f_*(\mathcal{F}_{f_pf_*\mathcal{O}})) \\
& = &
\Mor_{\textit{PMod}(f_*\mathcal{O})}(\mathcal{G}, f_*\mathcal{F}).
\end{eqnarray*}
로부터 따른다. 첫 번째 등식은 보조정리
\ref{sheaves-lemma-adjointness-tensor-restrict-presheaves}이고, 두 번째 등식은
보조정리 \ref{sheaves-lemma-adjoint-push-pull-presheaves-modules}이다.
세 번째 등식은 $f_*\mathcal{O}$-선형인 아벨층들의 등식
$f_*(\mathcal{F}_{f_pf_*\mathcal{O}}) = f_*\mathcal{F}$로 주어진다.
% P01-ERRATA: P01-E1261 -- these objects are abelian presheaves in the surrounding PMod/PAb argument, though the frozen source says sheaves.
실제로 보조정리 \ref{sheaves-lemma-pullback-presheaves}의 증명에서 설명한 수반
아래에서 $\text{id}_{f_*\mathcal{O}}$는 $c_\mathcal{O}$에 대응하고,
따라서
$\text{id}_{f_*\mathcal{O}} = f_*c_\mathcal{O} \circ i_{f_*\mathcal{O}}$이다.
\end{proof}

\begin{lemma}
\label{sheaves-lemma-pushforward-module}
$f : X \to Y$를 위상공간들의 연속사상이라 하자. $\mathcal{O}$를 $X$
위의 환의 층이라 하고, $\mathcal{F}$를 $\mathcal{O}$-가군의 층이라
하자. 보조정리 \ref{sheaves-lemma-pushforward-presheaf-module}에서 정의한 직접상
$f_*\mathcal{F}$는 $f_*\mathcal{O}$-가군의 층이다.
\end{lemma}

\begin{proof}
정의와 보조정리 \ref{sheaves-lemma-pushforward-sheaf}로부터 자명하다.
\end{proof}

\begin{lemma}
\label{sheaves-lemma-pullback-module}
$f : X \to Y$를 위상공간들의 연속사상이라 하자. $\mathcal{O}$를 $Y$
위의 환의 층이라 하고, $\mathcal{G}$를 $\mathcal{O}$-가군의 층이라
하자. 바탕 집합의 준층들의 자연스러운 사상
$$
f^{-1}\mathcal{O} \times f^{-1}\mathcal{G}
\longrightarrow
f^{-1}\mathcal{G}
$$
이 있으며, 이는 $f^{-1}\mathcal{G}$를 $f^{-1}\mathcal{O}$-가군의
층으로 만든다.
\end{lemma}

\begin{proof}
$f^{-1}$이 함자 $f_p$와 층화의 합성으로 정의됨을 상기하자. 따라서 이
보조정리는 보조정리 \ref{sheaves-lemma-pullback-presheaf-module}와 보조정리
\ref{sheaves-lemma-sheafification-presheaf-modules}를 결합한 것이다.
\end{proof}

\noindent
$f : X \to Y$를 연속사상이라 하자. $\mathcal{O}_X$를 $X$ 위의 환의
층이라 하고 $\mathcal{O}_Y$를 $Y$ 위의 환의 층이라 하자. 이제 함자들
\begin{eqnarray*}
f_* : \textit{Mod}(\mathcal{O}_X) &
\longrightarrow &
\textit{Mod}(f_*\mathcal{O}_X) \\
f^{-1} : \textit{Mod}(\mathcal{O}_Y) &
\longrightarrow &
\textit{Mod}(f^{-1}\mathcal{O}_Y)
\end{eqnarray*}
을 정의했다. 이들은 다음과 같은 양립성을 만족한다.

\begin{lemma}
\label{sheaves-lemma-adjoint-push-pull-modules}
$f : X \to Y$를 위상공간들의 연속사상이라 하자. $\mathcal{O}$를 $Y$
위의 환의 층, $\mathcal{G}$를 $\mathcal{O}$-가군의 층,
$\mathcal{F}$를 $f^{-1}\mathcal{O}$-가군의 층이라 하자. 그러면
$$
\Mor_{\textit{Mod}(f^{-1}\mathcal{O})}(f^{-1}\mathcal{G}, \mathcal{F})
=
\Mor_{\textit{Mod}(\mathcal{O})}(\mathcal{G}, f_*\mathcal{F}).
$$
이다. 여기서 보조정리 \ref{sheaves-lemma-pullback-module}과
\ref{sheaves-lemma-pushforward-module}을 사용하고, $f_*\mathcal{F}$를
$\mathcal{O} \to f_*f^{-1}\mathcal{O}$를 통한 제한에 의해
$\mathcal{O}$-가군으로 본다.
\end{lemma}

\begin{proof}
등식들
\begin{eqnarray*}
\Mor_{\textit{Mod}(f^{-1}\mathcal{O})}(f^{-1}\mathcal{G}, \mathcal{F})
& = &
\Mor_{\textit{Mod}(f_p\mathcal{O})}(f_p\mathcal{G}, \mathcal{F}) \\
& = &
\Mor_{\textit{Mod}(\mathcal{O})}(\mathcal{G}, f_*\mathcal{F}).
\end{eqnarray*}
로 논증한다. 여기서 두 번째 등식은 보조정리
\ref{sheaves-lemma-adjoint-push-pull-presheaves-modules}이고 첫 번째 등식은
보조정리 \ref{sheaves-lemma-sheafification-presheaf-modules}에 의한 것이다.
% P01-ERRATA: P01-E1262 -- the frozen source begins this sentence with lowercase "where" and uses plural "Lemmas" for one cited lemma.
\end{proof}

\begin{lemma}
\label{sheaves-lemma-adjoint-pull-push-modules}
$f : X \to Y$를 위상공간들의 연속사상이라 하자. $\mathcal{O}$를 $X$
위의 환의 층, $\mathcal{F}$를 $\mathcal{O}$-가군의 층,
$\mathcal{G}$를 $f_*\mathcal{O}$-가군의 층이라 하자. 그러면
$$
\Mor_{\textit{Mod}(\mathcal{O})}(
\mathcal{O} \otimes_{f^{-1}f_*\mathcal{O}} f^{-1}\mathcal{G}, \mathcal{F})
=
\Mor_{\textit{Mod}(f_*\mathcal{O})}(\mathcal{G}, f_*\mathcal{F}).
$$
이다. 여기서 보조정리 \ref{sheaves-lemma-pullback-module}과
\ref{sheaves-lemma-pushforward-module}을 사용하고, 텐서곱의 정의에서 표준적인
사상 $f^{-1}f_*\mathcal{O} \to \mathcal{O}$를 사용한다.
\end{lemma}

\begin{proof}
이는 등식들
\begin{eqnarray*}
\Mor_{\textit{Mod}(\mathcal{O})}(
\mathcal{O} \otimes_{f^{-1}f_*\mathcal{O}} f^{-1}\mathcal{G}, \mathcal{F})
& = &
\Mor_{\textit{Mod}(f^{-1}f_*\mathcal{O})}(
f^{-1}\mathcal{G}, \mathcal{F}_{f^{-1}f_*\mathcal{O}}) \\
& = &
\Mor_{\textit{Mod}(f_*\mathcal{O})}(\mathcal{G}, f_*\mathcal{F}).
\end{eqnarray*}
로부터 따르며, 이는 보조정리 \ref{sheaves-lemma-adjointness-tensor-restrict}와
\ref{sheaves-lemma-adjoint-push-pull-modules}를 결합한 것이다.
\end{proof}

\noindent
$f : X \to Y$를 연속사상이라 하자. $\mathcal{O}_X$를 $X$ 위의 환의
(준)층이라 하고 $\mathcal{O}_Y$를 $Y$ 위의 환의 (준)층이라 하자.
현재 함자들
\begin{eqnarray*}
f_* : \textit{PMod}(\mathcal{O}_X) &
\longrightarrow &
\textit{PMod}(f_*\mathcal{O}_X) \\
f_* : \textit{Mod}(\mathcal{O}_X) &
\longrightarrow &
\textit{Mod}(f_*\mathcal{O}_X) \\
f_p : \textit{PMod}(\mathcal{O}_Y) &
\longrightarrow &
\textit{PMod}(f_p\mathcal{O}_Y) \\
f^{-1} : \textit{Mod}(\mathcal{O}_Y) &
\longrightarrow &
\textit{Mod}(f^{-1}\mathcal{O}_Y)
\end{eqnarray*}
을 정의했다. 일반적으로 가군의 층 위의 함자쌍 $(f_*, f^{-1})$은 그
표적 범주들이 일치하지 않으므로 분명 수반이 아니다. 즉, 위에서 본 것처럼
어떤 기적에 의해 환의 층 $\mathcal{O}_X, \mathcal{O}_Y$가 관계
$\mathcal{O}_X = f^{-1}\mathcal{O}_Y$와
$\mathcal{O}_Y = f_*\mathcal{O}_X$를 만족할 때만 작동한다.
이는 실제로 거의 결코 참이 아니다. 가군의 층 위에서 적절한 수반함자쌍이
존재하게 하는 올바른 사상 개념을 정의하기 위해 논의를 잠시 중단한다.


\section{환 달린 공간}
\label{sheaves-section-ringed-spaces}

\noindent
$X$를 위상공간이라 하고 $\mathcal{O}_X$를 $X$ 위의 환의 층이라 하자.
환의 층 $\mathcal{O}_X$를 $X$ 위의 함수들의 층으로 생각해야 한다.
또한 $f : X \to Y$가 “적절한” 사상이면 $Y$ 위의 함수는 합성에 의해
$X$ 위의 함수가 된다. 따라서 $\mathcal{O}_Y$에서
$\mathcal{O}_X$로 가는 자연스러운 $f$-사상이 있어야 한다. 정의
\ref{sheaves-definition-f-map}과 보조정리 \ref{sheaves-lemma-f-map}을 보라.
구체적인 예는 아래의 예 \ref{sheaves-example-continuous-map-ringed}을 보라.
관련된 추상적 정의는 다음과 같다.

\begin{definition}
\label{sheaves-definition-ringed-space}
{\it 환 달린 공간}은 위상공간 $X$와 $X$ 위의 환의 층
$\mathcal{O}_X$로 이루어진 쌍 $(X, \mathcal{O}_X)$이다.
{\it 환 달린 공간의 사상}
$(X, \mathcal{O}_X) \to (Y, \mathcal{O}_Y)$은 연속사상
$f : X \to Y$와 환의 층들의 $f$-사상
$f^\sharp : \mathcal{O}_Y \to \mathcal{O}_X$로 이루어진 쌍이다.
\end{definition}

\begin{example}
\label{sheaves-example-continuous-map-ringed}
$f : X \to Y$를 위상공간들의 연속사상이라 하자. $X$ 위와 $Y$ 위의
실수값 연속함수들의 층 $\mathcal{C}^0_X$와 $\mathcal{C}^0_Y$를
생각하자. 예 \ref{sheaves-example-C0-sheaf-rings}를 보라.
% P01-ERRATA: P01-E1317 -- the frozen source leaves the compound modifier "real valued" unhyphenated.
$f$에 딸린 자연스러운 $f$-사상
$f^\sharp : \mathcal{C}^0_Y \to \mathcal{C}^0_X$가 있다고 주장한다.
실제로 이를 단순히 규칙
\begin{eqnarray*}
\mathcal{C}^0_Y(V) & \longrightarrow & \mathcal{C}^0_X(f^{-1}V) \\
h & \longmapsto & h \circ f
\end{eqnarray*}
으로 정의한다. 엄밀히 말하면
$f^\sharp(h) = h \circ f|_{f^{-1}(V)}$라고 써야 한다.
이것이 정의 \ref{sheaves-definition-f-map}에서와 같은 사상족이며
$\mathbf{R}$-대수 구조들과 양립함은 분명하다. 따라서 이는
$\mathbf{R}$-대수의 층들의 $f$-사상이다. 보조정리
\ref{sheaves-lemma-f-map-sets-algebraic-structures}를 보라.

\medskip\noindent
물론 기하학적인 유형의 사상에 딸린 표준적인 환 달린 공간의 사상이
존재하는 상황은 이 밖에도 많다. 예를 들어 $M$, $N$이
$\mathcal{C}^\infty$-다양체이고 $f : M \to N$이 무한번 미분 가능한
사상이면, $f$는 표준적인 환 달린 공간의 사상
$(M, \mathcal{C}_M^\infty) \to (N, \mathcal{C}^\infty_N)$을 유도한다.
위와 동일한 이 구성은 독자에게 맡긴다.
\end{example}

\noindent
환 달린 공간의 사상을 합성하는 방법이 완전히 분명하지 않을 수도 있으므로
여기에서 상세히 적는다.

\begin{definition}
\label{sheaves-definition-composition-maps-ringed-spaces}
$(f, f^\sharp) : (X, \mathcal{O}_X) \to (Y, \mathcal{O}_Y)$와
$(g, g^\sharp) : (Y, \mathcal{O}_Y) \to (Z, \mathcal{O}_Z)$를 환 달린
공간의 사상이라 하자. {\it 환 달린 공간의 사상의 합성}을 규칙
$$
(g, g^\sharp) \circ (f, f^\sharp) = (g \circ f, f^\sharp \circ g^\sharp).
$$
으로 정의한다. 여기서 정의 \ref{sheaves-definition-composition-f-maps}에서
정의한 $f$-사상의 합성을 사용한다.
\end{definition}


\section{환 달린 공간의 사상과 가군}
\label{sheaves-section-ringed-spaces-functoriality-modules}

\noindent
이제 환 달린 공간의 사상을 따라 가군을 역상과 직접상으로 보내는 것을
정의하기에 충분한 표기를 도입했다.

\begin{definition}
\label{sheaves-definition-pushforward}
$(f, f^\sharp) : (X, \mathcal{O}_X) \to (Y, \mathcal{O}_Y)$를 환 달린
공간의 사상이라 하자.
\begin{enumerate}
\item $\mathcal{F}$를 $\mathcal{O}_X$-가군의 층이라 하자.
$\mathcal{F}$의 {\it 직접상}을 아벨군의 층으로서는
$f_*\mathcal{F}$와 같고, 그 가군 구조는 보조정리
\ref{sheaves-lemma-pushforward-module}에서 준 가군 구조를
$f^\sharp : \mathcal{O}_Y \to f_*\mathcal{O}_X$를 통해 제한하여
얻는 $\mathcal{O}_Y$-가군의 층으로 정의한다.
\item $\mathcal{G}$를 $\mathcal{O}_Y$-가군의 층이라 하자.
{\it 역상} $f^*\mathcal{G}$를 공식
$$
f^*\mathcal{G}
=
\mathcal{O}_X \otimes_{f^{-1}\mathcal{O}_Y} f^{-1}\mathcal{G}
$$
으로 정의되는 $\mathcal{O}_X$-가군의 층이라 한다. 여기서 환 사상
$f^{-1}\mathcal{O}_Y \to \mathcal{O}_X$는 $f^\sharp$에 대응하는
사상이고, 가군 구조는 보조정리 \ref{sheaves-lemma-pullback-module}로 주어진다.
% P01-ERRATA: P01-E1318 -- the frozen source contains a doubled space in "the  module structure".
\end{enumerate}
\end{definition}

\noindent
따라서 함자들
\begin{eqnarray*}
f_* : \textit{Mod}(\mathcal{O}_X)
& \longrightarrow &
\textit{Mod}(\mathcal{O}_Y) \\
f^* : \textit{Mod}(\mathcal{O}_Y)
& \longrightarrow &
\textit{Mod}(\mathcal{O}_X)
\end{eqnarray*}
을 정의했다. 이 함자들에 대한 마지막 결과는 예상대로 실제로 수반이라는
것이다.

\begin{lemma}
\label{sheaves-lemma-adjoint-pullback-pushforward-modules}
$(f, f^\sharp) : (X, \mathcal{O}_X) \to (Y, \mathcal{O}_Y)$를 환 달린
공간의 사상이라 하자. $\mathcal{F}$를 $\mathcal{O}_X$-가군의 층이라
하고 $\mathcal{G}$를 $\mathcal{O}_Y$-가군의 층이라 하자.
표준적인 전단사
$$
\Hom_{\mathcal{O}_X}(f^*\mathcal{G}, \mathcal{F})
=
\Hom_{\mathcal{O}_Y}(\mathcal{G}, f_*\mathcal{F}).
$$
가 있다. 달리 말하면 함자 $f^*$는 $f_*$의 왼쪽 수반이다.
\end{lemma}

\begin{proof}
이는 앞에서 한 작업으로부터 따른다.
\begin{eqnarray*}
\Hom_{\mathcal{O}_X}(f^*\mathcal{G}, \mathcal{F})
& = &
\Mor_{\textit{Mod}(\mathcal{O}_X)}(
\mathcal{O}_X \otimes_{f^{-1}\mathcal{O}_Y} f^{-1}\mathcal{G}, \mathcal{F}) \\
& = &
\Mor_{\textit{Mod}(f^{-1}\mathcal{O}_Y)}(
f^{-1}\mathcal{G}, \mathcal{F}_{f^{-1}\mathcal{O}_Y}) \\
& = &
\Hom_{\mathcal{O}_Y}(\mathcal{G}, f_*\mathcal{F}).
\end{eqnarray*}
여기서 보조정리 \ref{sheaves-lemma-adjointness-tensor-restrict}와
\ref{sheaves-lemma-adjoint-push-pull-modules}를 사용한다.
\end{proof}

\begin{lemma}
\label{sheaves-lemma-push-pull-composition-modules}
$f : X \to Y$와 $g : Y \to Z$를 환 달린 공간의 사상이라 하자.
함자 $(g \circ f)_*$와 $g_* \circ f_*$는 같다. 함자들의 표준적인
동형사상 $(g \circ f)^* \cong f^* \circ g^*$가 있다.
\end{lemma}

\begin{proof}
직접상에 대한 결과는 보조정리 \ref{sheaves-lemma-pushforward-composition}와
우리의 정의들로부터 따른다. 역상에 대한 결과는 보조정리
\ref{sheaves-lemma-pullback-composition}의 증명과 같은 논증으로 이로부터 따른다.
\end{proof}


\noindent
환 달린 공간의 사상
$(f, f^\sharp) : (X, \mathcal{O}_X) \to (Y, \mathcal{O}_Y)$와
$\mathcal{O}_X$-가군의 층 $\mathcal{F}$, 그리고 $Y$ 위의
$\mathcal{O}_Y$-가군의 층 $\mathcal{G}$가 주어지면
가군의 층의 {\it $f$-사상 $\varphi : \mathcal{G} \to \mathcal{F}$}
이라는 개념이 의미를 갖는다. 이를 모든 열린집합 $V \subset Y$에 대해
사상
$$
\mathcal{G}(V) \longrightarrow \mathcal{F}(f^{-1}V)
$$
이 $\mathcal{O}_Y(V)$-가군 사상이 되는 아벨층의 $f$-사상
$\varphi : \mathcal{G} \to \mathcal{F}$로 정의하면 된다
(정의 \ref{sheaves-definition-f-map}과 보조정리 \ref{sheaves-lemma-f-map}을 보라).
여기서 $\mathcal{F}(f^{-1}V)$를 사상
$f^\sharp_V : \mathcal{O}_Y(V) \to \mathcal{O}_X(f^{-1}V)$를 통한
$\mathcal{O}_Y(V)$-가군으로 본다. $\mathcal{G}$와 $\mathcal{F}$ 사이의
$f$-사상들의 집합은 집합들
$\Mor_{\textit{Mod}(\mathcal{O}_X)}(f^*\mathcal{G}, \mathcal{F})$ 및
$\Mor_{\textit{Mod}(\mathcal{O}_Y)}(\mathcal{G}, f_*\mathcal{F})$와
표준적으로 전단사이다. 위를 보라.

\medskip\noindent
$f$-사상의 합성은 집합의 층의 $f$-사상의 경우와 정확히 같은 방식으로
정의한다. 더 나아가 위와 같은 $f$-사상
$\mathcal{G} \to \mathcal{F}$와 $x \in X$가 주어졌을 때 유도되는
줄기 위의 사상
$$
\varphi_x : \mathcal{G}_{f(x)} \longrightarrow \mathcal{F}_x
$$
은 $\mathcal{O}_{Y, f(x)}$-가군 사상이다. 여기서
$\mathcal{F}_x$ 위의 $\mathcal{O}_{Y, f(x)}$-가군 구조는 사상
$f^\sharp_x : \mathcal{O}_{Y, f(x)} \to \mathcal{O}_{X, x}$를 통해
$\mathcal{O}_{X, x}$-가군 구조로부터 온다. 관련된 보조정리는 다음과
같다.

\begin{lemma}
\label{sheaves-lemma-stalk-pullback-modules}
$(f, f^\sharp) : (X, \mathcal{O}_X) \to (Y, \mathcal{O}_Y)$를 환 달린
공간의 사상이라 하자. $\mathcal{G}$를 $\mathcal{O}_Y$-가군의 층이라
하고 $x \in X$라 하자. 그러면 오른쪽 텐서곱에서
$f^\sharp_x : \mathcal{O}_{Y, f(x)} \to \mathcal{O}_{X, x}$를 사용할 때,
$\mathcal{O}_{X, x}$-가군으로서
$$
(f^*\mathcal{G})_x =
\mathcal{G}_{f(x)}
\otimes_{\mathcal{O}_{Y, f(x)}}
\mathcal{O}_{X, x}
$$
이다.
\end{lemma}

\begin{proof}
이는 보조정리 \ref{sheaves-lemma-stalk-tensor-sheaf-modules}와, $x$에서의 역상층의
줄기를 $f(x)$에서의 대응하는 줄기와 동일시하는 것으로부터 따른다.
예를 들어 절 \ref{sheaves-section-presheaves-structures-functorial}의 공식들을 보라.
\end{proof}



\section{마천루층과 줄기}
\label{sheaves-section-skyscraper-sheaves}

\begin{definition}
\label{sheaves-definition-skyscraper-sheaf}
$X$를 위상공간이라 하자.
\begin{enumerate}
\item $x \in X$를 점이라 하고, 포함 사상을
$i_x : \{x\} \to X$로 나타내자.
% P01-ERRATA: P01-E1319 -- the frozen source says "Denote i_x ... the inclusion map" without the required preposition.
$A$를 집합이라 하고 $A$를 한 점 공간 $\{x\}$ 위의 층으로 생각하자.
$i_{x, *}A$를 {\it $x$에서 값이 $A$인 마천루층}이라 한다.
\item 위의 (1)에서 $A$가 아벨군이면 $i_{x, *}A$를 $X$ 위의 아벨군의
층으로 생각한다.
\item 위의 (1)에서 $A$가 대수 구조이면 $i_{x, *}A$를 대수 구조의
층으로 생각한다.
\item $(X, \mathcal{O}_X)$가 환 달린 공간이면
$i_x : \{x\} \to X$를 환 달린 공간의 사상
$(\{x\}, \mathcal{O}_{X, x}) \to (X, \mathcal{O}_X)$로 생각한다.
$A$가 $\mathcal{O}_{X, x}$-가군이면 $i_{x, *}A$를
$\mathcal{O}_X$-가군의 층으로 생각한다.
\item $X$의 점 $x$와 집합 $A$가 존재하여
$\mathcal{F} \cong i_{x, *}A$이면 집합의 층 $\mathcal{F}$를
{\it 마천루층}이라 한다.
\item $X$의 점 $x$와 아벨군 $A$가 존재하여 아벨군의 층으로서
$\mathcal{F} \cong i_{x, *}A$이면 아벨군의 층 $\mathcal{F}$를
{\it 마천루층}이라 한다.
\item $X$의 점 $x$와 대수 구조 $A$가 존재하여 대수 구조의 층으로서
$\mathcal{F} \cong i_{x, *}A$이면 대수 구조의 층 $\mathcal{F}$를
{\it 마천루층}이라 한다.
\item $(X, \mathcal{O}_X)$가 환 달린 공간이고 $\mathcal{F}$가
$\mathcal{O}_X$-가군의 층일 때, 점 $x \in X$와
$\mathcal{O}_{X, x}$-가군 $A$가 존재하여 $\mathcal{O}_X$-가군의
층으로서 $\mathcal{F} \cong i_{x, *}A$이면 $\mathcal{F}$를
{\it 마천루층}이라 한다.
\end{enumerate}
\end{definition}

\begin{lemma}
\label{sheaves-lemma-skyscraper-stalks}
$X$를 위상공간, $x \in X$를 점, $A$를 집합이라 하자. 임의의 점
$x' \in X$에 대하여 $x$에서 값이 $A$인 마천루층의 $x'$에서의 줄기는
$$
(i_{x, *}A)_{x'} =
\left\{
\begin{matrix}
A & \text{if} & x' \in \overline{\{x\}} \\
\{*\} & \text{if} & x' \not\in \overline{\{x\}}
\end{matrix}
\right.
$$
이다. 아벨군, 대수 구조, 가군의 층의 경우에도 비슷한 기술이 성립한다.
\end{lemma}

\begin{proof}
생략한다.
\end{proof}

\begin{lemma}
\label{sheaves-lemma-stalk-skyscraper-adjoint}
$X$를 위상공간이라 하고 $x \in X$를 점이라 하자.
% P01-ERRATA: P01-E1320 -- the frozen source says "let x ... a point" without "be".
함자 $\mathcal{F} \mapsto \mathcal{F}_x$와 $A \mapsto i_{x, *}A$는
수반이다. 공식으로는
$$
\Mor_{\textit{Sets}}(\mathcal{F}_x, A)
=
\Mor_{\Sh(X)}(\mathcal{F}, i_{x, *}A).
$$
이다. 아벨군과 대수 구조의 경우에도 비슷한 명제가 성립한다.
% P01-ERRATA: P01-E1321 -- the frozen source's list "abelian groups, algebraic structures" lacks a conjunction.
가군의 층의 경우에는
$$
\Hom_{\mathcal{O}_{X, x}}(\mathcal{F}_x, A)
=
\Hom_{\mathcal{O}_X}(\mathcal{F}, i_{x, *}A).
$$
이다.
\end{lemma}

\begin{proof}
생략한다. 힌트: 줄기 함자는 사상 $i_x : \{x\} \to X$에 대한 역상
함자로 볼 수 있다. 그러면 수반성은 $i_x^{-1}$과 $i_{x, *}$의
수반성(가군의 층의 경우에는 각각 $i_x^*$와 $i_{x, *}$의 수반성)으로부터
따른다.
\end{proof}



\section{준층의 극한과 여극한}
\label{sheaves-section-limits-presheaves}

\noindent
$X$를 위상공간이라 하자. 도표
$\mathcal{I} \to \textit{PSh}(X)$, $i \mapsto \mathcal{F}_i$가
주어졌다고 하자.
\begin{enumerate}
\item $\lim_i \mathcal{F}_i$와 $\colim_i \mathcal{F}_i$가 모두 존재한다.
\item 임의의 열린집합 $U \subset X$에 대하여
$$
(\lim_i \mathcal{F}_i)(U) =
\lim_i \mathcal{F}_i(U)
$$
이고
$$
(\colim_i \mathcal{F}_i)(U) =
\colim_i \mathcal{F}_i(U).
$$
\item $x \in X$를 점이라 하자. 일반적으로 $\lim_i \mathcal{F}_i$의
$x$에서의 줄기는 줄기들의 극한과 같지 않다. 그러나 지표범주가 유한하면
같다. 다시 말해 줄기 함자는 좌완전이다(Categories, 정의
\ref{categories-definition-exact}를 보라).
\item $x \in X$라 하자. 항상
$$
(\colim_i \mathcal{F}_i)_x =
\colim_i \mathcal{F}_{i, x}.
$$
이다.
\end{enumerate}
증명은 모두 쉽다.

\section{층의 극한과 여극한}
\label{sheaves-section-limits-sheaves}

\noindent
$X$를 위상공간이라 하자. 도표
$\mathcal{I} \to \Sh(X)$, $i \mapsto \mathcal{F}_i$가 주어졌다고 하자.
\begin{enumerate}
\item $\lim_i \mathcal{F}_i$와 $\colim_i \mathcal{F}_i$가 모두 존재한다.
\item 포함 함자 $i : \Sh(X) \to \textit{PSh}(X)$는 극한과 가환한다.
다시 말해 층의 범주에서의 극한을 준층의 범주에서의 극한으로 계산할 수
있다. 특히 임의의 열린집합 $U \subset X$에 대하여
$$
(\lim_i \mathcal{F}_i)(U) = \lim_i \mathcal{F}_i(U).
$$
이다.
\item 포함 함자 $i : \Sh(X) \to \textit{PSh}(X)$는 일반적으로
여극한과 가환하지 않는다(유한 여극한과도 가환하지 않는다---전사사상을
생각하라). 여극한은 준층의 범주에서의 여극한을 층화하여 계산한다:
$$
\colim_i \mathcal{F}_i =
\Big(U \mapsto \colim_i \mathcal{F}_i(U)\Big)^\#.
$$
\item $x \in X$를 점이라 하자. 일반적으로 $\lim_i \mathcal{F}_i$의
$x$에서의 줄기는 줄기들의 극한과 같지 않다. 그러나 지표범주가 유한하면
같다. 다시 말해 줄기 함자는 좌완전이다.
\item $x \in X$라 하자. 항상
$$
(\colim_i \mathcal{F}_i)_x = \colim_i \mathcal{F}_{i, x}.
$$
이다.
\item 층화 함자
${}^\# : \textit{PSh}(X) \to \Sh(X)$는 모든 여극한 및 유한 극한과
가환하지만 모든 극한과 가환하지는 않는다.
\end{enumerate}
증명은 모두 쉽다. 다음은 층의 유향 여극한에 대해 성립하는 명제의 한 예이다.

\begin{lemma}
\label{sheaves-lemma-directed-colimits-sections}
$X$를 위상공간이라 하고 $I$를 유향 집합이라 하자.
$(\mathcal{F}_i, \varphi_{ii'})$를 $I$ 위의 집합의 층들의 계라 하자.
Categories, Section \ref{categories-section-posets-limits}를 보라.
$U \subset X$를 열린 부분집합이라 하자. 표준 사상
$$
\Psi :
\colim_i \mathcal{F}_i(U)
\longrightarrow
\left(\colim_i \mathcal{F}_i\right)(U)
$$
을 생각하자.
\begin{enumerate}
\item 모든 전이 사상이 단사이면 임의의 열린집합 $U$에 대하여
$\Psi$는 단사이다.
\item $U$가 준콤팩트이면 $\Psi$는 단사이다.
\item $U$가 준콤팩트이고 모든 전이 사상이 단사이면 $\Psi$는 동형사상이다.
\item $U$가 열린 덮개
$\mathcal{U} : U = \bigcup_{j\in J} U_j$들의 공종계를 가지며,
$J$가 유한하고 모든 $j, j' \in J$에 대하여
$U_j \cap U_{j'}$가 준콤팩트이면 $\Psi$는 전단사이다.
\end{enumerate}
\end{lemma}

\begin{proof}
모든 전이 사상이 단사라고 가정하자. 이 경우 준층
$\mathcal{F}' : V \mapsto \colim_i \mathcal{F}_i(V)$는 분리되어 있다
(정의 \ref{sheaves-definition-separated}를 보라). 위의 논의에 의해
$(\mathcal{F}')^\# = \colim_i \mathcal{F}_i$이다. 보조정리
\ref{sheaves-lemma-separated-presheaf-into-sheaf}에 의해
$\mathcal{F}' \to (\mathcal{F}')^\#$가 단사임을 알 수 있다.
이로써 (1)이 증명된다.

\medskip\noindent
$U$가 준콤팩트라고 가정하자. $s \in \mathcal{F}_i(U)$와
$s' \in \mathcal{F}_{i'}(U)$가 왼쪽의 원소들을 주며, 이 원소들이
$\Psi$에 의해 같은 상으로 간다고 하자. $U$가 준콤팩트이므로 이는
유한 열린 덮개 $U = \bigcup_{j = 1, \ldots, m} U_j$가 존재하고 각
$j$마다 지표 $i_j \in I$, $i_j \geq i$, $i_j \geq i'$가 존재하여
$\varphi_{ii_j}(s)$와 $\varphi_{i'i_j}(s')$의 $U_j$로의 제한이 같다는
뜻이다. 모든 $i_j$에 대해 $\geq$인 $i''\in I$를 택하자. 모든 $j$에 대하여
$\varphi_{ii''}(s)$와 $\varphi_{i'i''}(s')$가 열린집합 $U_j$ 위에서
일치하므로 $\varphi_{ii''}(s) = \varphi_{i'i''}(s')$이다.
이로써 (2)가 증명된다.

\medskip\noindent
$U$가 준콤팩트이고 모든 전이 사상이 단사라고 가정하자.
$s$를 $\Psi$의 공역의 원소라 하자. $U$가 준콤팩트이므로 유한 열린 덮개
$U = \bigcup_{j = 1, \ldots, m} U_j$가 존재하고, 각 $j$마다 지표
$i_j \in I$와 $s_j \in \mathcal{F}_{i_j}(U_j)$가 존재하여 모든 $j$에
대하여 $s|_{U_j}$가 $s_j$에서 온다. 모든 $i_j$에 대해 $\geq$인 $i \in I$를
택하자. (1)에 의해 각 절단 $\varphi_{i_ji}(s_j)$는 교집합
$U_j \cap U_{j'}$ 위에서 일치한다. 따라서 이들은
$s' \in \mathcal{F}_i(U)$인 절단으로 붙고, 이 절단은 $\Psi$에 의해
$s$로 간다. 이로써 (3)이 증명된다.

\medskip\noindent
(4)의 가정을 하자. 특히 $U$는 준콤팩트이므로 (2)에 의해 $\Psi$는
단사이다. $s$를 $\Psi$의 공역의 원소라 하자. 가정에 의해 유한 열린 덮개
$U = \bigcup_{j = 1, \ldots, m} U_j$가 존재하며, 모든
$j, j' \in J$에 대하여 $U_j \cap U_{j'}$는 준콤팩트이고, 각 $j$마다
지표 $i_j \in I$와 $s_j \in \mathcal{F}_{i_j}(U_j)$가 존재하여 모든
$j$에 대해 $s|_{U_j}$는 $s_j$의 상이다. $U_j \cap U_{j'}$가
준콤팩트이므로 (2)를 적용할 수 있고, $i_{jj'} \in I$가 존재하여
$i_{jj'} \geq i_j$, $i_{jj'} \geq i_{j'}$이며
$\varphi_{i_ji_{jj'}}(s_j)$와 $\varphi_{i_{j'}i_{jj'}}(s_{j'})$가
$U_j \cap U_{j'}$ 위에서 일치함을 알 수 있다. 모든 $i_{jj'}$ 이상인
지표 $i \in I$를 택하자. 그러면 $\mathcal{F}_i$의 절단들
$\varphi_{i_ji}(s_j)$가 $U$ 위의 $\mathcal{F}_i$의 절단으로 붙는다.
이 절단은 요구한 대로 원소 $s$로 간다.
\end{proof}

\begin{example}
\label{sheaves-example-conditions-needed-colimit}
$X = \{s_1, s_2, \xi_1, \xi_2, \xi_3, \ldots\}$를 집합이라 하자.
부분집합 $U \subset X$가 열린집합이라는 것을, $s_1 \in U$ 또는
$s_2 \in U$이면 $U$가 모든 $\xi_i$를 포함한다는 것으로 정의하자.
$U_n = \{\xi_n, \xi_{n + 1}, \ldots\}$라 하고,
$j_n : U_n \to X$를 포함 사상이라 하자.
$\mathcal{F}_n = j_{n, *}\underline{\mathbf{Z}}$라 놓자. 전이 사상
$\mathcal{F}_n \to \mathcal{F}_{n + 1}$이 있다.
$\mathcal{F} = \colim \mathcal{F}_n$이라 놓자. $\{\xi_m\}$은 $U_n$과
만나지 않는 $X$의 열린 부분집합이므로 $m < n$이면
$\mathcal{F}_{n, \xi_m} = 0$임에 유의하라. 따라서 모든 $n$에 대하여
$\mathcal{F}_{\xi_n} = 0$임을 알 수 있다. 한편 $i = 1, 2$일 때 줄기
$\mathcal{F}_{s_i}$는 여극한
$$
M = \colim_n \prod\nolimits_{m \geq n} \mathbf{Z}
$$
이고, 이는 영이 아니다. 따라서 층 $\mathcal{F}$는 닫힌점 $s_1$과
$s_2$에서 값이 $M$인 마천루층들의 직합이다. 따라서
$\Gamma(X, \mathcal{F}) = M \oplus M$이다. 한편 독자는
$\colim_n \Gamma(X, \mathcal{F}_n) = M$임을 확인할 수 있다.
따라서 위의 보조정리 \ref{sheaves-lemma-directed-colimits-sections}의 (4)에는
어떤 조건이 필요하다.
\end{example}

\noindent
앞 보조정리에는 스펙트럴 공간들의 도표 위의 층을 다루는 판본이 있다.
이를 진술하기 위해 몇 가지 표기를 도입하자. $\mathcal{I}$를 쌍대여과
지표범주라 하자. $i \mapsto X_i$를 $\mathcal{I}$ 위의 스펙트럴 공간들의
도표라 하고, $\mathcal{I}$의 $a : j \to i$에 대응하는 사상
$f_a : X_j \to X_i$가 스펙트럴 사상이라고 하자. $X = \lim X_i$라 놓고
$p_i : X \to X_i$를 사영사상이라 하자.

\begin{lemma}
\label{sheaves-lemma-compute-pullback-to-limit}
위의 상황에서 $i \in \Ob(\mathcal{I})$라 하고 $\mathcal{G}$를
$X_i$ 위의 층이라 하자. 준콤팩트 열린집합 $U_i \subset X_i$에 대하여
$$
p_i^{-1}\mathcal{G}(p_i^{-1}(U_i)) =
\colim_{a : j \to i} f_a^{-1}\mathcal{G}(f_a^{-1}(U_i))
$$
이다.
\end{lemma}

\begin{proof}
표준 사상
$\colim_{a : j \to i} f_a^{-1}\mathcal{G}(f_a^{-1}(U_i)) \to
p_i^{-1}\mathcal{G}(p_i^{-1}(U_i))$가 단사임을 증명하자.
어떤 $a : j \to i$에 대하여 $s, s'$를 $f_a^{-1}(U_i)$ 위의
$f_a^{-1}\mathcal{G}$의 절단이라 하자. $b : k \to j$에 대하여
$Z_k \subset f_{a \circ b}^{-1}(U_i)$를, $s$와 $s'$의 상이 줄기
$(f_{a \circ b}^{-1}\mathcal{G})_x$에서 서로 다른 점 $x$들의 닫힌
부분집합이라 하자. 모든 $b : k \to j$에 대하여 $Z_k$가 비어 있지
않다면 Topology, Lemma
\ref{topology-lemma-inverse-limit-spectral-spaces-nonempty}에 의해
$\lim_{b : k \to j} Z_k$도 비어 있지 않음을 알 수 있다. 그러면
$x \in \lim_{b : k \to j} Z_k \subset X$에 대하여
($\mathcal{I}/j \to \mathcal{I}$가 초기임에 유의하라) 위의 모든
$b : k \to j$에 대해
% P01-ERRATA: P01-E0048 -- frozen source reverses the composable arrows here as b\circ a; retain source formula.
$(p_i^{-1}\mathcal{G})_x = (f_{b \circ a}^{-1}\mathcal{G})_{p_k(x)}$
이므로 $s$와 $s'$의 상은 $p_i^{-1}\mathcal{G}$의 $x$에서의 줄기에서도
서로 다르다. 따라서 $s$와 $s'$의 상이
$p_i^{-1}\mathcal{G}(p_i^{-1}(U_i))$에서 같다면 어떤
$b : k \to j$에 대하여 $Z_k$는 공집합이다. 이로써 단사성이 증명된다.

\medskip\noindent
전사성을 보이자. $s$를 $p_i^{-1}(U_i)$ 위의
$p_i^{-1}\mathcal{G}$의 절단이라 하자. Topology, Lemma
\ref{topology-lemma-directed-inverse-limit-spectral-spaces}에 의해
집합 $p_i^{-1}(U_i)$는 스펙트럴 공간 $X$의 준콤팩트 열린집합이다.
역상층의 구성에 의해 열린 덮개
$p_i^{-1}(U_i) = \bigcup_{l \in L} W_l$, 열린집합
$V_{l, i} \subset X_i$, 절단 $s_{l, i} \in \mathcal{G}(V_{l, i})$를
찾을 수 있어서 $p_i(W_l) \subset V_{l, i}$이고
$p_i^{-1}s_{l, i}|_{W_l} = s|_{W_l}$이다. $X$와 $X_i$가 스펙트럴이고
$p_i^{-1}(U_i)$가 준콤팩트 열린집합이므로 $L$이 유한하며 모든 $l$에
대하여 $W_l$과 $V_{l, i}$가 준콤팩트 열린집합이라고 가정할 수 있다.
그러면 Topology, Lemma \ref{topology-lemma-descend-opens}를 적용하여
$a : j \to i$와 준콤팩트 열린집합들에 의한 열린 덮개
$f_a^{-1}(U_i) = \bigcup_{l \in L} W_{l, j}$를 찾을 수 있다. 이 덮개의
$X$로의 역상은 덮개 $p_i^{-1}(U_i) = \bigcup_{l \in L} W_l$이고, 또한
$W_{l, j} \subset f_a^{-1}(V_{l, i})$이다. $s_{l, j}$를 $s_{l, i}$의
$f_a$에 의한 역상을 $W_{l, j}$로 제한한 것으로 쓰자. 그러면
$s_{l, j}$와 $s_{l', j}$는
$(f_a^{-1}\mathcal{G})(W_{l, j} \cap W_{l', j})$의 원소들로 제한되며,
그 역상들은 $(p_i^{-1}\mathcal{G})(W_l \cap W_{l'})$의 같은 원소,
즉 $s$의 제한이 된다. 따라서 단사성에 의해, 절단들
$f_b^{-1}s_{l, j}$가 요구한 대로 $f_{a \circ b}^{-1}(U_i)$ 위의
절단으로 붙게 하는 $b : k \to j$를 찾을 수 있다.
\end{proof}

\noindent
다음으로 스펙트럴 공간들의 쌍대여과계 $X_i$에 더하여 다음이 주어졌다고
가정하자.
\begin{enumerate}
\item 모든 $i \in \Ob(\mathcal{I})$에 대하여 $X_i$ 위의 층
$\mathcal{F}_i$,
\item $a : j \to i$에 대하여 $f_a$-사상
$\varphi_a : \mathcal{F}_i \to \mathcal{F}_j$.
\end{enumerate}
$c = a \circ b$일 때마다 $\varphi_c = \varphi_b \circ \varphi_a$라고
가정하자. $X$ 위에서 $\mathcal{F} = \colim p_i^{-1}\mathcal{F}_i$라 놓자.

\begin{lemma}
\label{sheaves-lemma-descend-opens}
위의 상황에서 $i \in \Ob(\mathcal{I})$라 하고
$U_i \subset X_i$를 준콤팩트 열린집합이라 하자. 그러면
$$
\colim_{a : j \to i} \mathcal{F}_j(f_a^{-1}(U_i)) = \mathcal{F}(p_i^{-1}(U_i))
$$
이다.
\end{lemma}

\begin{proof}
$p_i^{-1}(U_i)$가 스펙트럴 공간 $X$의 준콤팩트 열린집합임을 상기하라.
Topology, Lemma
\ref{topology-lemma-directed-inverse-limit-spectral-spaces}를 보라.
따라서 보조정리 \ref{sheaves-lemma-directed-colimits-sections}를 적용할 수 있고
$$
\mathcal{F}(p_i^{-1}(U_i)) =
\colim_{a : j \to i} p_j^{-1}\mathcal{F}_j(p_i^{-1}(U_i)).
$$
이다. 형식적인 논증으로
$$
\colim_{a : j \to i} \mathcal{F}_j(f_a^{-1}(U_i)) =
\colim_{a : j \to i} \colim_{b : k \to j}
f_b^{-1}\mathcal{F}_j(f_{a \circ b}^{-1}(U_i))
$$
임을 알 수 있다. 따라서
$$
p_j^{-1}\mathcal{F}_j(p_i^{-1}(U_i)) =
\colim_{b : k \to j} f_b^{-1}\mathcal{F}_j(f_{a \circ b}^{-1}(U_i))
$$
임을 보이면 충분하다. 이는 $\mathcal{F}_j$와 준콤팩트 열린집합
$f_a^{-1}(U_i)$에 적용한 보조정리
\ref{sheaves-lemma-compute-pullback-to-limit}이다.
\end{proof}

















\section{기저와 층}
\label{sheaves-section-bases}

\noindent
때로는 일반적인 열린집합보다 다루기 쉬운 열린집합들로 이루어진 위상의
기저가 존재한다. 이러한 상황에서 (준)층을 다루기 쉽도록 몇 가지 정의와
간단한 보조정리를 편의를 위해 제시한다.

\begin{definition}
\label{sheaves-definition-presheaf-basis}
$X$를 위상공간이라 하자. $\mathcal{B}$를 $X$의 위상의 기저라 하자.
\begin{enumerate}
\item {\it $\mathcal{B}$ 위의 집합의 준층 $\mathcal{F}$}란 각
$U \in \mathcal{B}$에 집합 $\mathcal{F}(U)$를 대응시키고,
$\mathcal{B}$의 원소들의 각 포함관계 $V \subset U$에 사상
$\rho^U_V : \mathcal{F}(U) \to \mathcal{F}(V)$를 대응시키는 규칙으로서,
모든 $U \in \mathcal{B}$에 대해 $\rho^U_U = \text{id}_{\mathcal{F}(U)}$이고,
$\mathcal{B}$에서 $W \subset V \subset U$일 때마다
$\rho^U_W = \rho^V_W \circ \rho ^U_V$인 것을 말한다.
\item {\it $\mathcal{B}$ 위의 집합의 준층의 사상
$\varphi : \mathcal{F} \to \mathcal{G}$}이란 각 원소
$U \in \mathcal{B}$에 제한 사상들과 가환하는 집합의 사상
% P01-ERRATA: P01-E0049 -- frozen source writes each component as phi rather than phi_U.
$\varphi : \mathcal{F}(U) \to \mathcal{G}(U)$를 대응시키는 규칙이다.
\end{enumerate}
\end{definition}

\noindent
보통의 준층의 경우와 마찬가지로 절단, 절단의 제한 등의 용어를 쓴다.
특히 점 $x \in X$에서 $\mathcal{F}$의 {\it 줄기}를 여극한
$$
\mathcal{F}_x = \colim_{U\in \mathcal{B}, x\in U} \mathcal{F}(U).
$$
으로 정의할 수 있다. $X$ 위의 준층의 줄기와 마찬가지로 이 극한은
유향이다. 이는 $U\in \mathcal{B}$, $x \in U$인 열린집합들의 모임이
$x$의 열린 근방의 기본계를 이루기 때문이다.

\medskip\noindent
$X$ 위의 준층이라는 개념이 $X$의 위상의 기저 위의 준층이라는 개념과
매우 다름을 보여 주는 예를 만들기는 쉽다. 층의 경우에는 이런 일이
일어나지 않는다. 따라서 다음 개념이 훨씬 더 유용하다.

\begin{definition}
\label{sheaves-definition-sheaf-basis}
$X$를 위상공간이라 하자. $\mathcal{B}$를 $X$의 위상의 기저라 하자.
\begin{enumerate}
\item {\it $\mathcal{B}$ 위의 집합의 층 $\mathcal{F}$}란 다음 추가
성질을 만족하는 $\mathcal{B}$ 위의 집합의 준층이다. 임의의
$U \in \mathcal{B}$와 $U_i \in \mathcal{B}$인 임의의 덮개
$U = \bigcup_{i \in I} U_i$, 그리고 $U_{ijk} \in \mathcal{B}$인
임의의 덮개 $U_i \cap U_j = \bigcup_{k \in I_{ij}} U_{ijk}$가
주어졌을 때 다음 층 조건이 성립한다.
\begin{itemize}
\item[(**)] 절단들의 임의의 모임 $s_i \in \mathcal{F}(U_i)$,
$i \in I$가 주어져서 $\forall i, j\in I$, $\forall k\in I_{ij}$에 대해
$$
s_i|_{U_{ijk}} = s_j|_{U_{ijk}}
$$
이면, 모든 $i \in I$에 대하여 $s_i = s|_{U_i}$가 되는 유일한 절단
$s \in \mathcal{F}(U)$가 존재한다.
\end{itemize}
\item {\it $\mathcal{B}$ 위의 집합의 층의 사상}은 단순히 집합의
준층의 사상이다.
\end{enumerate}
\end{definition}

\noindent
먼저 층 조건 $(**)$을 덮개들의 공종계에서 확인하는 것으로 충분함을
설명하자. 정의의 상황에서 $U \in \mathcal{B}$라 하자. 잠시
$\text{Cov}_\mathcal{B}(U)$를 $\mathcal{B}$의 원소들에 의한 $U$의
모든 덮개들의 집합이라 쓰자. $\text{Cov}_\mathcal{B}(U)$는 세분에 의해
전순서가 주어짐에 유의하라. 부분집합
$C \subset \text{Cov}_\mathcal{B}(U)$가 공종계라는 것은, 모든
$\mathcal{U} \in \text{Cov}_\mathcal{B}(U)$에 대하여
$\mathcal{U}$를 세분하는 덮개 $\mathcal{V} \in C$가 존재한다는 뜻이다.

\begin{lemma}
\label{sheaves-lemma-cofinal-systems-coverings}
위의 표기를 사용하자. 각 $U \in \mathcal{B}$에 대하여
$C(U) \subset \text{Cov}_\mathcal{B}(U)$를 공종계라 하자. 각
$U \in \mathcal{B}$와 $C(U)$ 안의 각
$\mathcal{U} : U = \bigcup U_i$에 대하여 덮개
$\mathcal{U}_{ij} : U_i \cap U_j = \bigcup U_{ijk}$,
$U_{ijk} \in \mathcal{B}$가 주어졌다고 하자. $\mathcal{F}$를
$\mathcal{B}$ 위의 집합의 준층이라 하자. 다음은 서로 동치이다.
\begin{enumerate}
\item 준층 $\mathcal{F}$는 $\mathcal{B}$ 위의 층이다.
\item 모든 $U \in \mathcal{B}$와 $C(U)$ 안의 모든 덮개
$\mathcal{U} : U = \bigcup U_i$에 대하여 층 조건 $(**)$이 성립한다
(주어진 덮개들 $\mathcal{U}_{ij}$에 대하여).
\end{enumerate}
\end{lemma}

\begin{proof}
(2)가 (1)을 함의함을 보이면 된다. $U \in \mathcal{B}$라 하고
$\mathcal{U} : U = \bigcup_{i\in I} U_i$를 $\mathcal{B}$의 원소들에
의한 임의의 덮개라 하자. 계 $C(U)$가 공종이므로
$\mathcal{U}$를 세분하는 $C(U)$의 원소
$\mathcal{V} : U = \bigcup_{j \in J} V_j$를 찾을 수 있다. 이는
$V_j \subset U_{\alpha(j)}$가 되게 하는 사상 $\alpha : J \to I$가
존재한다는 뜻이다.

\medskip\noindent
$s, s' \in \mathcal{F}(U)$가 $s|_{U_i} = s'|_{U_i}$인 절단들이면
$$
s|_{V_j}
= (s|_{U_{\alpha(j)}})|_{V_j}
= (s'|_{U_{\alpha(j)}})|_{V_j}
= s'|_{V_j}
$$
가 모든 $j$에 대하여 성립함에 유의하라. 따라서 덮개
$\mathcal{V}$에 대한 $(**)$의 유일성으로부터 $s = s'$임을 얻는다.
이로써 임의의 덮개 $\mathcal{U}$에 대한 $(**)$의 유일성 부분을 증명했다.

\medskip\noindent
더 나아가 $U_{ii'k} \in \mathcal{B}$에 의한
$U_i \cap U_{i'} = \bigcup_{k \in I_{ii'}} U_{ii'k}$가 임의의 덮개들이라
하자. 계 $(\mathcal{U}, \mathcal{U}_{ij})$에 대한 $(**)$의 존재성 부분을
증명해 보자. 즉 $s_i \in \mathcal{F}(U_i)$라 하고
$$
s_i|_{U_{ii'k}} = s_{i'}|_{U_{ii'k}}
$$
가 모든 $i, i', k$에 대해 성립한다고 하자. 위의 $\mathcal{V}$와
$\alpha$에 대하여 $t_j = s_{\alpha(j)}|_{V_j}$라 놓자.

\medskip\noindent
여기서 논증에는 작은 난점이 하나 있다. 보조정리의 명제로 주어진 덮개를
$\mathcal{V}_{jj'} : V_j \cap V_{j'} = \bigcup_{l \in J_{jj'}} V_{jj'l}$라
하자. 모든 $j, j', l$에 대하여
$$
t_j|_{V_{jj'l}} = t_{j'}|_{V_{jj'l}}
$$
임은 선험적으로 분명하지 않다. 이를 보려면, 원소들의 족 $s_i$에 대한
가정으로부터 모든 $k \in I_{\alpha(j)\alpha(j')}$에 대해
$$
t_j|_W = t_{j'}|_W \text{ for all } W \in \mathcal{B},
W \subset V_{jj'l} \cap U_{\alpha(j)\alpha(j')k}
$$
임에 유의하라. 또한
$V_j \cap V_{j'} \subset U_{\alpha(j)} \cap U_{\alpha(j')}$이므로
$t_j|_{V_{jj'l}}$와 $t_{j'}|_{V_{jj'l}}$는 $\mathcal{B}$의 원소들에 의한
$V_{jj'l}$의 한 덮개의 각 원소 위에서 일치한다. 따라서 위에서 증명한
유일성 부분에 의해 마침내 $t_j|_{V_{jj'l}}$와
$t_{j'}|_{V_{jj'l}}$가 같다는 원하는 결론을 얻는다. 이제
$(\mathcal{V}, \mathcal{V}_{jj'})$에 대한 성질 $(**)$에 의해 원소
$t \in \mathcal{F}(U)$가 존재한다.

\medskip\noindent
다시 작은 난점이 있다. $t$가 $V_j$에서 $t_j$로 제한됨은 알지만,
$t$가 $U_i$에서 $s_i$로 제한됨은 아직 모른다. 이를 결론내기 위해 집합들
$U_i \cap V_j$, $j \in J$가 $U_i$를 덮는다는 것에 유의하라. 따라서
$U_{i \alpha(j) k} \cap V_j$, $j\in J$, $k \in I_{i\alpha(j)}$도
$U_i$를 덮는다. 이 열린집합들 중 하나에 포함되는 임의의
$W \in \mathcal{B}$ 위에서 $t$와 $s_i$가 $\mathcal{F}$의 같은
절단으로 제한됨을 확인하는 것은 독자에게 맡긴다. 그 열린집합은
$U_{i \alpha(j) k} \cap V_j$, $j\in J$, $k \in I_{i\alpha(j)}$ 중
하나이다. 따라서 위에서 본 유일성 부분에 의해 원하는 결론을 얻는다.
\end{proof}

\begin{lemma}
\label{sheaves-lemma-cofinal-systems-coverings-standard-case}
$X$를 위상공간이라 하자. $\mathcal{B}$를 $X$의 위상의 기저라 하자.
$U' \subset U$ 및 $U'' \subset U$인 모든 세 원소
$U, U', U'' \in \mathcal{B}$에 대하여
$U' \cap U'' \in \mathcal{B}$라고 가정하자. 각
$U \in \mathcal{B}$에 대하여
$C(U) \subset \text{Cov}_\mathcal{B}(U)$를 공종계라 하자.
$\mathcal{F}$를 $\mathcal{B}$ 위의 집합의 준층이라 하자. 다음은
서로 동치이다.
\begin{enumerate}
\item 준층 $\mathcal{F}$는 $\mathcal{B}$ 위의 층이다.
\item 모든 $U \in \mathcal{B}$와 $C(U)$ 안의 모든 덮개
$\mathcal{U} : U = \bigcup U_i$, 그리고
$s_i|_{U_i \cap U_j} = s_j|_{U_i \cap U_j}$인 모든 절단들의 족
$s_i \in \mathcal{F}(U_i)$에 대하여, $U_i$에서 $s_i$로 제한되는
유일한 절단 $s \in \mathcal{F}(U)$가 존재한다.
\end{enumerate}
\end{lemma}

\begin{proof}
이는 각 덮개 $\mathcal{U}_{ij}$가 하나의 원소로 이루어진 특별한 경우에
대한 위 보조정리 \ref{sheaves-lemma-cofinal-systems-coverings}의 재진술이다.
그러나 이 경우는 훨씬 쉬우며 직접 푸는 것도 간단한 연습문제이다.
\end{proof}

\begin{lemma}
\label{sheaves-lemma-condition-star-sections}
$X$를 위상공간이라 하자. $\mathcal{B}$를 $X$의 위상의 기저라 하자.
$U \in \mathcal{B}$라 하자. $\mathcal{F}$를 $\mathcal{B}$ 위의 집합의
층이라 하자. 사상
$$
\mathcal{F}(U) \to \prod\nolimits_{x \in U} \mathcal{F}_x
$$
은 $\mathcal{F}(U)$를 다음 성질을 가진 원소들
$(s_x)_{x\in U}$와 동일시한다.
\begin{itemize}
\item[(*)] 임의의 $x \in U$에 대하여 $x \in V \subset U$인
$V \in \mathcal{B}$와 절단 $\sigma \in \mathcal{F}(V)$가 존재하여,
모든 $y \in V$에 대해 $\mathcal{F}_y$ 안에서
$s_y = (V, \sigma)$이다.
\end{itemize}
\end{lemma}

\begin{proof}
먼저 사상 $\mathcal{F}(U) \to \prod\nolimits_{x \in U} \mathcal{F}_x$는
정의 \ref{sheaves-definition-sheaf-basis}의 층 조건의 유일성에 의해 단사이다.
오른쪽에서 $(*)$를 만족하는 임의의 원소 $(s_x)$를 택하자. 이는 분명히
$U = \bigcup U_i$, $U_i \in \mathcal{B}$인 덮개를 찾아서
$(s_x)_{x \in U_i}$가 어떤 $\sigma_i \in \mathcal{F}(U_i)$에서 오게 할
수 있다는 뜻이다. 모든 $y \in U_i \cap U_j$에 대하여 절단
$\sigma_i$와 $\sigma_j$는 줄기 $\mathcal{F}_y$에서 일치한다. 따라서
$\sigma_i|_{V_{ijy}} = \sigma_j|_{V_{ijy}}$가 되는
$y \in V_{ijy}$인 원소 $V_{ijy} \in \mathcal{B}$가 존재한다. 그러므로
정의 \ref{sheaves-definition-sheaf-basis}의 층 조건 $(**)$을 $\sigma_i$들의
계에 적용하여 원하는 성질을 가진 절단 $s \in \mathcal{F}(U)$를 얻는다.
\end{proof}

\noindent
$X$를 위상공간이라 하자. $\mathcal{B}$를 $X$의 위상의 기저라 하자.
$X$ 위의 집합의 층들의 범주에서 $\mathcal{B}$ 위의 집합의 층들의
범주로 가는 자연스러운 제한 함자가 있다. 이 함자는 사실 범주의
동치이다. 구체적으로 이는 다음을 뜻한다.

\begin{lemma}
\label{sheaves-lemma-extend-off-basis}
$X$를 위상공간이라 하자. $\mathcal{B}$를 $X$의 위상의 기저라 하자.
$\mathcal{F}$를 $\mathcal{B}$ 위의 집합의 층이라 하자. 모든
$U \in \mathcal{B}$에 대하여 제한 사상들과 가환하게
$\mathcal{F}^{ext}(U) = \mathcal{F}(U)$인 $X$ 위의 유일한 집합의 층
$\mathcal{F}^{ext}$가 존재한다.
\end{lemma}

\begin{proof}
먼저 원하는 성질을 가진 준층 $\mathcal{F}^{ext}$를 구성하자. 즉 임의의
열린집합 $U \subset X$에 대하여 $\mathcal{F}^{ext}(U)$를 보조정리
\ref{sheaves-lemma-condition-star-sections}의 $(*)$를 만족하는 원소들
$(s_x)_{x \in U}$의 집합으로 정의한다. 명백한 제한 사상들에 의해
$\mathcal{F}^{ext}$는 집합의 준층이 된다. 또한 보조정리
\ref{sheaves-lemma-condition-star-sections}에 의해 $U$가 기저 $\mathcal{B}$의
원소일 때마다 $\mathcal{F}(U) = \mathcal{F}^{ext}(U)$임을 알 수 있다.
$\mathcal{F}^{ext}$가 층임을 보이려면 보조정리
\ref{sheaves-lemma-sheafification-sheaf}의 증명과 같이 논증하면 된다.
\end{proof}

\noindent
보조정리의 상황에서
$$
\mathcal{F}_x = \mathcal{F}_x^{ext}
$$
임에 유의하라. 이는 $x$를 포함하는 $\mathcal{B}$의 원소들의 모임이
$x$의 열린 근방의 기본계를 이루기 때문이다.

\begin{lemma}
\label{sheaves-lemma-restrict-basis-equivalence}
$X$를 위상공간이라 하자. $\mathcal{B}$를 $X$의 위상의 기저라 하자.
% P01-ERRATA: P01-E1322 -- parallel frozen naming clauses omit "by" after Denote.
$\Sh(\mathcal{B})$로 $\mathcal{B}$ 위의 층들의 범주를 나타내자.
다음 범주의 동치가 있다.
$$
\Sh(X) \longrightarrow \Sh(\mathcal{B})
$$
이 동치는 $X$ 위의 층에 그 층을 $\mathcal{B}$의 원소들로 제한한 것을
대응시킨다.
\end{lemma}

\begin{proof}
% P01-ERRATA: P01-E0051 -- frozen source says "inverse functor in given" at three parallel loci.
역함자는 위의 보조정리 \ref{sheaves-lemma-extend-off-basis}에서 주어진다.
명백한 함자성의 확인은 독자에게 맡긴다.
\end{proof}

\noindent
이로써 기저 $\mathcal{B}$ 위의 집합의 층에 대한 논의를 마친다.
$(\mathcal{C}, F)$를 대수 구조의 한 유형이라 하자. 이 절의 끝에서는
위 구성들이 $\mathcal{C}$에 값을 갖는 층에도 적용됨을 지적하고자 한다.
관련된 개념들을 간단히 정의하자.

\begin{definition}
\label{sheaves-definition-sheaf-structures-basis}
$X$를 위상공간이라 하자. $\mathcal{B}$를 $X$의 위상의 기저라 하자.
$(\mathcal{C}, F)$를 대수 구조의 한 유형이라 하자.
\begin{enumerate}
\item {\it $\mathcal{B}$ 위에서 $\mathcal{C}$에 값을 갖는 준층
$\mathcal{F}$}란 각 $U \in \mathcal{B}$에 $\mathcal{C}$의 대상
$\mathcal{F}(U)$를 대응시키고, $\mathcal{B}$의 원소들의 각 포함관계
$V \subset U$에 $\mathcal{C}$의 사상
$\rho^U_V : \mathcal{F}(U) \to \mathcal{F}(V)$를 대응시키는 규칙으로서,
모든 $U \in \mathcal{B}$에 대해
$\rho^U_U = \text{id}_{\mathcal{F}(U)}$이고 $\mathcal{B}$에서
$W \subset V \subset U$일 때마다
$\rho^U_W = \rho^V_W \circ \rho ^U_V$인 것을 말한다.
\item {\it $\mathcal{B}$ 위에서 $\mathcal{C}$에 값을 갖는 준층의 사상
$\varphi : \mathcal{F} \to \mathcal{G}$}이란 각 원소
$U \in \mathcal{B}$에 제한 사상들과 가환하는 대수 구조의 사상
% P01-ERRATA: P01-E0049 -- frozen source writes each component as phi rather than phi_U.
$\varphi : \mathcal{F}(U) \to \mathcal{G}(U)$를 대응시키는 규칙이다.
\item $\mathcal{B}$ 위에서 $\mathcal{C}$에 값을 갖는 준층
$\mathcal{F}$가 주어졌을 때 $U \mapsto F(\mathcal{F}(U))$를 집합의
바탕 준층이라 한다.
\item {\it $\mathcal{B}$ 위에서 $\mathcal{C}$에 값을 갖는 층
$\mathcal{F}$}란 그 집합의 바탕 준층이 층인, $\mathcal{B}$ 위에서
$\mathcal{C}$에 값을 갖는 준층이다.
\end{enumerate}
\end{definition}

\noindent
이제 $\mathcal{B}$ 위에서 $\mathcal{C}$에 값을 갖는 준층의
$x \in X$에서의 {\it 줄기}를 유향 여극한
$$
\mathcal{F}_x = \colim_{U\in \mathcal{B}, x\in U} \mathcal{F}(U).
$$
으로 정의할 수 있다. 이는 $\mathcal{C}$에 대한 우리의 가정으로 인해
$\mathcal{C}$의 대상으로 존재한다. 또한 $\mathcal{F}_x$의 바탕집합은
$\mathcal{B}$ 위의 집합의 바탕 준층의 줄기임을 알 수 있다.

\medskip\noindent
보조정리들 \ref{sheaves-lemma-cofinal-systems-coverings},
\ref{sheaves-lemma-cofinal-systems-coverings-standard-case} 및
\ref{sheaves-lemma-condition-star-sections}는 연관된 집합의 준층으로 정의한
층 성질을 언급한다. 따라서 이 보조정리들은 $\mathcal{C}$에 값을 갖는
준층이라는 개념으로 변경 없이 일반화된다. 보조정리
\ref{sheaves-lemma-extend-off-basis}의 유사 명제에는 약간의 주의가 필요하다.
% P01-ERRATA: P01-E0050 -- frozen source has singular "analogue" with plural verb "need".
그 명제는 다음과 같다.

\begin{lemma}
\label{sheaves-lemma-extend-off-basis-structures}
$X$를 위상공간이라 하자. $(\mathcal{C}, F)$를 대수 구조의 한 유형이라
하자. $\mathcal{B}$를 $X$의 위상의 기저라 하자. $\mathcal{F}$를
$\mathcal{B}$ 위에서 $\mathcal{C}$에 값을 갖는 층이라 하자.
모든 $U \in \mathcal{B}$에 대하여 제한 사상들과 가환하게
$\mathcal{F}^{ext}(U) = \mathcal{F}(U)$인, $X$ 위에서
$\mathcal{C}$에 값을 갖는 유일한 층 $\mathcal{F}^{ext}$가 존재한다.
\end{lemma}

\begin{proof}
쌍 $(\mathcal{C}, F)$에 부과한 조건에 의해 집합의 바탕 준층의 수준에서
올바르게 작동하는 준층 $\mathcal{F}^{ext}$를 찾는 것으로 충분하다.
따라서 우리의 첫 과제는 모든 열린집합 $U \subset X$에 대하여 적절한
대상 $\mathcal{F}^{ext}(U)$를 구성하는 것이다. 이는 $\mathcal{B}$ 위의
준층의 맥락에서 보조정리 \ref{sheaves-lemma-diagram-fibre-product}를 모방하여
할 수 있다. 그러나 약간 다르지만 본질적으로 동치인 다음 방법을 쓰자.
이를 $U_i \in \mathcal{B}$인 모든 덮개
$\mathcal{U} : U = \bigcup_{i\in I} U_i$에 걸쳐 섬유곱들의 유향 여극한
$$
\mathcal{F}^{ext}(U)
:=
\colim_\mathcal{U} FIB(\mathcal{U})
$$
으로 정의한다. 여기서 섬유곱은
$$
\xymatrix{
FIB(\mathcal{U}) \ar[r] \ar[d] &
\prod\nolimits_{x\in U} \mathcal{F}_x \ar[d] \\
\prod\nolimits_{i\in I} \mathcal{F}(U_i) \ar[r] &
\prod\nolimits_{i \in I} \prod\nolimits_{x\in U_i} \mathcal{F}_x
}
$$
이다. 보조정리 \ref{sheaves-lemma-image-contained-in}과 예
\ref{sheaves-example-application-lemma-image-contained-in}을 보라. 통상적인
논증에 의해, 바탕집합에 대한 이 구성이 위의 $(**)$를 사용한 정의와
같음을 보이면 충분하다. 세부사항은 독자에게 맡긴다.
\end{proof}

\noindent
보조정리의 상황에서 $\mathcal{C}$의 대상들로서
$$
\mathcal{F}_x = \mathcal{F}_x^{ext}
$$
임에 유의하라. 이는 $x$를 포함하는 $\mathcal{B}$의 원소들의 모임이
$x$의 열린 근방의 기본계를 이루기 때문이다.

\begin{lemma}
\label{sheaves-lemma-restrict-basis-equivalence-structures}
$X$를 위상공간이라 하자. $\mathcal{B}$를 $X$의 위상의 기저라 하자.
$(\mathcal{C}, F)$를 대수 구조의 한 유형이라 하자.
% P01-ERRATA: P01-E1322 -- parallel frozen naming clauses omit "by" after Denote.
$\Sh(\mathcal{B}, \mathcal{C})$로 $\mathcal{B}$ 위에서
$\mathcal{C}$에 값을 갖는 층들의 범주를 나타내자. 다음 범주의 동치가 있다.
$$
\Sh(X, \mathcal{C})
\longrightarrow
\Sh(\mathcal{B}, \mathcal{C})
$$
이 동치는 $X$ 위의 층에 그 층을 $\mathcal{B}$의 원소들로 제한한 것을
대응시킨다.
\end{lemma}

\begin{proof}
% P01-ERRATA: P01-E0051 -- frozen source says "inverse functor in given" at three parallel loci.
역함자는 위의 보조정리 \ref{sheaves-lemma-extend-off-basis-structures}에서
주어진다. 명백한 함자성의 확인은 독자에게 맡긴다.
\end{proof}

\noindent
마지막으로 기저 위의 가군의 (준)층의 경우를 다룬다. 기저 위의 집합의
준층들의 범주가 곱을 가지며, 이 곱은 기저의 원소에서의 값들의 곱을
취하는 것으로 기술된다는 쉬운 사실을 사용할 것이다.

\begin{definition}
\label{sheaves-definition-sheaf-modules-basis}
$X$를 위상공간이라 하자. $\mathcal{B}$를 $X$의 위상의 기저라 하자.
$\mathcal{O}$를 $\mathcal{B}$ 위의 환의 준층이라 하자.
\begin{enumerate}
\item {\it $\mathcal{B}$ 위의 $\mathcal{O}$-가군의 준층
$\mathcal{F}$}란 $\mathcal{B}$ 위의 아벨군의 준층과 집합의 준층의 사상
$\mathcal{O} \times \mathcal{F} \to \mathcal{F}$를 합친 것으로서,
모든 $U \in \mathcal{B}$에 대해 사상
$\mathcal{O}(U) \times \mathcal{F}(U) \to \mathcal{F}(U)$가 군
$\mathcal{F}(U)$를 $\mathcal{O}(U)$-가군으로 만드는 것을 말한다.
\item {\it $\mathcal{B}$ 위의 $\mathcal{O}$-가군의 준층의 사상
$\varphi : \mathcal{F} \to \mathcal{G}$}이란 모든
$U \in \mathcal{B}$에 대하여 $\mathcal{O}(U)$-가군 준동형
$\mathcal{F}(U) \to \mathcal{G}(U)$을 유도하는 $\mathcal{B}$ 위의
아벨 준층의 사상이다.
\item $\mathcal{O}$가 $\mathcal{B}$ 위의 환의 층이라고 하자.
{\it $\mathcal{B}$ 위의 $\mathcal{O}$-가군의 층 $\mathcal{F}$}란
그 바탕 아벨군의 준층이 층인 $\mathcal{B}$ 위의
$\mathcal{O}$-가군의 준층이다.
\end{enumerate}
\end{definition}

\noindent
$\mathcal{B}$ 위의 $\mathcal{O}$-가군의 준층의 $x \in X$에서의
{\it 줄기}를 유향 여극한
$$
\mathcal{F}_x = \colim_{U\in \mathcal{B}, x\in U} \mathcal{F}(U).
$$
으로 정의할 수 있다. 이는 $\mathcal{O}_x$-가군이다.

\medskip\noindent
보조정리들 \ref{sheaves-lemma-cofinal-systems-coverings},
\ref{sheaves-lemma-cofinal-systems-coverings-standard-case} 및
\ref{sheaves-lemma-condition-star-sections}는 연관된 집합의 준층으로 정의한
층 성질을 언급한다. 따라서 이 보조정리들은 $\mathcal{O}$-가군의
준층이라는 개념으로 변경 없이 일반화된다. 보조정리
\ref{sheaves-lemma-extend-off-basis}의 유사 명제는 다음과 같다.

\begin{lemma}
\label{sheaves-lemma-extend-off-basis-module}
$X$를 위상공간이라 하자. $\mathcal{B}$를 $X$의 위상의 기저라 하자.
$\mathcal{O}$를 $\mathcal{B}$ 위의 환의 층이라 하자.
$\mathcal{F}$를 $\mathcal{B}$ 위의 $\mathcal{O}$-가군의 층이라 하자.
$\mathcal{O}^{ext}$를 $\mathcal{O}$를 확장하는 $X$ 위의 환의 층이라 하고,
$\mathcal{F}^{ext}$를 $\mathcal{F}$를 확장하는 $X$ 위의 아벨층이라 하자.
보조정리 \ref{sheaves-lemma-extend-off-basis-structures}를 보라. $\mathcal{B}$의
원소들 위에서 주어진 사상과 일치하며 $\mathcal{F}^{ext}$에
$\mathcal{O}^{ext}$-가군 구조를 부여하는 표준 사상
$$
\mathcal{O}^{ext} \times \mathcal{F}^{ext}
\longrightarrow
\mathcal{F}^{ext}
$$
이 존재한다.
\end{lemma}

\begin{proof}
집합의 준층의 수준에서 곱셈 사상을 구성하면 충분하다. 이를 보는 아마
가장 쉬운 방법은 $(f_x)_{x \in U}$, $f_x \in \mathcal{O}_x$와
$(m_x)_{x \in U}$, $m_x \in \mathcal{F}_x$가 $(*)$를 만족하면 원소
$(f_xm_x)_{x \in U}$도 $(*)$를 만족함을 직접 증명하는 것이다.
그러면 원하는 결과를 얻는다. 실제로 보조정리
\ref{sheaves-lemma-extend-off-basis}의 증명에서 확장을 $(*)$를 만족하는 줄기의
원소들의 족을 사용하여 구성하기 때문이다.
\end{proof}

\noindent
보조정리의 상황에서 $\mathcal{O}_x$-가군으로서
$$
\mathcal{F}_x = \mathcal{F}_x^{ext}
$$
임에 유의하라. 이는 $x$를 포함하는 $\mathcal{B}$의 원소들의 모임이
$x$의 열린 근방의 기본계를 이루기 때문이거나, 단순히 바탕집합에 대해
참이기 때문이다.

\begin{lemma}
\label{sheaves-lemma-restrict-basis-equivalence-modules}
$X$를 위상공간이라 하자. $\mathcal{B}$를 $X$의 위상의 기저라 하자.
$\mathcal{O}$를 $X$ 위의 환의 층이라 하자.
% P01-ERRATA: P01-E1322 -- parallel frozen naming clauses omit "by" after Denote.
$\textit{Mod}(\mathcal{O}|_\mathcal{B})$로 $\mathcal{B}$ 위의
$\mathcal{O}|_\mathcal{B}$-가군의 층들의 범주를 나타내자. 다음 범주의
동치가 있다.
$$
\textit{Mod}(\mathcal{O})
\longrightarrow
\textit{Mod}(\mathcal{O}|_\mathcal{B})
$$
이 동치는 $X$ 위의 $\mathcal{O}$-가군의 층에 그 층을
$\mathcal{B}$의 원소들로 제한한 것을 대응시킨다.
\end{lemma}

\begin{proof}
% P01-ERRATA: P01-E0051 -- frozen source says "inverse functor in given" at three parallel loci.
역함자는 위의 보조정리 \ref{sheaves-lemma-extend-off-basis-module}에서 주어진다.
명백한 함자성의 확인은 독자에게 맡긴다.
\end{proof}

\noindent
마지막으로 이 논의와 연속 사상 사이의 관계를 다룬다. 위의 작업 덕분에
이는 이제 매우 쉽다. 먼저 공역에 기저가 주어진 경우를 다룬다.

\begin{lemma}
\label{sheaves-lemma-f-map-basis-below-structures}
$f : X \to Y$를 위상공간들의 연속 사상이라 하자. $(\mathcal{C}, F)$를
대수 구조의 한 유형이라 하자. $\mathcal{F}$를 $X$ 위에서
$\mathcal{C}$에 값을 갖는 층이라 하고, $\mathcal{G}$를 $Y$ 위에서
$\mathcal{C}$에 값을 갖는 층이라 하자. $\mathcal{B}$를 $Y$의 위상의
기저라 하자. 모든 $V \in \mathcal{B}$에 대하여 제한 사상들과 가환하는
$\mathcal{C}$의 사상
$$
\varphi_V :
\mathcal{G}(V)
\longrightarrow
\mathcal{F}(f^{-1}V)
$$
이 주어졌다고 하자. 그러면 $V \in \mathcal{B}$에 대하여
$\varphi_V$를 복원하는 유일한 $f$-사상
(정의 \ref{sheaves-definition-f-map} 및 Section
\ref{sheaves-section-presheaves-structures-functorial}의 $f$-사상에 대한 논의를
보라) $\varphi : \mathcal{G} \to \mathcal{F}$가 존재한다.
\end{lemma}

\begin{proof}
이 사상들의 모임은 $\mathcal{G}$와 $f_*\mathcal{F}$를
$\mathcal{B}$로 제한한 것 사이의 사상을 주는 것과 같으므로 자명하다.
보조정리 \ref{sheaves-lemma-restrict-basis-equivalence-structures}에 의해 이는
$\mathcal{G}$에서 $f_*\mathcal{F}$로 가는 사상을 주는 것과 같고,
보조정리 \ref{sheaves-lemma-f-map}에 의해 다시 $f$-사상을 주는 것과 같다.
대수 구조의 층의 경우를 다루는 방법에 대해서는 보조정리
\ref{sheaves-lemma-f-map-sets-algebraic-structures}와 그 앞의 논의도 보라.
\end{proof}

\noindent
환 달린 공간에 대한 유사 명제는 다음과 같다.

\begin{lemma}
\label{sheaves-lemma-f-map-basis-below-modules}
$(f, f^\sharp) : (X, \mathcal{O}_X) \to (Y, \mathcal{O}_Y)$를 환 달린
공간의 사상이라 하자. $\mathcal{F}$를 $\mathcal{O}_X$-가군의 층이라
하고, $\mathcal{G}$를 $\mathcal{O}_Y$-가군의 층이라 하자.
$\mathcal{B}$를 $Y$의 위상의 기저라 하자. 모든
$V \in \mathcal{B}$에 대하여 제한 사상들과 가환하는
$\mathcal{O}_Y(V)$-가군 사상
$$
\varphi_V :
\mathcal{G}(V)
\longrightarrow
\mathcal{F}(f^{-1}V)
$$
이 주어졌다고 하자(여기서 $\mathcal{F}(f^{-1}V)$는
$f^\sharp_V : \mathcal{O}_Y(V) \to \mathcal{O}_X(f^{-1}V)$를 사용한
가군 구조를 갖는다). 그러면 $V \in \mathcal{B}$에 대해
$\varphi_V$를 복원하는 유일한 $f$-사상
(Section \ref{sheaves-section-ringed-spaces-functoriality-modules}의
$f$-사상에 대한 논의를 보라)
$\varphi : \mathcal{G} \to \mathcal{F}$가 존재한다.
\end{lemma}

\begin{proof}
위의 대수 구조의 층에 대한 대응하는 보조정리의 증명과 같다.
\end{proof}

\begin{lemma}
\label{sheaves-lemma-f-map-basis-above-and-below-structures}
$f : X \to Y$를 위상공간들의 연속 사상이라 하자. $(\mathcal{C}, F)$를
대수 구조의 한 유형이라 하자. $\mathcal{F}$를 $X$ 위에서
$\mathcal{C}$에 값을 갖는 층이라 하고, $\mathcal{G}$를 $Y$ 위에서
$\mathcal{C}$에 값을 갖는 층이라 하자. $\mathcal{B}_Y$를 $Y$의 위상의
기저라 하고, $\mathcal{B}_X$를 $X$의 위상의 기저라 하자. 모든
$V \in \mathcal{B}_Y$ 및 $f(U) \subset V$인
$U \in \mathcal{B}_X$에 대하여 제한 사상들과 가환하는
$\mathcal{C}$의 사상
$$
\varphi_V^U :
\mathcal{G}(V)
\longrightarrow
\mathcal{F}(U)
$$
이 주어졌다고 하자. 그러면 위와 같은 모든 쌍 $(U, V)$에 대하여 합성
$$
\mathcal{G}(V) \xrightarrow{\varphi_V}
\mathcal{F}(f^{-1}(V)) \xrightarrow{\text{restr.}}
\mathcal{F}(U)
$$
으로 $\varphi_V^U$를 복원하는 유일한 $f$-사상
(정의 \ref{sheaves-definition-f-map} 및 Section
\ref{sheaves-section-presheaves-structures-functorial}의 $f$-사상에 대한 논의를
보라) $\varphi : \mathcal{G} \to \mathcal{F}$가 존재한다.
\end{lemma}

\begin{proof}
% P01-ERRATA: P01-E0052 -- frozen proof has two subject-verb agreement errors.
우선 집합의 층에 대해 이를 증명하자. $V \subset Y$를 열린집합으로
고정하고 $s \in \mathcal{G}(V)$를 택하자. 원소
$\varphi_V(s) \in \mathcal{F}(f^{-1}V)$를 구성할 것이다. 모든
$x \in f^{-1}V$에 대하여 $x \in U \subset f^{-1}V$인
$U \in \mathcal{B}_X$를 택하여 줄기 $\mathcal{F}_x$ 안의 값
$\varphi(s)_x$를 정의할 수 있다. 구체적으로 $\varphi(s)_x$를
$(U, \varphi_V^U(s))$의 동치류로 놓는다. 변하는 $U$에 대한 사상들
$\varphi_V^U$가 층
$\mathcal{F}$의 제한들과 가환하므로 족
$(\varphi(s)_x)_{x \in f^{-1}V}$는 분명히 조건 $(*)$를 만족한다.
따라서 보조정리 \ref{sheaves-lemma-extend-off-basis}의 증명에 의해
$(\varphi(s)_x)_{x \in f^{-1}V}$는
$\mathcal{F}(f^{-1}V)$의 유일한 원소 $\varphi_V(s)$에 대응한다.
이로써 집합 사상
$\varphi_V : \mathcal{G}(V) \to \mathcal{F}(f^{-1}V)$를 정의했다.
$\varphi_V$와 $\varphi_V^U$ 사이의 가환성은 보조정리
\ref{sheaves-lemma-condition-star-sections}로부터 따른다.

\medskip\noindent
$V \in \mathcal{B}_Y$를 변화시킬 때 $\varphi_V$의 구성이 제한 사상들과
가환함을 보이는 것은 독자에게 맡긴다. 따라서 위의 보조정리
\ref{sheaves-lemma-f-map-basis-below-structures}를 적용하여 이들을 원하는
$f$-사상으로 ‘붙일’ 수 있다.

\medskip\noindent
마지막으로 이렇게 구성한 집합의 층의 사상이 줄기 위에서 주는 사상
$$
\mathcal{G}_{f(x)} \longrightarrow \mathcal{F}_x
$$
은 $V \in \mathcal{B}_Y$가 $f(x)$를 포함하는 원소들 사이에서 변하고
$U \in \mathcal{B}_X$가 $x$를 포함하는 원소들 사이에서 변할 때의
사상들 $\varphi_V^U$의 여극한이라는 성질을 만족함에 유의하라. 특히
$\mathcal{G}$와 $\mathcal{F}$가 대수 구조의 층들의 집합의 바탕 층이면,
% P01-ERRATA: P01-E0052 -- frozen source says plural "maps" with singular "is" here.
줄기 위의 사상들이 대수 구조의 사상임을 알 수 있다. 따라서 연관된
집합의 바탕 층의 사상 $f^{-1}\mathcal{G} \to \mathcal{F}$는 보조정리
\ref{sheaves-lemma-f-map-sets-algebraic-structures}의 가정을 만족한다.
그러므로 $f^{-1}\mathcal{G} \to \mathcal{F}$는 $\mathcal{C}$에 값을
갖는 층들의 사상이다. 수반성에 의해 이는 $\varphi$가 대수 구조의
층들의 $f$-사상이라는 뜻이다.
\end{proof}

\begin{lemma}
\label{sheaves-lemma-f-map-basis-above-and-below-modules}
$(f, f^\sharp) : (X, \mathcal{O}_X) \to (Y, \mathcal{O}_Y)$를 환 달린
공간의 사상이라 하자. $\mathcal{F}$를 $\mathcal{O}_X$-가군의 층이라
하고, $\mathcal{G}$를 $\mathcal{O}_Y$-가군의 층이라 하자.
$\mathcal{B}_Y$를 $Y$의 위상의 기저라 하고, $\mathcal{B}_X$를
$X$의 위상의 기저라 하자. 모든 $V \in \mathcal{B}_Y$ 및
$f(U) \subset V$인 $U \in \mathcal{B}_X$에 대하여 제한 사상들과
가환하는 $\mathcal{O}_Y(V)$-가군 사상
$$
\varphi_V^U :
\mathcal{G}(V)
\longrightarrow
\mathcal{F}(U)
$$
이 주어졌다고 하자. 여기서 $\mathcal{F}(U)$ 위의
$\mathcal{O}_Y(V)$-가군 구조는 $\mathcal{O}_X(U)$-가군 구조와 사상
$f^\sharp_V : \mathcal{O}_Y(V)
\to \mathcal{O}_X(f^{-1}V) \to \mathcal{O}_X(U)$로부터 온다.
그러면 위와 같은 모든 쌍 $(U, V)$에 대하여 합성
$$
\mathcal{G}(V) \xrightarrow{\varphi_V}
\mathcal{F}(f^{-1}(V)) \xrightarrow{\text{restr.}}
\mathcal{F}(U)
$$
으로 $\varphi_V^U$를 복원하는 유일한 가군의 층의 $f$-사상
(정의 \ref{sheaves-definition-f-map} 및 Section
\ref{sheaves-section-ringed-spaces-functoriality-modules}의 $f$-사상에 대한
논의를 보라) $\varphi : \mathcal{G} \to \mathcal{F}$가 존재한다.
\end{lemma}

\begin{proof}
위와 비슷하므로 생략한다.
\end{proof}

\section{열린 몰입과 (준)층}
\label{sheaves-section-open-immersions}

\noindent
$X$를 위상공간이라 하자. $j : U \to X$를 열린부분집합 $U$의 $X$로의
포함사상이라 하자. Section \ref{sheaves-section-presheaves-functorial}에서
$j_*$가 $j^{-1}$의 오른쪽 수반이 되도록 함자 $j_*$와 $j^{-1}$을
정의했다. 열린 몰입의 경우에는 $j^{-1}$에 왼쪽 수반도 있음이
밝혀지며, 이를 $j_!$로 나타낼 것이다. 먼저 열린 몰입의 경우
$j^{-1}$이 특히 간단하게 기술됨을 지적하자.

\begin{lemma}
\label{sheaves-lemma-j-pullback}
$X$를 위상공간이라 하자. $j : U \to X$를 열린부분집합 $U$의 $X$로의
포함사상이라 하자.
\begin{enumerate}
\item $\mathcal{G}$를 $X$ 위의 집합의 준층이라 하자. 준층
$j_p\mathcal{G}$ (Section \ref{sheaves-section-presheaves-functorial} 참조)는
열린집합 $V \subset U$에 대해 $V \mapsto \mathcal{G}(V)$라는 규칙으로
주어진다.
\item $\mathcal{G}$를 $X$ 위의 집합의 층이라 하자. 층
$j^{-1}\mathcal{G}$는 열린집합 $V \subset U$에 대해
$V \mapsto \mathcal{G}(V)$라는 규칙으로 주어진다.
\item 임의의 점 $u \in U$와 $X$ 위의 임의의 층 $\mathcal{G}$에 대해
줄기의 표준적 동일시
$$
j^{-1}\mathcal{G}_u = (\mathcal{G}|_U)_u = \mathcal{G}_u.
$$
가 성립한다.
\item $U$ 위의 준층의 범주에서 $j_pj_* = \text{id}$이다.
\item $U$ 위의 층의 범주에서 $j^{-1}j_* = \text{id}$이다.
\end{enumerate}
아벨 군의 (준)층, 대수 구조의 (준)층, 그리고 가군의 (준)층에도 같은
기술이 성립한다.
\end{lemma}

\begin{proof}
% P01-ERRATA: P01-E1323 -- frozen source is missing an article before “collection”.
$j_p\mathcal{G}(V)$의 정의에 나오는 여극한은 $V \subset W$를 만족하는
모든 열린집합 $W \subset X$의 모임을 역포함관계로 순서 지은 것
위에서 취한다. 따라서 이 모임에는 최대 원소, 곧 $V$가 있다. 이것이 (1)을
증명한다. $\mathcal{G}$가 층이면 열린집합 $V \subset U$에
$V \mapsto \mathcal{G}(V)$를 대응시키는 것이 분명히 층이므로 (2)도
따른다. $U$에 포함되고 $u$를 포함하는 열린 근방들의 모임은 $X$에서
$u$의 모든 열린 근방들의 모임 안에서 공종이므로 (3)은 (2)로부터
따른다. (4)와 (5)는
$j^{-1}j_*\mathcal{F}(V) = j_*\mathcal{F}(V) = \mathcal{F}(V)$를
계산하면 따른다.

\medskip\noindent
완전히 같은 논증이 아벨 군의 (준)층과 대수 구조의 (준)층에도
성립한다.
\end{proof}

\begin{definition}
\label{sheaves-definition-restriction}
$X$를 위상공간이라 하자. $j : U \to X$를 열린부분집합의
포함사상이라 하자.
\begin{enumerate}
\item $\mathcal{G}$를 $X$ 위의 집합, 아벨 군 또는 대수 구조의
준층이라 하자. 보조정리 \ref{sheaves-lemma-j-pullback}에서 기술한 준층
$j_p\mathcal{G}$를 {\it $\mathcal{G}$의 $U$로의 제한}이라 하고
$\mathcal{G}|_U$로 나타낸다.
\item $\mathcal{G}$를 $X$ 위의 집합의 층, 또는 $X$ 위의 아벨 군이나
대수 구조의 층이라 하자. 층 $j^{-1}\mathcal{G}$를
{\it $\mathcal{G}$의 $U$로의 제한}이라 하고 $\mathcal{G}|_U$로
나타낸다.
\item $(X, \mathcal{O})$가 환 달린 공간이면 쌍
$(U, \mathcal{O}|_U)$를 {\it $U$에 대응하는 $(X, \mathcal{O})$의
열린부분공간}이라 한다.
\item $\mathcal{G}$가 $\mathcal{O}$-가군의 준층이면
$\mathcal{G}|_U$에 곱셈 사상
$\mathcal{O}|_U \times \mathcal{G}|_U \to \mathcal{G}|_U$
(보조정리 \ref{sheaves-lemma-pullback-module} 참조)을 함께 갖춘 것을
{\it $\mathcal{G}$의 $U$로의 제한}이라 한다.
\end{enumerate}
\end{definition}

\noindent
% P01-ERRATA: P01-E0053 -- frozen prose leaves to the reader what item (4) just defined.
가군의 준층의 제한에 대한 정의는 독자에게 맡긴다. 이제 이 절에서는
제한 함자의 왼쪽 수반을 논의할 것이다. 먼저 집합의 (준)층의 경우에
그 정의를 주자.

\begin{definition}
\label{sheaves-definition-j-shriek}
$X$를 위상공간이라 하자. $j : U \to X$를 열린부분집합의
포함사상이라 하자.
\begin{enumerate}
\item $\mathcal{F}$를 $U$ 위의 집합의 준층이라 하자.
{\it $\mathcal{F}$의 공집합에 의한 확장 $j_{p!}\mathcal{F}$}을
다음 규칙으로 정의되는 $X$ 위의 집합의 준층으로 정의한다.
$$
j_{p!}\mathcal{F}(V) =
\left\{
\begin{matrix}
\emptyset & \text{if} & V \not \subset U \\
\mathcal{F}(V) & \text{if} & V \subset U
\end{matrix}
\right.
$$
제한 사상은 자명한 것으로 취한다.
\item $\mathcal{F}$를 $U$ 위의 집합의 층이라 하자.
{\it $\mathcal{F}$의 공집합에 의한 확장 $j_!\mathcal{F}$}을 준층
$j_{p!}\mathcal{F}$의 층화로 정의한다.
\end{enumerate}
\end{definition}

\begin{lemma}
\label{sheaves-lemma-j-shriek}
$X$를 위상공간이라 하자. $j : U \to X$를 열린부분집합의
포함사상이라 하자.
\begin{enumerate}
\item 함자 $j_{p!}$는 제한 함자 $j_p$의 왼쪽 수반이다
(보조정리 \ref{sheaves-lemma-j-pullback} 참조).
\item 함자 $j_!$는 제한의 왼쪽 수반이다. 식으로 쓰면
$$
\Mor_{\Sh(X)}(j_!\mathcal{F}, \mathcal{G})
=
\Mor_{\Sh(U)}(\mathcal{F}, j^{-1}\mathcal{G})
=
\Mor_{\Sh(U)}(\mathcal{F}, \mathcal{G}|_U)
$$
이며, 이는 $\mathcal{F}$와 $\mathcal{G}$에 대해 쌍함자적이다.
\item $\mathcal{F}$를 $U$ 위의 집합의 층이라 하자. 층
$j_!\mathcal{F}$의 줄기는 다음과 같이 기술된다.
$$
j_{!}\mathcal{F}_x =
\left\{
\begin{matrix}
\emptyset & \text{if} & x \not \in U \\
\mathcal{F}_x & \text{if} & x \in U
\end{matrix}
\right.
$$
\item $U$ 위의 준층의 범주에서 $j_pj_{p!} = \text{id}$이다.
\item $U$ 위의 층의 범주에서 $j^{-1}j_! = \text{id}$이다.
\end{enumerate}
\end{lemma}

\begin{proof}
$j_{p!}\mathcal{F}$에서 $\mathcal{G}$로 사상하려면 모든
$V \subset U$에 대해 제한 사상들과 가환하도록
$\mathcal{F}(V) \to \mathcal{G}(V)$를 사상하는 것으로 충분하다.
보조정리 \ref{sheaves-lemma-j-pullback}에 의해
$\mathcal{F} \to \mathcal{G}|_U$인 사상에도 같은 기술이 성립한다.
$j_!$와 제한 사이의 수반성은 이 사실과 층화의 성질들로부터 따른다.
줄기의 동일시는 공집합에 의한 확장의 정의와 줄기의 정의로부터
자명하다. 명제 (4)와 (5)는 $U$의 임의의 열린집합에서 층의 값을
계산하면 따른다.
\end{proof}

\noindent
$\mathcal{F}$가 $U$ 위의 아벨 군의 층이면 위에서 정의한
$j_!\mathcal{F}$는 일반적으로 아벨 군의 층이 아님에 유의하라.
예를 들어 그 줄기 가운데 일부가 공집합인데, 공집합은 결코 아벨 군이
아니다. 따라서 고려하는 층의 종류에 따라 $j_!$의 정의를 수정해야
한다. 공집합에 의한 확장의 정의에서 공집합을 택한 까닭은 이것이
집합의 범주에서 초기 대상이기 때문이다. 따라서 아벨 군의 경우에는
$0$을 사용한다(더 일반적으로는 임의의 아벨 범주에 값을 갖는 층에도
마찬가지이다).

\begin{definition}
\label{sheaves-definition-j-shriek-structures}
$X$를 위상공간이라 하자. $j : U \to X$를 열린부분집합의
포함사상이라 하자.
\begin{enumerate}
\item $\mathcal{F}$를 $U$ 위의 아벨 준층이라 하자.
{\it $\mathcal{F}$의 $0$에 의한 확장
$j_{p!}\mathcal{F}$}을 다음 규칙으로 정의되는 $X$ 위의 아벨
준층으로 정의한다.
$$
j_{p!}\mathcal{F}(V) =
\left\{
\begin{matrix}
0 & \text{if} & V \not \subset U \\
\mathcal{F}(V) & \text{if} & V \subset U
\end{matrix}
\right.
$$
제한 사상은 자명한 것으로 취한다.
\item $\mathcal{F}$를 $U$ 위의 아벨 층이라 하자.
{\it $\mathcal{F}$의 $0$에 의한 확장 $j_!\mathcal{F}$}을 아벨 준층
$j_{p!}\mathcal{F}$의 층화로 정의한다.
\item $\mathcal{C}$를 초기 대상 $e$를 갖는 범주라 하자.
$\mathcal{F}$를 $\mathcal{C}$에 값을 갖는 $U$ 위의 준층이라 하자.
{\it $\mathcal{F}$의 $e$에 의한 확장
$j_{p!}\mathcal{F}$}을 다음 규칙으로 정의되는, $\mathcal{C}$에 값을
갖는 $X$ 위의 준층으로 정의한다.
$$
j_{p!}\mathcal{F}(V) =
\left\{
\begin{matrix}
e & \text{if} & V \not \subset U \\
\mathcal{F}(V) & \text{if} & V \subset U
\end{matrix}
\right.
$$
제한 사상은 자명한 것으로 취한다.
\item $(\mathcal{C}, F)$를 $\mathcal{C}$가 초기 대상 $e$를 갖는
대수 구조의 종류라 하자. $\mathcal{F}$를 $U$ 위의 (주어진 종류의)
대수 구조의 층이라 하자.
% P01-ERRATA: P01-E1324 -- frozen source says “give type” for “given type”.
{\it $\mathcal{F}$의 $e$에 의한 확장 $j_!\mathcal{F}$}을 위에서
정의한 준층 $j_{p!}\mathcal{F}$의 층화로 정의한다.
\item $\mathcal{O}$를 $X$ 위의 환의 준층이라 하자.
$\mathcal{F}$를 $\mathcal{O}|_U$-가군의 준층이라 하자. 이 경우
{\it $0$에 의한 확장}을, 아벨 준층으로서 $j_{p!}\mathcal{F}$와 같고
곱셈 사상
$\mathcal{O} \times j_{p!}\mathcal{F} \to j_{p!}\mathcal{F}$를 갖춘
$\mathcal{O}$-가군의 준층으로 정의한다.
\item $\mathcal{O}$를 $X$ 위의 환의 층이라 하자.
$\mathcal{F}$를 $\mathcal{O}|_U$-가군의 층이라 하자. 이 경우
{\it $0$에 의한 확장}을, 아벨 층으로서 $j_!\mathcal{F}$와 같고
곱셈 사상 $\mathcal{O} \times j_!\mathcal{F} \to j_!\mathcal{F}$를
갖춘 $\mathcal{O}$-가군으로 정의한다.
\end{enumerate}
\end{definition}

\noindent
대수 구조의 층의 맥락에서도 $j_!$를 정의할 수 있다는 것은
사실이다(아래를 보라). 그러나 그 결과로 얻는 대상은 해당 대수
구조의 종류에 의존한다. 예를 들어 $\mathcal{O}$가 $U$ 위의 환의
층이면
$j_{!, rings}\mathcal{O} \not = j_{!, abelian}\mathcal{O}$이다.
환의 범주에서 초기 대상은 $\mathbf{Z}$이고 아벨 군의 범주에서 초기
대상은 $0$이기 때문이다. 특히 지금까지 본 것과는 달리 함자 $j_!$는
{\it 바탕 집합의 층을 취하는 연산과 가환하지 않는다}! 평소처럼 아벨
군의 (준)층, 대수 구조의 (준)층, 그리고 가군의 (준)층의 경우를
따로 다룬다.

\begin{lemma}
\label{sheaves-lemma-j-shriek-abelian}
$X$를 위상공간이라 하자. $j : U \to X$를 열린부분집합의
포함사상이라 하자. 아벨 (준)층에 대한 제한 함자와 $0$에 의한 확장
함자를 생각하자.
\begin{enumerate}
\item 함자 $j_{p!}$는 제한 함자 $j_p$의 왼쪽 수반이다
(보조정리 \ref{sheaves-lemma-j-pullback} 참조).
\item 함자 $j_!$는 제한의 왼쪽 수반이다. 식으로 쓰면
$$
\Mor_{\textit{Ab}(X)}(j_!\mathcal{F}, \mathcal{G})
=
\Mor_{\textit{Ab}(U)}(\mathcal{F}, j^{-1}\mathcal{G})
=
\Mor_{\textit{Ab}(U)}(\mathcal{F}, \mathcal{G}|_U)
$$
이며, 이는 $\mathcal{F}$와 $\mathcal{G}$에 대해 쌍함자적이다.
\item $\mathcal{F}$를 $U$ 위의 아벨 층이라 하자. 층
$j_!\mathcal{F}$의 줄기는 다음과 같이 기술된다.
$$
j_{!}\mathcal{F}_x =
\left\{
\begin{matrix}
0 & \text{if} & x \not \in U \\
\mathcal{F}_x & \text{if} & x \in U
\end{matrix}
\right.
$$
\item $U$ 위의 아벨 준층의 범주에서
$j_pj_{p!} = \text{id}$이다.
\item $U$ 위의 아벨 층의 범주에서 $j^{-1}j_! = \text{id}$이다.
\end{enumerate}
\end{lemma}

\begin{proof}
생략한다.
\end{proof}

\begin{lemma}
\label{sheaves-lemma-j-shriek-structures}
$X$를 위상공간이라 하자. $j : U \to X$를 열린부분집합의
포함사상이라 하자. $(\mathcal{C}, F)$를 $\mathcal{C}$가 초기 대상
$e$를 갖는 대수 구조의 종류라 하자. 위에서 정의한 대수 구조의
(준)층에 대한 제한 함자와 $e$에 의한 확장 함자를 생각하자.
\begin{enumerate}
\item 함자 $j_{p!}$는 제한 함자 $j_p$의 왼쪽 수반이다
(보조정리 \ref{sheaves-lemma-j-pullback} 참조).
\item 함자 $j_!$는 제한의 왼쪽 수반이다. 식으로 쓰면
$$
\Mor_{\Sh(X, \mathcal{C})}(j_!\mathcal{F}, \mathcal{G})
=
\Mor_{\Sh(U, \mathcal{C})}(\mathcal{F}, j^{-1}\mathcal{G})
=
\Mor_{\Sh(U, \mathcal{C})}(\mathcal{F}, \mathcal{G}|_U)
$$
이며, 이는 $\mathcal{F}$와 $\mathcal{G}$에 대해 쌍함자적이다.
\item $\mathcal{F}$를 $U$ 위의 층이라 하자. 층
$j_!\mathcal{F}$의 줄기는 다음과 같이 기술된다.
$$
j_{!}\mathcal{F}_x =
\left\{
\begin{matrix}
e & \text{if} & x \not \in U \\
\mathcal{F}_x & \text{if} & x \in U
\end{matrix}
\right.
$$
\item $U$ 위의 대수 구조의 준층의 범주에서
$j_pj_{p!} = \text{id}$이다.
\item $U$ 위의 대수 구조의 층의 범주에서
$j^{-1}j_! = \text{id}$이다.
\end{enumerate}
\end{lemma}

\begin{proof}
생략한다.
\end{proof}

\begin{lemma}
\label{sheaves-lemma-j-shriek-modules}
$(X, \mathcal{O})$를 환 달린 공간이라 하자.
$j : (U, \mathcal{O}|_U) \to (X, \mathcal{O})$를
열린부분공간이라 하자. 위에서 정의한 가군의 (준)층에 대한 제한
함자와 $0$에 의한 확장 함자를 생각하자.
\begin{enumerate}
\item 함자 $j_{p!}$는 제한의 왼쪽 수반이다. 식으로 쓰면
$$
\Mor_{\textit{PMod}(\mathcal{O})}(j_{p!}\mathcal{F}, \mathcal{G})
=
\Mor_{\textit{PMod}(\mathcal{O}|_U)}(\mathcal{F}, \mathcal{G}|_U)
$$
이며, 이는 $\mathcal{F}$와 $\mathcal{G}$에 대해 쌍함자적이다.
\item 함자 $j_!$는 제한의 왼쪽 수반이다. 식으로 쓰면
$$
\Mor_{\textit{Mod}(\mathcal{O})}(j_!\mathcal{F}, \mathcal{G})
=
\Mor_{\textit{Mod}(\mathcal{O}|_U)}(\mathcal{F}, \mathcal{G}|_U)
$$
이며, 이는 $\mathcal{F}$와 $\mathcal{G}$에 대해 쌍함자적이다.
\item $\mathcal{F}$를 $U$ 위의 $\mathcal{O}$-가군의 층이라 하자.
층 $j_!\mathcal{F}$의 줄기는 다음과 같이 기술된다.
$$
j_{!}\mathcal{F}_x =
\left\{
\begin{matrix}
0 & \text{if} & x \not \in U \\
\mathcal{F}_x & \text{if} & x \in U
\end{matrix}
\right.
$$
\item $U$ 위의 $\mathcal{O}|_U$-가군의 층의 범주에서
$j^{-1}j_! = \text{id}$이다.
\end{enumerate}
\end{lemma}

\begin{proof}
생략한다.
\end{proof}

\noindent
위의 보조정리들에 의해 함자 $j_*$와 $j_!$는 모두 $U$ 위의 층의
범주를 $X$ 위의 층의 범주로 보내는 충만충실 매입임에 유의하라.
이 함자의 본질적 상을 쉽게 기술할 수 있는 것은 함자 $j_!$에
대해서뿐이다.

\begin{lemma}
\label{sheaves-lemma-equivalence-categories-open}
$X$를 위상공간이라 하자. $j : U \to X$를 열린부분집합의
포함사상이라 하자. 함자
$$
j_! : \Sh(U) \longrightarrow \Sh(X)
$$
는 충만충실이다. 그 본질적 상은 모든 $x \in X \setminus U$에 대해
$\mathcal{G}_x = \emptyset$인 층 $\mathcal{G}$들로 정확히
이루어진다.
\end{lemma}

\begin{proof}
% P01-ERRATA: P01-E1325 -- frozen source says “Fully faithfulness”.
충만충실성은 $j^{-1} j_! = \text{id}$로부터 형식적으로 따른다.
함자의 상에 속하는 모든 층이 보조정리에 서술된 줄기의 성질을
가짐을 이미 보았다. 역으로 $\mathcal{G}$가 그 성질을 갖는다고 하자.
그러면
$$
j_! j^{-1} \mathcal{G} \to \mathcal{G}
$$
가 모든 줄기 위에서 동형이고 따라서 동형임을 쉽게 확인할 수 있다.
\end{proof}

\begin{lemma}
\label{sheaves-lemma-equivalence-categories-open-abelian}
$X$를 위상공간이라 하자. $j : U \to X$를 열린부분집합의
포함사상이라 하자. 함자
$$
j_! : \textit{Ab}(U) \longrightarrow \textit{Ab}(X)
$$
는 충만충실이다. 그 본질적 상은 모든 $x \in X \setminus U$에 대해
$\mathcal{G}_x = 0$인 층 $\mathcal{G}$들로 정확히 이루어진다.
\end{lemma}

\begin{proof}
생략한다.
\end{proof}

\begin{lemma}
\label{sheaves-lemma-equivalence-categories-open-structures}
$X$를 위상공간이라 하자. $j : U \to X$를 열린부분집합의
포함사상이라 하자. $(\mathcal{C}, F)$를 $\mathcal{C}$가 초기 대상
$e$를 갖는 대수 구조의 종류라 하자. 함자
$$
j_! : \Sh(U, \mathcal{C}) \longrightarrow \Sh(X, \mathcal{C})
$$
는 충만충실이다. 그 본질적 상은 모든 $x \in X \setminus U$에 대해
$\mathcal{G}_x = e$인 층 $\mathcal{G}$들로 정확히 이루어진다.
\end{lemma}

\begin{proof}
생략한다.
\end{proof}


\begin{lemma}
\label{sheaves-lemma-equivalence-categories-open-modules}
$(X, \mathcal{O})$를 환 달린 공간이라 하자.
$j : (U, \mathcal{O}|_U) \to (X, \mathcal{O})$를
열린부분공간이라 하자. 함자
$$
j_! : \textit{Mod}(\mathcal{O}|_U) \longrightarrow \textit{Mod}(\mathcal{O})
$$
는 충만충실이다. 그 본질적 상은 모든 $x \in X \setminus U$에 대해
$\mathcal{G}_x = 0$인 층 $\mathcal{G}$들로 정확히 이루어진다.
\end{lemma}

\begin{proof}
생략한다.
\end{proof}

\begin{remark}
\label{sheaves-remark-j-shriek-not-exact}
$j : U \to X$를 위와 같은 위상공간의 열린 몰입이라 하자.
$x \in X$, $x \not \in U$라 하자. $\mathcal{F}$를 $U$ 위의 집합의
층이라 하자. 보조정리 \ref{sheaves-lemma-j-shriek}에 의해
$j_!\mathcal{F}_x = \emptyset$이다. 따라서 $U = X$가 아닌 한
$j_!$는 $\Sh(U)$의 끝 대상을 $\Sh(X)$의 끝 대상으로 보내지
않는다. Categories, Section
\ref{categories-section-exact-functor}의 규약에 따르면 이는 집합의
층의 범주 사이의 함자로서 $j_!$가 좌완전이 아니라는 뜻이다. 아벨
층 위의 $j_!$는 완전함을 나중에 보일 것이다. Modules, Lemma
\ref{modules-lemma-j-shriek-exact}를 보라.
\end{remark}

\section{닫힌 몰입과 (준)층}
\label{sheaves-section-closed-immersions}

\noindent
$X$를 위상공간이라 하자. $i : Z \to X$를 닫힌부분집합 $Z$의 $X$로의
포함사상이라 하자. Section \ref{sheaves-section-presheaves-functorial}에서
$i_*$가 $i^{-1}$의 오른쪽 수반이 되도록 함자 $i_*$와 $i^{-1}$을
정의했다.

\begin{lemma}
\label{sheaves-lemma-stalks-closed-pushforward}
$X$를 위상공간이라 하자. $i : Z \to X$를 닫힌부분집합 $Z$의 $X$로의
포함사상이라 하자. $\mathcal{F}$를 $Z$ 위의 집합의 층이라 하자.
$i_*\mathcal{F}$의 줄기는 다음과 같이 기술된다.
$$
i_*\mathcal{F}_x =
\left\{
\begin{matrix}
\{*\} & \text{if} & x \not \in Z \\
\mathcal{F}_x & \text{if} & x \in Z
\end{matrix}
\right.
$$
여기서 $\{*\}$는 한원소 집합을 나타낸다. 또한 $Z$ 위의 집합의 층의
범주에서 $i^{-1}i_* = \text{id}$이다. 더 나아가 $Z$ 위의 아벨 층,
각각 $Z$ 위의 대수 구조의 층에도 같은 명제가 성립한다. 이때
$\{*\}$를 각각 $0$, 대수 구조의 범주의 끝 대상으로 바꾸어야 한다.
\end{lemma}

\begin{proof}
$x \not \in Z$이면 $Z$와 만나지 않으면서 임의로 작은 $x$의 열린
근방 $U$들이 존재한다. $\mathcal{F}$가 층이므로 그러한 모든 $U$에
대해 $\mathcal{F}(i^{-1}(U)) = \{*\}$이다. 주석
\ref{sheaves-remark-confusion}을 보라. 이것이 첫째 경우를 증명한다. 둘째
경우는 $z \in Z$에 대해 $z$의 임의의 열린 근방이 $X$의 어떤
열린집합 $U$에 대한 $Z \cap U$ 꼴이라는 사실로부터 나온다.
$i^{-1}i_* = \text{id}$라는 명제를 위해 표준 사상
$i^{-1}i_*\mathcal{F} \to \mathcal{F}$를 생각하자. 이는 줄기 위에서
동형이고(위를 보라), 따라서 동형이다.

\medskip\noindent
아벨 군의 층과 대수 구조의 층의 경우에도 같은 방식으로 논증한다.
\end{proof}


\begin{lemma}
\label{sheaves-lemma-equivalence-categories-closed}
$X$를 위상공간이라 하자. $i : Z \to X$를 닫힌부분집합의
포함사상이라 하자. 함자
$$
i_* : \Sh(Z) \longrightarrow \Sh(X)
$$
는 충만충실이다. 그 본질적 상은 모든 $x \in X \setminus Z$에 대해
$\mathcal{G}_x = \{*\}$인 층 $\mathcal{G}$들로 정확히 이루어진다.
\end{lemma}

\begin{proof}
% P01-ERRATA: P01-E1325 -- frozen source again says “Fully faithfulness”.
충만충실성은 $i^{-1} i_* = \text{id}$로부터 형식적으로 따른다.
함자의 상에 속하는 모든 층이 보조정리에 서술된 줄기의 성질을
가짐을 이미 보았다. 역으로 $\mathcal{G}$가 그 성질을 갖는다고 하자.
그러면
$$
\mathcal{G} \to i_* i^{-1} \mathcal{G}
$$
가 모든 줄기 위에서 동형이고 따라서 동형임을 쉽게 확인할 수 있다.
\end{proof}

\begin{lemma}
\label{sheaves-lemma-equivalence-categories-closed-abelian}
$X$를 위상공간이라 하자. $i : Z \to X$를 닫힌부분집합의
포함사상이라 하자. 함자
$$
i_* : \textit{Ab}(Z) \longrightarrow \textit{Ab}(X)
$$
는 충만충실이다. 그 본질적 상은 모든 $x \in X \setminus Z$에 대해
$\mathcal{G}_x = 0$인 층 $\mathcal{G}$들로 정확히 이루어진다.
\end{lemma}

\begin{proof}
생략한다.
\end{proof}

\begin{lemma}
\label{sheaves-lemma-equivalence-categories-closed-structures}
$X$를 위상공간이라 하자. $i : Z \to X$를 닫힌부분집합의
포함사상이라 하자. $(\mathcal{C}, F)$를 끝 대상 $0$을 갖는 대수
구조의 종류라 하자. 함자
$$
i_* : \Sh(Z, \mathcal{C}) \longrightarrow \Sh(X, \mathcal{C})
$$
는 충만충실이다. 그 본질적 상은 모든 $x \in X \setminus Z$에 대해
$\mathcal{G}_x = 0$인 층 $\mathcal{G}$들로 정확히 이루어진다.
\end{lemma}

\begin{proof}
생략한다.
\end{proof}

\begin{remark}
\label{sheaves-remark-i-star-not-exact}
$i : Z \to X$를 위와 같은 위상공간의 닫힌 몰입이라 하자.
$x \in X$, $x \not \in Z$라 하자. $\mathcal{F}$를 $Z$ 위의 집합의
층이라 하자. 보조정리 \ref{sheaves-lemma-stalks-closed-pushforward}에 의해
$(i_*\mathcal{F})_x = \{ * \}$이다. 따라서 $\mathcal{F} = * \amalg *$이고
$*$가 한원소 층이면
% P01-ERRATA: P01-E0055 -- frozen second stalk expression drops parentheses.
$i_*\mathcal{F}_x = \{*\} \not = i_*(*)_x \amalg i_*(*)_x$이다.
뒤의 집합은 두 원소 집합이기 때문이다. Categories, Section
\ref{categories-section-exact-functor}의 규약에 따르면 이는 집합의
층의 범주 사이의 함자로서 $i_*$가 우완전이 아니라는 뜻이다.
특히 이 함자는 오른쪽 수반을 가질 수 없다. Categories, Lemma
\ref{categories-lemma-exact-adjoint}를 보라.

\medskip\noindent
한편 아벨 층 위의 $i_*$는 완전하고 오른쪽 수반도 가짐을 나중에
보일 것이다(Modules, Lemma
\ref{modules-lemma-i-star-right-adjoint} 참조). 그 오른쪽 수반은
$X$ 위의 아벨 층에 $Z$에서 지지되는 절단들의 층을 대응시키는
함자이다.
\end{remark}

\begin{remark}
\label{sheaves-remark-closed-immersion-spaces}
닫힌 몰입과 환 달린 공간 사이의 관계는 논의하지 않았다. 환 달린
공간의 닫힌 몰입이라는 개념은 준연접층의 맥락에서 논의하는 것이
가장 적절하기 때문이다. Modules, Section
\ref{modules-section-closed-immersion}을 보라.
\end{remark}

\section{층 붙이기}
\label{sheaves-section-glueing-sheaves}

\noindent
이 절에서는 $X$의 덮개의 원소들 위에 정의된 층들을 붙인다. 먼저
사상을 다룬다.

\begin{lemma}
\label{sheaves-lemma-glue-maps}
$X$를 위상공간이라 하자. $X = \bigcup U_i$를 열린 덮개라 하자.
$\mathcal{F}$, $\mathcal{G}$를 $X$ 위의 집합의 층이라 하자. 모든
$i, j \in I$에 대해 사상 $\varphi_i, \varphi_j$의 제한이 같은 사상
$\mathcal{F}|_{U_i \cap U_j} \to \mathcal{G}|_{U_i \cap U_j}$이 되도록
층의 사상들의 족
$$
\varphi_i :
\mathcal{F}|_{U_i}
\longrightarrow
\mathcal{G}|_{U_i}
$$
가 주어졌다고 하자. 그러면 각 $U_i$로의 제한이 $\varphi_i$와
일치하는 유일한 층의 사상
$$
\varphi :
\mathcal{F}
\longrightarrow
\mathcal{G}
$$
가 존재한다.
\end{lemma}

\begin{proof}
각 열린부분집합 $U \subset X$에 대해
$$
\varphi_U : \mathcal{F}(U) \to \mathcal{G}(U), \quad
s \mapsto \varphi_U(s)
$$
를 정의하자. 여기서 $\varphi_U(s)$는
$$
(\varphi_U(s))|_{U \cap U_i}
= (\varphi_i)_{U \cap U_i}(s|_{U \cap U_i}).
$$
를 만족하는 유일한 절단이다. 그러한 절단의 존재성과 유일성은 층
공리로부터 따른다. 실제로
\begin{align*}
((\varphi_i)_{U \cap U_i}(s|_{U \cap U_i}))|_{U \cap U_i \cap U_j}
&= (\varphi_i)_{U \cap U_i \cap U_j}(s|_{U \cap U_i \cap U_j})\\
&= (\varphi_j)_{U \cap U_i \cap U_j}(s|_{U \cap U_i \cap U_j})\\
&= ((\varphi_j)_{U \cap U_j}(s|_{U \cap U_j}))|_{U \cap U_i \cap U_j}.
\end{align*}
이다. 이 사상들의 족은 실제로 층의 사상을 준다. 열린부분집합
$V \subset U \subset X$에 대해
$$
(\varphi_U(s))|_V
= \varphi_V(s|_V)
$$
이다. 각 $i \in I$에 대해 다음이 성립하기 때문이다.
% P01-ERRATA: P01-E0056 -- frozen display uses undefined s_V instead of s|_V.
\begin{align*}
(\varphi_U(s))|_{V \cap U_i}
&= ((\varphi_U(s))|_{U \cap U_i})|_{V \cap U_i}\\
&= ((\varphi_i)_{U \cap U_i}(s|_{U \cap U_i}))|_{V \cap U_i}\\
&= (\varphi_i)_{V \cap U_i}(s|_{V \cap U_i})\\
&= \varphi_V(s_{V})|_{V \cap U_i}.
\end{align*}
또한 각 $U_i$로의 제한은 $\varphi_i$와 일치한다. 실제로 열린부분집합
% P01-ERRATA: P01-E1326 -- frozen source omits an article in “given U ... open subset”.
$U \subset X$와 $s \in \mathcal{F}(U \cap U_i)$가 주어지면
\begin{align*}
\varphi_{U \cap U_i}(s)
&= \varphi_{U \cap U_i}(s)|_{U \cap U_i}\\
&= (\varphi_i)_{U \cap U_i}(s|_{U \cap U_i})\\
&= (\varphi_i)_{U \cap U_i}(s).
\end{align*}
\end{proof}

\noindent
앞의 보조정리는 위상공간 $X$ 위의 두 층 $\mathcal{F}$,
$\mathcal{G}$가 주어졌을 때 규칙
$$
U \longmapsto
\Mor_{\Sh(U)}(
\mathcal{F}|_U, \mathcal{G}|_U)
$$
가 층을 정의함을 뜻한다. 이것은 일종의 {\it 내부 hom 층}이다.
집합의 층의 맥락에서는 거의 쓰이지 않고, 가군의 층의 맥락에서 더
흔히 쓰인다. Modules, Section
\ref{modules-section-internal-hom}을 보라.

\medskip\noindent
$X$를 위상공간이라 하자. $X = \bigcup_{i\in I} U_i$를 열린 덮개라
하자. 각 $i \in I$에 대해 $\mathcal{F}_i$를 $U_i$ 위의 집합의 층이라
하자. 각 쌍 $i, j \in I$에 대해
$$
\varphi_{ij} :
\mathcal{F}_i|_{U_i \cap U_j}
\longrightarrow
\mathcal{F}_j|_{U_i \cap U_j}
$$
를 집합의 층의 동형이라 하자. 또한 모든 세 지표
$i, j, k \in I$에 대해 다음 도표가 가환한다고 가정하자.
$$
\xymatrix{
\mathcal{F}_i|_{U_i \cap U_j \cap U_k}
\ar[rr]_{\varphi_{ik}}
\ar[rd]_{\varphi_{ij}} & &
\mathcal{F}_k|_{U_i \cap U_j \cap U_k} \\
&
\mathcal{F}_j|_{U_i \cap U_j \cap U_k}
\ar[ru]_{\varphi_{jk}}
}
$$
% P01-ERRATA: P01-E1327 -- frozen source repeatedly uses singular article with plural “data”.
이러한 자료의 족 $(\mathcal{F}_i, \varphi_{ij})$를 덮개
$X = \bigcup U_i$에 대한 {\it 집합의 층의 붙이기 자료}라 한다.

\begin{lemma}
\label{sheaves-lemma-glue-sheaves}
$X$를 위상공간이라 하자. $X = \bigcup_{i\in I} U_i$를 열린 덮개라
하자. 덮개 $X = \bigcup U_i$에 대한 집합의 층의 임의의 붙이기 자료
$(\mathcal{F}_i, \varphi_{ij})$가 주어지면, $X$ 위의 집합의 층
$\mathcal{F}$와 동형
$$
\varphi_i : \mathcal{F}|_{U_i} \to \mathcal{F}_i
$$
가 존재하여 도표
$$
\xymatrix{
\mathcal{F}|_{U_i \cap U_j} \ar[r]_{\varphi_i} \ar[d]_{\text{id}} &
\mathcal{F}_i|_{U_i \cap U_j} \ar[d]^{\varphi_{ij}} \\
\mathcal{F}|_{U_i \cap U_j} \ar[r]^{\varphi_j} &
\mathcal{F}_j|_{U_i \cap U_j}
}
$$
가 가환한다.
\end{lemma}

\begin{proof}
첫째 증명. 이 증명에서는 열린집합 $W \subset X$ 위에서
$\mathcal{F}$의 절단들의 집합을 식으로 준다. 즉,
$$
\mathcal{F}(W) =
\{
(s_i)_{i \in I} \mid
s_i \in \mathcal{F}_i(W \cap U_i),
\varphi_{ij}(s_i|_{W \cap U_i \cap U_j}) = s_j|_{W \cap U_i \cap U_j}
\}.
$$
로 정의한다. $W' \subset W$에 대한 제한 사상은 각 $s_i$를
% P01-ERRATA: P01-E1328 -- frozen source says “by the restricting”.
$W' \cap U_i$로 제한하여 정의한다. $\mathcal{F}$의 층 조건은 각
$\mathcal{F}_i$의 층 조건으로부터 즉시 따른다.

\medskip\noindent
이제 $\mathcal{F}|_{U_i}$가 $\mathcal{F}_i$로 동형하게 사상됨을
증명해야 한다. $W \subset U_i$라 하자. 이 경우 $\mathcal{F}(W)$의
정의에 나오는 조건은
$s_j = \varphi_{ij}(s_i|_{W \cap U_j})$임을 뜻한다. 또한 붙이기 자료의
정의에 나오는 도표들의 가환성에 의해 {\it 임의의} 절단
$s \in \mathcal{F}_i(W)$에서 시작하여 $s_i = s$ 및
$s_j = \varphi_{ij}(s_i|_{W \cap U_j})$로 놓음으로써 양립하는 족을
얻을 수 있다.

\medskip\noindent
둘째 증명(개요). $\mathcal{B}$를 어떤 $i \in I$에 대해
$U \subset U_i$인 열린집합 $U \subset X$들의 집합이라 하자. 그러면
$\mathcal{B}$는 $X$의 위상의 기저이다. $U \in \mathcal{B}$에 대해
$U \subset U_i$인 $i \in I$를 택하고
$\mathcal{F}(U) = \mathcal{F}_i(U)$로 놓는다. 동형
$\varphi_{ij}$를 사용하면 이 처방이 “$i$의 선택과 무관함”을 알 수
있다. $\mathcal{F}_i$의 제한 사상들을 사용하면 $\mathcal{F}$가
$\mathcal{B}$ 위의 층임을 알 수 있다. 마지막으로 보조정리
\ref{sheaves-lemma-extend-off-basis}를 사용하여 $\mathcal{F}$를 $X$ 위의
유일한 층 $\mathcal{F}$로 확장한다.
\end{proof}

\begin{lemma}
\label{sheaves-lemma-glue-sheaves-structures}
$X$를 위상공간이라 하자. $X = \bigcup U_i$를 열린 덮개라 하자.
$(\mathcal{F}_i, \varphi_{ij})$를 아벨 군의 층, 각각 대수 구조의 층,
각각 $X$ 위의 어떤 환의 층 $\mathcal{O}$에 대한
$\mathcal{O}$-가군의 층의 붙이기 자료라 하자. 그러면 위의 보조정리
\ref{sheaves-lemma-glue-sheaves}의 증명에 나오는 구성은 각각 아벨 군의 층,
대수 구조의 층, $\mathcal{O}$-가군의 층을 준다.
\end{lemma}

\begin{proof}
이는 구성에서 열린집합 $W$ 위의 절단들의 집합
$\mathcal{F}(W)$가 사상들
$$
\xymatrix{
\prod_{i \in I} \mathcal{F}_i(W \cap U_i)
\ar@<1ex>[r]
\ar@<-1ex>[r]
&
\prod_{i, j\in I} \mathcal{F}_i(W \cap U_i \cap U_j)
}
$$
의 등화자로 주어지기 때문에 참이다. 상정한 각 경우에 이 등화자는
해당 범주의 대상이며, 그 바탕 집합은 인용한 보조정리에서 고려한
대상이다.
\end{proof}

\begin{lemma}
\label{sheaves-lemma-mapping-property-glue}
$X$를 위상공간이라 하자. $X = \bigcup_{i\in I} U_i$를 열린 덮개라
하자. 집합의 층 $\mathcal{F}$에 덮개 $X = \bigcup U_i$에 대한
다음 붙이기 자료의 족을 대응시키는 함자
$$
(\mathcal{F}|_{U_i},
(\mathcal{F}|_{U_i})|_{U_i \cap U_j}
\to
(\mathcal{F}|_{U_j})|_{U_i \cap U_j}
)
$$
는 $\Sh(X)$와 붙이기 자료의 범주 사이의 범주의 동치를 정의한다.
아벨 층, 각각 대수 구조의 층, 각각 $\mathcal{O}$-가군의 층에도
비슷한 명제가 성립한다.
\end{lemma}

\begin{proof}
함자는 보조정리 \ref{sheaves-lemma-glue-maps}에 의해 충만충실이고, 보조정리
\ref{sheaves-lemma-glue-sheaves}에 의해 (명시적으로 주어진 준역 함자를
통하여) 본질적으로 전사이다.
\end{proof}

\noindent
이 보조정리는 층 $\mathcal{F}$가 붙이기 자료
$(\mathcal{F}_i, \varphi_{ij})$로부터 구성되고 $\mathcal{G}$가 $X$
위의 층이면, 사상 $f : \mathcal{F} \to \mathcal{G}$를 준다는 것이
붙이기 사상 $\varphi_{ij}$와 양립하는 층의 사상들의 족
$$
f_i : \mathcal{F}_i \longrightarrow \mathcal{G}|_{U_i}
$$
를 주는 것과 같음을 뜻한다. 마찬가지로 층의 사상
$g : \mathcal{G} \to \mathcal{F}$를 준다는 것은 붙이기 사상
$\varphi_{ij}$와 양립하는 층의 사상들의 족
$$
g_i : \mathcal{G}|_{U_i} \longrightarrow \mathcal{F}_i
$$
를 주는 것과 같다.
